% =============================================================================
% ΚΕΦΑΛΑΙΟ 1: ΕΙΣΑΓΩΓΗ
% =============================================================================
% Αυτό το αρχείο περιέχει το πλήρες Κεφάλαιο 1 της διπλωματικής εργασίας
% "Μοντελοποίηση Αγορών Ηλεκτρικής Ενέργειας με Julia/JuMP"
% 
% Για ενσωμάτωση στο thesis.tex, χρησιμοποιήστε: % =============================================================================
% ΚΕΦΑΛΑΙΟ 1: ΕΙΣΑΓΩΓΗ
% =============================================================================
% Αυτό το αρχείο περιέχει το πλήρες Κεφάλαιο 1 της διπλωματικής εργασίας
% "Μοντελοποίηση Αγορών Ηλεκτρικής Ενέργειας με Julia/JuMP"
% 
% Για ενσωμάτωση στο thesis.tex, χρησιμοποιήστε: % =============================================================================
% ΚΕΦΑΛΑΙΟ 1: ΕΙΣΑΓΩΓΗ
% =============================================================================
% Αυτό το αρχείο περιέχει το πλήρες Κεφάλαιο 1 της διπλωματικής εργασίας
% "Μοντελοποίηση Αγορών Ηλεκτρικής Ενέργειας με Julia/JuMP"
% 
% Για ενσωμάτωση στο thesis.tex, χρησιμοποιήστε: % =============================================================================
% ΚΕΦΑΛΑΙΟ 1: ΕΙΣΑΓΩΓΗ
% =============================================================================
% Αυτό το αρχείο περιέχει το πλήρες Κεφάλαιο 1 της διπλωματικής εργασίας
% "Μοντελοποίηση Αγορών Ηλεκτρικής Ενέργειας με Julia/JuMP"
% 
% Για ενσωμάτωση στο thesis.tex, χρησιμοποιήστε: \input{chapter1_introduction}
% =============================================================================

\chapter{Εισαγωγή}
\label{ch:introduction}

\section{Κλιματική Αλλαγή και μέτρα για την αντιμετώπισή της}
\label{sec:climate-change}

Η Κλιματική Αλλαγή αναδεικνύεται ως η κορυφαία πρόκληση της εποχής μας, συνιστώντας μια υπαρξιακή απειλή για το φυσικό περιβάλλον, την συνοχή των ανθρώπινων κοινωνιών και την παγκόσμια οικονομική σταθερότητα. Η επιστημονική κοινότητα συμφωνεί ότι η ραγδαία αύξηση της μέσης θερμοκρασίας του πλανήτη και η εντατικοποίηση των ακραίων καιρικών φαινομένων οφείλονται, κατά κύριο λόγο, στις ανθρωπογενείς εκπομπές αερίων του θερμοκηπίου~\cite{lynas2021}.

Μέρος αυτών των εκπομπών προέρχεται από την χρήση ορυκτών καυσίμων όπως ο άνθρακας, το πετρέλαιο, ο λιγνίτης και το φυσικό αέριο, για την παραγωγή ηλεκτρικής ενέργειας. Η πρωτοκαθεδρία των ορυκτών καυσίμων παραμένει αδιαμφισβήτητη, αντιπροσωπεύοντας συνολικά σχεδόν το 60\% της παγκόσμιας παραγωγής ηλεκτρικής ενέργειας για το έτος 2024. Η εν λόγω αναλογία υπογραμμίζει την συνεχιζόμενη εξάρτηση των παγκόσμιων ενεργειακών συστημάτων από συμβατικές πηγές που επιτείνουν την κλιματική αλλαγή και τις επιπτώσεις αυτής. Αξίζει όμως να σημειωθεί ότι το ενεργειακό μίγμα βελτιώνεται σταδιακά, καθότι οι Ανανεώσιμες Πηγές Ενέργειας (ΑΠΕ) ανήλθαν συνολικά στο ένα τρίτο (33\%) της ηλεκτροπαραγωγής, όπου μαζί με την πυρηνική ενέργεια (9\%), άγγιξαν για πρώτη φορά το 40\% της παγκόσμιας ηλεκτροπαραγωγής ενέργειας~\cite{iea2024}.

\subsection{Παγκόσμιο επίπεδο}
\label{subsec:global-level}

Η Σύμβαση-Πλαίσιο των Ηνωμένων Εθνών για την Κλιματική Αλλαγή (UNFCCC), που συμφωνήθηκε το 1992 στη Διάσκεψη του Ρίο, αποτελεί τη βασική διεθνή συνθήκη για την αντιμετώπιση της κλιματικής αλλαγής. Στόχος της είναι η πρόληψη της επικίνδυνης ανθρωπογενούς παρέμβασης στο παγκόσμιο κλιματικό σύστημα. Από αυτήν, προέκυψαν συμφωνίες όπως το Πρωτόκολλο του Κιότο (1997), το οποίο εισήγαγε νομικά δεσμευτικούς στόχους μείωσης εκπομπών για τις ανεπτυγμένες χώρες και καθιέρωσε μηχανισμούς όπως η εμπορία δικαιωμάτων εκπομπών (Emissions Trading), ο μηχανισμός ``καθαρής'' ανάπτυξης (Clean Development Mechanism) και ο μηχανισμός κοινής εφαρμογής (Joint Implementation)~\cite{ypen-kyoto}. Το Πρωτόκολλο του Κιότο ήταν πρωτοπόρο για την εποχή, αλλά η αποτελεσματικότητά του περιορίστηκε από τη μη συμμετοχή ή αποχώρηση μεγάλων εκπομπέων όπως οι ΗΠΑ, η Ιαπωνία και ο Καναδάς~\cite{bassetti2022}.

Το 2015, η Συμφωνία του Παρισιού υιοθετήθηκε ως η νέα παγκόσμια συμφωνία για το κλίμα. Ο πρωταρχικός της στόχος είναι να συγκρατήσει ``την αύξηση της μέσης παγκόσμιας θερμοκρασίας πολύ κάτω από τους 2°C σε σχέση με τα προβιομηχανικά επίπεδα'' και να καταβάλει προσπάθειες ``να περιορίσει την αύξηση της θερμοκρασίας στον 1,5°C σε σχέση με τα προβιομηχανικά επίπεδα''~\cite{unfccc-paris}.

\subsection{Ευρωπαϊκό επίπεδο}
\label{subsec:european-level}

Πέραν των ανωτέρω συμφωνιών, η αντιμετώπιση της κλιματικής αλλαγής σε ευρωπαϊκό επίπεδο και πιο συγκεκριμένα όσον αφορά την Ευρωπαϊκή Ένωση, εστιάζει γύρω από τον στόχο της κλιματικής ουδετερότητας έως το 2050 και τις ενδιάμεσες δεσμεύσεις μείωσης των καθαρών εκπομπών τουλάχιστον 55\% έως το 2030 και τουλάχιστον 90\% έως το 2040, σε σχέση με τα επίπεδα του 1990. Η πρωτοβουλία-πλαίσιο που συγκεντρώνει αυτά τα πολιτικά μέτρα είναι το Ευρωπαϊκό Σχέδιο για το Περιβάλλον (European Green Deal), το οποίο προβλέπει αναθεώρηση της υπάρχουσας νομοθεσίας και εισαγωγή νέων πολιτικών σε τομείς όπως η ενεργειακή μετάβαση, η κυκλική οικονομία, η ανακαίνιση κτιρίων, η βιοποικιλότητα και η αγροτική πολιτική~\cite{ec-green-deal}.

Το πακέτο πολιτικών του Green Deal συνδέει ρητά την ενεργειακή αποδέσμευση από τα ορυκτά καύσιμα με μέτρα οικονομικής στήριξης, επενδύσεις σε τεχνολογικές καινοτομίες και κοινωνικά μέτρα για την αντιμετώπιση των επιπτώσεων στη συνοχή των περιοχών και των εργαζομένων. Η νομοθετική δέσμευση ενισχύεται από τον Ευρωπαϊκό Νόμο για το Κλίμα, ο οποίος καθιστά τον στόχο της ουδετερότητας νομικά δεσμευτικό και απαιτεί σαφείς ενδιάμεσους στόχους και μηχανισμούς παρακολούθησης για τα κράτη-μέλη~\cite{eu-climate-law}.

Ως αποτέλεσμα, η ΕΕ με το Green Deal συνδυάζει ρυθμιστικά εργαλεία (όπως αναθεωρήσεις αγοράς δικαιωμάτων εκπομπών, πρότυπα ενεργειακής απόδοσης και κανονισμούς για τη μεταφορά και τη βιομηχανία) με χρηματοδοτικά μέσα (ταμεία δίκαιης μετάβασης, προγράμματα επενδύσεων) και πολιτικές υποστήριξης των ΑΠΕ και της έρευνας. Η ευρωπαϊκή προσέγγιση προωθεί τόσο τη μείωση εκπομπών όσο και την ενίσχυση της κλιματικής ανθεκτικότητας των υποδομών και των κοινοτήτων, με έμφαση στη συνοχή των πολιτικών σε όλα τα επίπεδα διακυβέρνησης~\cite{sikora2020}.

\subsection{Εθνικό επίπεδο}
\label{subsec:national-level}

Σε εθνικό επίπεδο, η Ελλάδα, ως μέλος της ΕΕ, έχει υιοθετήσει ένα πλαίσιο πολιτικών και στρατηγικών που ευθυγραμμίζονται πλήρως με τους στόχους της Ευρωπαϊκής Ένωσης για την κλιματική ουδετερότητα έως το 2050 και τη μείωση των εκπομπών έως το 2030, ενσωματώνοντας μέτρα για την ενεργειακή μετάβαση, την απολιγνιτοποίηση, την επιτάχυνση των ΑΠΕ και την αναβάθμιση της ενεργειακής απόδοσης των κτιρίων.

Το παραπάνω πλαίσιο είναι το αναθεωρημένο Εθνικό Σχέδιο για την Ενέργεια και το Κλίμα (ΕΣΕΚ) το οποίο κυρώθηκε το 2024. Σε αυτό το έγγραφο αναγνωρίζονται οι επιπτώσεις που η κλιματική αλλαγή θα επιφέρει ειδικά στην Ελλάδα, όπως η ``μειωμένη διαθεσιμότητα υδάτων'' η οποία ``θέτει σε κίνδυνο τη λειτουργία θερμοηλεκτρικών μονάδων. Τα δίκτυα μεταφοράς και διανομής ηλεκτρικής ενέργειας, οι σταθμοί υψηλής τάσης, και λοιπά ενεργειακά δίκτυα και εγκαταστάσεις είναι τρωτά σε ακραία καιρικά φαινόμενα, ενώ οι παράκτιες ενεργειακές υποδομές, απειλούνται από την άνοδο της στάθμης της θάλασσας. Η αύξηση της μέσης θερμοκρασίας αναμένεται να μειώσει τις ενεργειακές ανάγκες για θέρμανση κατά τη χειμερινή περίοδο και να αυξήσει τις ανάγκες για ψύξη κατά τη θερινή περίοδο. Οι μεταβολές στη ζήτηση θα οδηγήσουν σε μεγαλύτερη διακύμανση φορτίων δοκιμάζοντας το σύστημα ηλεκτρικής ενέργειας. Επίσης, η κλιματική αλλαγή αναμένεται να επηρεάσει την απόδοση των συστημάτων ΑΠΕ''~\cite{esek2024}.

\subsection{Τεχνητή νοημοσύνη και ενεργειακή κατανάλωση}
\label{subsec:ai-energy}

Η ραγδαία επέκταση της Τεχνητής Νοημοσύνης (AI) και του υπολογιστικού νέφους (cloud computing) θέτει μία αυξανόμενη πρόκληση για την κλιματική αλλαγή, καθώς οδηγεί σε πρωτοφανή αύξηση της κατανάλωσης ενέργειας των κέντρων δεδομένων (data centers). Τα κέντρα δεδομένων, τα οποία ήδη καταναλώνουν περίπου το 1\% της παγκόσμιας ηλεκτρικής ενέργειας, αναμένεται να διπλασιάσουν αυτόν τον αριθμό έως το 2030, κυρίως λόγω των εξαιρετικά απαιτητικών εργασιών της τεχνητής νοημοσύνης, όπως η εκπαίδευση μεγάλων γλωσσικών μοντέλων (Large Language Models - LLMs). 

Η εκπαίδευση και η χρήση αυτών των μοντέλων απαιτεί την εντατική χρήση ενεργοβόρου υλικού όπως GPUs και TPUs, καθώς και μεγάλη ανάγκη για συστήματα ψύξης, τα οποία ευθύνονται για σχεδόν το 40\% της συνολικής ενεργειακής κατανάλωσης. Αυτό σημαίνει ότι η περιβαλλοντική επίπτωση αυξάνεται σημαντικά, εκτός εάν υπάρξει μια σημαντική στροφή σε ανανεώσιμες πηγές ενέργειας και προηγμένες τεχνικές ψύξης. Κατά συνέπεια, η έγκαιρη αντιμετώπιση της ενεργειακής ζήτησης της AI είναι κρίσιμη για τη διασφάλιση της βιωσιμότητας της ψηφιακής υποδομής και τον μετριασμό των αρνητικών επιπτώσεων στο περιβάλλον~\cite{chauhan2024}.


% =============================================================================
\section{Συστήματα Ηλεκτρικής Ενέργειας}
\label{sec:power-systems}

Τα συστήματα ηλεκτρικής ενέργειας αποτελούν το σύνολο των εγκαταστάσεων και μέσων για την παροχή ηλεκτρικής ενέργειας σε εξυπηρετούμενες περιοχές κατανάλωσης. Για την εύρυθμη λειτουργία ενός συστήματος ηλεκτρικής ενέργειας είναι απαραίτητη η παροχή ηλεκτρικής ενέργειας οπουδήποτε υπάρχει ζήτηση με το ελάχιστο δυνατό κόστος και τις ελάχιστες οικολογικές επιπτώσεις, εξασφαλίζοντας σταθερή συχνότητα, σταθερή τάση και υψηλή αξιοπιστία τροφοδότησης~\cite{vournas-kontaxis}.

Η τροφοδότηση των καταναλωτών με ηλεκτρική ενέργεια προϋποθέτει τρεις ξεχωριστές λειτουργίες του συστήματος ηλεκτρικής ενέργειας: την παραγωγή, τη μεταφορά, και τη διανομή. Η παραγόμενη ηλεκτρική ενέργεια πρέπει ιδανικά να απορροφάται από την κατανάλωση σε ίδιο χρόνο, καθώς η αποθήκευσή της παραμένει δύσκολη και τις περισσότερες φορές οικονομικά ασύμφορη, παρά την πρόοδο στην τεχνολογία αποθήκευσης ηλεκτρικής ενέργειας τα τελευταία χρόνια~\cite{hiesl2020}.

% TODO: Εισαγωγή σχήματος συστήματος ηλεκτρικής ενέργειας
\begin{figure}[htbp]
\centering
% \includegraphics[width=0.9\textwidth]{figures/power_system_overview.pdf}
\fbox{\parbox{0.8\textwidth}{\centering\vspace{2cm}[Σχήμα: Γενική δομή συστήματος ηλεκτρικής ενέργειας]\vspace{2cm}}}
\caption{Γενική δομή συστήματος ηλεκτρικής ενέργειας με τις τρεις βασικές λειτουργίες: παραγωγή, μεταφορά και διανομή.}
\label{fig:power-system-overview}
\end{figure}

\subsection{Παραγωγή Ηλεκτρικής Ενέργειας}
\label{subsec:generation}

Η παραγωγή ηλεκτρικής ενέργειας αφορά τη μετατροπή διαφόρων μορφών πρωτογενούς ενέργειας σε ηλεκτρική. Τα εργοστάσια παραγωγής διακρίνονται με βάση την τεχνολογία και την πηγή ενέργειας που χρησιμοποιούν.

Οι \textbf{θερμοηλεκτρικοί σταθμοί παραγωγής} αποτελούν τη συμβατική μορφή παραγωγής ηλεκτρικής ενέργειας. Σε αυτούς, η θερμική ενέργεια που προέρχεται από την καύση ορυκτών καυσίμων (λιγνίτης, φυσικό αέριο, πετρέλαιο) ή από πυρηνικές αντιδράσεις μετατρέπεται σε μηχανική μέσω ατμοστροβίλων ή αεριοστροβίλων, και στη συνέχεια σε ηλεκτρική μέσω γεννητριών. Οι μονάδες αυτές χαρακτηρίζονται από υψηλή διαθεσιμότητα και ελεγχόμενη παραγωγή, αλλά συνεπάγονται σημαντικές εκπομπές αερίων του θερμοκηπίου.

Οι \textbf{υδροηλεκτρικοί σταθμοί} αξιοποιούν τη δυναμική ενέργεια του νερού για την κίνηση υδροστροβίλων. Διακρίνονται σε συμβατικούς (με φράγμα και ταμιευτήρα) και αντλησιοταμιευτικούς, οι οποίοι λειτουργούν και ως συστήματα αποθήκευσης ενέργειας.

Οι \textbf{Ανανεώσιμες Πηγές Ενέργειας (ΑΠΕ)} περιλαμβάνουν τεχνολογίες όπως τα αιολικά πάρκα, τα φωτοβολταϊκά συστήματα, τη γεωθερμία και τη βιομάζα. Οι ΑΠΕ χαρακτηρίζονται από μηδενικές ή ελάχιστες εκπομπές κατά τη λειτουργία τους, αλλά η παραγωγή τους εξαρτάται από τις καιρικές συνθήκες, καθιστώντας τη μεταβλητή και εν μέρει απρόβλεπτη.

\subsection{Μεταφορά Ηλεκτρικής Ενέργειας}
\label{subsec:transmission}

Η μεταφορά ηλεκτρικής ενέργειας αφορά τη μεταφορά μεγάλων ποσοτήτων ηλεκτρικής ενέργειας από τους σταθμούς παραγωγής προς τα κέντρα κατανάλωσης μέσω του δικτύου υψηλής τάσης. Η υψηλή τάση (150kV, 400kV) χρησιμοποιείται για τη μείωση των απωλειών μεταφοράς, καθώς για δεδομένη ισχύ, η αύξηση της τάσης οδηγεί σε μείωση του ρεύματος και συνεπώς των απωλειών λόγω θέρμανσης των αγωγών.

Αρμόδιος για την αξιόπιστη μεταφορά ηλεκτρικής ενέργειας στην ελληνική επικράτεια είναι ο Ανεξάρτητος Διαχειριστής Μεταφοράς Ηλεκτρικής Ενέργειας (ΑΔΜΗΕ). Ο ΑΔΜΗΕ διαχειρίζεται το Διασυνδεδεμένο Σύστημα Μεταφοράς, το οποίο περιλαμβάνει βασικά στοιχεία όπως εναέριες γραμμές μεταφοράς 400kV, 150kV και 66kV, υπόγειες και υποβρύχιες καλωδιακές γραμμές 150kV και 400kV (AC και DC), υποσταθμούς 150/20kV, καθώς και Κέντρα Υπερυψηλής Τάσης (ΚΥΤ) 400/150kV.

Το Ελληνικό Σύστημα Μεταφοράς λειτουργεί παράλληλα με το διασυνδεδεμένο Ευρωπαϊκό Σύστημα, υπό τον ευρύτερο συντονισμό του Ευρωπαϊκού Δικτύου Διαχειριστών Συστημάτων Μεταφοράς Ηλεκτρικής Ενέργειας (European Network of Transmission System Operators for Electricity - ENTSO-E). Η σύνδεση αυτή πραγματοποιείται μέσω διασυνδετικών γραμμών μεταφοράς 400kV, που ενώνουν το ελληνικό σύστημα με εκείνα της Αλβανίας, της Βουλγαρίας, της Βόρειας Μακεδονίας και της Τουρκίας. Επιπλέον, υπάρχει ασύγχρονη διασύνδεση με την Ιταλία μέσω υποβρυχίου συνδέσμου συνεχούς ρεύματος τάσης 400kV~\cite{admie-system}.

\subsection{Διανομή Ηλεκτρικής Ενέργειας}
\label{subsec:distribution}

Από τους υποσταθμούς υποβιβασμού τάσης ξεκινά η διανομή ηλεκτρικής ενέργειας, το σύνολο των διαδικασιών που απαιτούνται για τον αξιόπιστο διαμοιρασμό της στους μετρητές των καταναλωτών, διατηρώντας την τάση στους ζυγούς και τις ροές στις γραμμές του δικτύου εντός επιτρεπτών ορίων.

Αρμόδιος για τη διαχείριση του δικτύου διανομής ηλεκτρικής ενέργειας στην ελληνική επικράτεια είναι ο Διαχειριστής Ελληνικού Δικτύου Διανομής Ηλεκτρικής Ενέργειας (ΔΕΔΔΗΕ). Ο ΔΕΔΔΗΕ είναι επίσης υπεύθυνος για τη διαχείριση των αγορών των Μη Διασυνδεδεμένων Νησιών (ΜΔΝ)~\cite{nomos4001}.

Εντός του δικτύου διανομής ηλεκτρικής ενέργειας έχει ξεκινήσει και η ανάπτυξη των τεχνολογιών έξυπνων δικτύων (smart grids) με χρήση έξυπνων μετρητών (smart meters), συστημάτων εποπτικού ελέγχου (SCADA) και ολοκληρωμένων πληροφοριακών συστημάτων διανομής (DMS), που λειτουργούν καταλυτικά στην εκμετάλλευση των οικιακών εγκαταστάσεων ΑΠΕ, συστημάτων αποθήκευσης καθώς και τις τεχνολογίες Vehicle to Grid (V2G)~\cite{fang2012}.

Μέσω της αμφίδρομης ροής πληροφοριών και ενέργειας, τα έξυπνα δίκτυα επιτρέπουν στους διαχειριστές δικτύου να παρακολουθούν σε πραγματικό χρόνο την κατανάλωση, να εντοπίζουν βλάβες άμεσα και να διαχειρίζονται καλύτερα την κατανομή φορτίου. Επιπλέον, παρέχουν δυνατότητες συμμετοχής των τελικών καταναλωτών στην αγορά μέσω demand response, όπου οι χρήστες προσαρμόζουν τη ζήτησή τους ανάλογα με τα σήματα της αγοράς ή τις τιμές ενέργειας. Η μετάβαση προς τα έξυπνα δίκτυα είναι αναγκαία για την υλοποίηση των στόχων της ενεργειακής μετάβασης, την ενίσχυση της ενεργειακής απόδοσης και την επίτευξη ενός βιώσιμου, ψηφιοποιημένου και αποκεντρωμένου ενεργειακού συστήματος.


% =============================================================================
\section{Αγορές Ηλεκτρικής Ενέργειας}
\label{sec:energy-markets}

Για την εξασφάλιση της ισορροπίας προσφοράς-ζήτησης και της απρόσκοπτης παροχής ενέργειας στους καταναλωτές στο ελάχιστο δυνατό κόστος, τα συστήματα ηλεκτρικής ενέργειας ενθυλακώνονται σε αγορές ηλεκτρικής ενέργειας. Οι αγορές ηλεκτρικής ενέργειας χαρακτηρίζονται από ιδιαιτερότητες σε σχέση με άλλες αγορές εμπορευμάτων ή καταναλωτικών αγαθών, λόγω του υψηλού κόστους αποθήκευσης της ηλεκτρικής ενέργειας~\cite{mo2025}.

% TODO: Εισαγωγή σχήματος δομής αγοράς
\begin{figure}[htbp]
\centering
% \includegraphics[width=0.9\textwidth]{figures/market_structure.pdf}
\fbox{\parbox{0.8\textwidth}{\centering\vspace{2cm}[Σχήμα: Δομή αγορών ηλεκτρικής ενέργειας]\vspace{2cm}}}
\caption{Χρονική αλληλουχία των αγορών ηλεκτρικής ενέργειας, από την προθεσμιακή αγορά έως την αγορά εξισορρόπησης.}
\label{fig:market-structure}
\end{figure}

\subsection{Μονοπωλιακή Αγορά}
\label{subsec:monopoly-market}

Ιστορικά, η ανάπτυξη του τομέα της ηλεκτρικής ενέργειας απαιτούσε τεράστια κεφαλαιουχική δαπάνη, μακροχρόνιο σχεδιασμό και υψηλό επίπεδο τεχνικής εξειδίκευσης, στοιχεία που καθιστούσαν πρακτικά αδύνατη τη συμμετοχή ιδιωτών. Για τον λόγο αυτό, το κράτος ή κρατικά ελεγχόμενες επιχειρήσεις αναλάμβαναν την ευθύνη για την κατασκευή, λειτουργία και διαχείριση του συστήματος παραγωγής, μεταφοράς και διανομής ηλεκτρικής ενέργειας. Το μοντέλο που διαμορφώθηκε ήταν μονοπωλιακό, με χαρακτηριστικό γνώρισμα την κάθετη ολοκλήρωση, δηλαδή τον έλεγχο όλων των σταδίων της αλυσίδας αξίας --- από την παραγωγή έως την προμήθεια στον τελικό καταναλωτή --- από έναν και μόνο φορέα. Στην ελληνική επικράτεια το ρόλο αυτό είχε αναλάβει η Δημόσια Επιχείρηση Ηλεκτρισμού (ΔΕΗ).

Σε αυτό το πλαίσιο απουσίας ανταγωνισμού οι καταναλωτές δεν είχαν τη δυνατότητα επιλογής παρόχου. Οι τιμές της ηλεκτρικής ενέργειας καθορίζονταν διοικητικά από τις αρμόδιες ρυθμιστικές αρχές, με βάση το κόστος και ένα σταθερό περιθώριο κέρδους. Επιπλέον, οι επενδύσεις και η επέκταση του δικτύου σχεδιάζονταν κεντρικά, με στόχο την κάλυψη της ζήτησης και την ασφάλεια εφοδιασμού. Παρά τα πλεονεκτήματα σταθερότητας και ελέγχου, το μονοπωλιακό αυτό μοντέλο συχνά οδηγούσε σε αναποτελεσματικότητα, καθυστερήσεις στην υιοθέτηση καινοτομιών και υψηλότερο κόστος για τον τελικό καταναλωτή~\cite{vlados2021}.

Συγκεκριμένα, η ρύθμιση ενός μονοπωλιακού φορέα, όπως μια κρατική ή ρυθμιζόμενη επιχείρηση ηλεκτρικής ενέργειας, αποτελεί χαρακτηριστικό πρόβλημα εντολέα-εντολοδόχου (principal-agent). Ο ρυθμιστής (εντολέας) επιδιώκει τη μεγιστοποίηση του κοινωνικού πλεονάσματος, ενώ η επιχείρηση (εντολοδόχος) στοχεύει στη μεγιστοποίηση του κέρδους. Ωστόσο, ο ρυθμιστής δεν διαθέτει πλήρη πληροφόρηση για το κόστος και τα πραγματικά δεδομένα λειτουργίας της επιχείρησης, η οποία έχει κίνητρο να τα υπερεκτιμά, προβάλλοντας αυξημένες ανάγκες για ασφάλεια και επενδύσεις. Η ασυμμετρία πληροφόρησης αυτή δημιουργεί προβλήματα όπως ηθικός κίνδυνος και αρνητική επιλογή, οδηγώντας σε αποκλίσεις από την αποδοτικότητα του τέλειου ανταγωνισμού~\cite{grossman-hart1992}.

Για την αντιμετώπιση αυτών των στρεβλώσεων, ο ρυθμιστής μπορεί είτε να μειώσει την ασυμμετρία πληροφόρησης, μέσω ελέγχων ή συγκριτικής αξιολόγησης (benchmarking), είτε να εφαρμόσει στρατηγικές κινήτρων, ώστε να ενισχύσει την αποδοτικότητα της επιχείρησης. Παρά τις προσπάθειες αυτές, η ισορροπία που επιτυγχάνεται μεταξύ ρυθμιστή και επιχείρησης δεν φτάνει ποτέ το επίπεδο αποδοτικότητας που προκύπτει σε καθεστώς τέλειου ανταγωνισμού, με αποτέλεσμα την ύπαρξη κόστους πρακτόρευσης (agency cost) για την κοινωνία~\cite{panda-leepsa2017}.

\subsection{Άνοιγμα της αγοράς}
\label{subsec:market-liberalization}

Για την αντιμετώπιση των προβλημάτων αποδοτικότητας που προκύπτουν λόγω του κόστους πρακτόρευσης στο μονοπωλιακό μοντέλο αγοράς, από τη δεκαετία του 1990 η Ευρωπαϊκή Ένωση ξεκίνησε μια διαδικασία απελευθέρωσης των αγορών ηλεκτρικής ενέργειας, με στόχο τη δημιουργία μιας ενιαίας, ανταγωνιστικής αγοράς. Τα κύρια στάδια αυτής της διαδικασίας περιλαμβάνουν το Πρώτο Ενεργειακό Πακέτο (1996), όπου η Οδηγία 96/92/ΕΚ έθεσε τα θεμέλια για το άνοιγμα των αγορών, εισάγοντας την έννοια του διαχωρισμού των δραστηριοτήτων και την πρόσβαση τρίτων στο δίκτυο. Το Δεύτερο Ενεργειακό Πακέτο (2003) με τις Οδηγίες 2003/54/ΕΚ και 2003/55/ΕΚ επιτάχυνε το άνοιγμα της αγοράς, καθιστώντας υποχρεωτικό το νομικό διαχωρισμό των δραστηριοτήτων μεταφοράς και διανομής από την παραγωγή και προμήθεια. Το Τρίτο Ενεργειακό Πακέτο (2009) εισήγαγε αυστηρότερους κανόνες για το διαχωρισμό της ιδιοκτησίας, ενίσχυσε τις αρμοδιότητες των εθνικών ρυθμιστικών αρχών και ίδρυσε τον Οργανισμό Συνεργασίας των Ρυθμιστικών Αρχών Ενέργειας (ACER). Τέλος, το Τέταρτο Ενεργειακό Πακέτο (2019), γνωστό και ως ``Καθαρή Ενέργεια για όλους τους Ευρωπαίους'', εστιάζει στην ενσωμάτωση των ανανεώσιμων πηγών ενέργειας, την ενεργειακή αποδοτικότητα και νέους κανόνες για την αγορά ηλεκτρικής ενέργειας~\cite{eikeland-duffield2011, eu-energy-packages}.

Με βάση τα παραπάνω, το άνοιγμα της αγοράς οδήγησε στον ιδιοκτησιακό διαχωρισμό των δραστηριοτήτων, τη δημιουργία ανταγωνιστικών αγορών χονδρικής και λιανικής, την ελευθερία επιλογής προμηθευτή για τους καταναλωτές και την ανάπτυξη νέων επιχειρηματικών μοντέλων και υπηρεσιών.

\subsection{Φορείς της αγοράς}
\label{subsec:market-participants}

Στο σύγχρονο μοντέλο των απελευθερωμένων αγορών ηλεκτρικής ενέργειας, διάφοροι φορείς διαδραματίζουν καθοριστικό ρόλο. Οι \textbf{Παραγωγοί} είναι εταιρείες που διαθέτουν και λειτουργούν μονάδες παραγωγής ηλεκτρικής ενέργειας, περιλαμβάνοντας τόσο συμβατικούς παραγωγούς (θερμικές μονάδες) όσο και παραγωγούς ΑΠΕ. Ο \textbf{Διαχειριστής Συστήματος Μεταφοράς} (TSO - Transmission System Operator) είναι υπεύθυνος για τη λειτουργία, συντήρηση και ανάπτυξη του συστήματος μεταφοράς υψηλής τάσης, καθώς και για την εξισορρόπηση του συστήματος σε πραγματικό χρόνο. Για την Ελλάδα, αυτός ο φορέας είναι ο ΑΔΜΗΕ. Ο \textbf{Διαχειριστής Δικτύου Διανομής} (DSO - Distribution System Operator) είναι υπεύθυνος για τη λειτουργία, συντήρηση και ανάπτυξη του δικτύου διανομής μέσης και χαμηλής τάσης. Για την Ελλάδα, αυτός ο φορέας είναι ο ΔΕΔΔΗΕ.

Οι \textbf{Προμηθευτές} αγοράζουν ηλεκτρική ενέργεια από την αγορά χονδρικής και την πωλούν στους τελικούς καταναλωτές. Οι \textbf{Έμποροι} (Traders) αγοράζουν και πωλούν ηλεκτρική ενέργεια στις χονδρικές αγορές με σκοπό το κέρδος, χωρίς να είναι απαραίτητα παραγωγοί ή προμηθευτές. Τα \textbf{Χρηματιστήρια Ενέργειας} (Power Exchanges) είναι οργανωμένες αγορές όπου διαπραγματεύονται προϊόντα ηλεκτρικής ενέργειας. Παραδείγματα περιλαμβάνουν το EPEX SPOT, το Nord Pool, το GME και για την Ελλάδα το HENEX (EnExGroup). Ο \textbf{Διαχειριστής Αγοράς} είναι υπεύθυνος για την οργάνωση και λειτουργία των διαφόρων τμημάτων της αγοράς ηλεκτρικής ενέργειας. Οι \textbf{Ρυθμιστικές Αρχές} εποπτεύουν τη λειτουργία της αγοράς, θεσπίζουν κανόνες και διασφαλίζουν τον υγιή ανταγωνισμό. Σε ευρωπαϊκό επίπεδο, ο ACER συντονίζει τις εθνικές ρυθμιστικές αρχές. Οι \textbf{Καταναλωτές} περιλαμβάνουν βιομηχανικούς, εμπορικούς και οικιακούς χρήστες. Στις σύγχρονες αγορές, οι καταναλωτές μπορούν επίσης να συμμετέχουν ενεργά στην αγορά μέσω της απόκρισης ζήτησης και της αυτοπαραγωγής.

\subsection{Προθεσμιακή Αγορά - OTC}
\label{subsec:forward-market}

Η Προθεσμιακή Αγορά αποτελεί το πρώτο στάδιο της χονδρικής αγοράς ηλεκτρικής ενέργειας, όπου οι συμμετέχοντες μπορούν να διαπραγματευτούν και να συνάψουν συμβόλαια για μελλοντική παράδοση ηλεκτρικής ενέργειας. Διακρίνεται σε \textbf{Οργανωμένη Προθεσμιακή Αγορά}, η οποία λειτουργεί μέσω χρηματιστηρίων ενέργειας (π.χ. EEX - European Energy Exchange) όπου διαπραγματεύονται τυποποιημένα προθεσμιακά συμβόλαια (futures, forwards, options), και σε \textbf{Εξωχρηματιστηριακή Αγορά} (OTC - Over The Counter), όπου οι συναλλαγές πραγματοποιούνται απευθείας μεταξύ των αντισυμβαλλομένων, με συμβόλαια προσαρμοσμένα στις ανάγκες τους.

Τα βασικά χαρακτηριστικά της Προθεσμιακής Αγοράς περιλαμβάνουν τον χρονικό ορίζοντα (τα συμβόλαια μπορεί να αφορούν παράδοση ηλεκτρικής ενέργειας από μερικές ημέρες έως αρκετά χρόνια στο μέλλον), τη διαχείριση κινδύνου (επιτρέπει στους συμμετέχοντες να αντισταθμίσουν τον κίνδυνο μεταβολής των τιμών), την ευελιξία (ιδιαίτερα στην αγορά OTC, τα συμβόλαια μπορούν να προσαρμοστούν ως προς τον όγκο, τη διάρκεια και τους όρους παράδοσης), τη ρευστότητα (στις οργανωμένες αγορές, τα τυποποιημένα προϊόντα έχουν συνήθως υψηλότερη ρευστότητα) και τον πιστωτικό κίνδυνο (στις συναλλαγές OTC υπάρχει κίνδυνος αθέτησης από τον αντισυμβαλλόμενο).

Η Προθεσμιακή Αγορά διαδραματίζει σημαντικό ρόλο στη διαμόρφωση μακροπρόθεσμων τιμολογιακών σημάτων, συμβάλλει στην οικονομική σταθερότητα των συμμετεχόντων και παρέχει ενδείξεις για μελλοντικές επενδύσεις στον τομέα της ηλεκτρικής ενέργειας.

\subsection{Προημερήσια Αγορά Ηλεκτρικής Ενέργειας}
\label{subsec:day-ahead-market}

Η Προημερήσια Αγορά (Day-Ahead Market - DAM) αποτελεί κεντρικό πυλώνα της χονδρικής αγοράς ηλεκτρικής ενέργειας στην Ευρώπη. Σε αυτή την αγορά, οι συμμετέχοντες υποβάλλουν προσφορές για αγορά και πώληση ηλεκτρικής ενέργειας για κάθε ώρα της επόμενης ημέρας.

Η αγορά λειτουργεί συνήθως με \textbf{μηχανισμό δημοπρασίας μοναδικής τιμής} (uniform pricing auction), όπου όλες οι αποδεκτές προσφορές εκκαθαρίζονται στην οριακή τιμή συστήματος (System Marginal Price - SMP). Όσον αφορά το χρονοδιάγραμμα, οι προσφορές υποβάλλονται μέχρι μια συγκεκριμένη ώρα (π.χ. 12:00 CET) και τα αποτελέσματα ανακοινώνονται λίγο αργότερα, καθορίζοντας το πρόγραμμα παραγωγής και κατανάλωσης για την επόμενη ημέρα. Τα προϊόντα συνήθως αφορούν ωριαίες ή ημίωρες ποσότητες ηλεκτρικής ενέργειας, ενώ σε ορισμένες αγορές υπάρχουν και πιο σύνθετα προϊόντα όπως μπλοκ προσφορές και συνδεδεμένες προσφορές.

Στην Ευρώπη, οι περισσότερες εθνικές προημερήσιες αγορές είναι συζευγμένες μέσω του μηχανισμού \textbf{Price Coupling of Regions} (PCR), που βελτιστοποιεί τη χρήση των διασυνδέσεων και προωθεί τη σύγκλιση των τιμών. Για την επίλυση της συζευγμένης προημερήσιας αγοράς χρησιμοποιείται ο \textbf{αλγόριθμος EUPHEMIA} (EU Pan-European Hybrid Electricity Market Integration Algorithm), ο οποίος λαμβάνει υπόψη τις προσφορές αγοράς και πώλησης, τους περιορισμούς μεταφοράς και τα διαθέσιμα περιθώρια των διασυνδέσεων~\cite{euphemia-public-description}.

Η Προημερήσια Αγορά παρέχει ένα αξιόπιστο σημείο αναφοράς για τις τιμές ηλεκτρικής ενέργειας και μία βάση για τη διαμόρφωση του προγράμματος λειτουργίας του συστήματος. Επίσης αποτελεί έναν μηχανισμό για την αποτελεσματική κατανομή των πόρων και των διασυνοριακών ικανοτήτων μεταφοράς, καθώς και ένα σημαντικό εργαλείο για τη βελτιστοποίηση του κόστους παραγωγής σε πανευρωπαϊκό επίπεδο.

Οι τιμές που προκύπτουν από την Προημερήσια Αγορά αντανακλούν το οριακό κόστος παραγωγής ηλεκτρικής ενέργειας και επηρεάζονται από παράγοντες όπως η ζήτηση, η διαθεσιμότητα των μονάδων παραγωγής, η παραγωγή από ΑΠΕ και οι περιορισμοί του δικτύου.

\subsection{Ενδοημερήσια Αγορά Ηλεκτρικής Ενέργειας}
\label{subsec:intraday-market}

Η Ενδοημερήσια Αγορά (Intraday Market - IDM) επιτρέπει στους συμμετέχοντες να προσαρμόσουν τις θέσεις τους μετά το κλείσιμο της Προημερήσιας Αγοράς και πριν την πραγματική ώρα παράδοσης της ηλεκτρικής ενέργειας. Παρέχει τη δυνατότητα διαπραγμάτευσης ηλεκτρικής ενέργειας σε χρονικό ορίζοντα ``σχεδόν πραγματικού χρόνου''.

Σε αντίθεση με την Προημερήσια Αγορά που λειτουργεί με δημοπρασία, η Ενδοημερήσια Αγορά συνήθως λειτουργεί με μηχανισμό \textbf{συνεχούς διαπραγμάτευσης} (continuous trading), όπου οι συναλλαγές πραγματοποιούνται συνεχώς μέχρι το κλείσιμο της πύλης (gate closure) που είναι συνήθως 5-60 λεπτά πριν την παράδοση. Η αγορά επιτρέπει τη διαχείριση αβεβαιότητας και την αντιμετώπιση απρόβλεπτων γεγονότων όπως βλάβες μονάδων παραγωγής, μεταβολές στην πρόβλεψη παραγωγής ΑΠΕ ή απροσδόκητες διακυμάνσεις στη ζήτηση. Το Ευρωπαϊκό μοντέλο \textbf{XBID} (Cross-Border Intraday) παρέχει μια ενιαία πλατφόρμα για διασυνοριακές ενδοημερήσιες συναλλαγές σε όλη την Ευρώπη, ενισχύοντας τη ρευστότητα και την αποτελεσματικότητα της αγοράς.

Η Ενδοημερήσια Αγορά προσφέρει βελτιωμένη αποτελεσματικότητα του συστήματος, καθώς οι συμμετέχοντες μπορούν να προσαρμόσουν τις θέσεις τους με βάση πιο πρόσφατες πληροφορίες. Αποτελεί ένα σημαντικό εργαλείο για την ενσωμάτωση μεταβλητών ΑΠΕ, επιτρέποντας την αντιμετώπιση σφαλμάτων πρόβλεψης και βοηθάει στη μείωση της ανάγκης για ενεργοποίηση πόρων εξισορρόπησης, καθώς αποκλίσεις μπορούν να διορθωθούν μέσω της αγοράς αυτής. Η σημασία της Ενδοημερήσιας Αγοράς αυξάνεται συνεχώς, ιδιαίτερα με την αυξανόμενη διείσδυση των ΑΠΕ και τις προκλήσεις που αυτή δημιουργεί στην πρόβλεψη της παραγωγής.

\subsection{Αγορά Εξισορρόπησης}
\label{subsec:balancing-market}

Η Αγορά Εξισορρόπησης (Balancing Market) αποτελεί το τελευταίο χρονικά στάδιο των αγορών ηλεκτρικής ενέργειας και λειτουργεί σε πραγματικό ή κοντά σε πραγματικό χρόνο. Ο βασικός της στόχος είναι να διασφαλίσει την ισορροπία μεταξύ παραγωγής και κατανάλωσης ηλεκτρικής ενέργειας ανά πάσα στιγμή, διατηρώντας τη συχνότητα του συστήματος στα επιτρεπτά όρια.

Η αγορά διαχειρίζεται από τον Διαχειριστή του Συστήματος Μεταφοράς (TSO), ο οποίος είναι υπεύθυνος για την εξισορρόπηση του συστήματος σε πραγματικό χρόνο. Τα προϊόντα εξισορρόπησης περιλαμβάνουν τις Εφεδρείες Διατήρησης Συχνότητας (FCR - Frequency Containment Reserves) με αυτόματη απόκριση εντός δευτερολέπτων, τις Εφεδρείες Αποκατάστασης Συχνότητας (FRR - Frequency Restoration Reserves) που διακρίνονται σε αυτόματες (aFRR) και χειροκίνητες (mFRR), καθώς και τις Εφεδρείες Αντικατάστασης (RR - Replacement Reserves) για την αποκατάσταση των προηγούμενων εφεδρειών~\cite{entsoe-balancing}.

Για την αποτελεσματικότερη λειτουργία της αγοράς, έχουν αναπτυχθεί πανευρωπαϊκές πλατφόρμες εξισορρόπησης όπως η πλατφόρμα PICASSO για την ανταλλαγή aFRR, η πλατφόρμα MARI για την ανταλλαγή mFRR και η TERRE για την ανταλλαγή RR. Η τιμολόγηση ενέργειας εξισορρόπησης μπορεί να βασίζεται είτε στην τιμή προσφοράς (pay-as-bid) είτε στην οριακή τιμή (pay-as-cleared), με την τελευταία να προτιμάται στις περισσότερες ευρωπαϊκές αγορές. Οι συμμετέχοντες που αποκλίνουν από το πρόγραμμά τους χρεώνονται ή πιστώνονται με την τιμή απόκλισης (imbalance price).

Η Αγορά Εξισορρόπησης είναι καθοριστικής σημασίας για τη διασφάλιση της αξιοπιστίας και ασφάλειας του συστήματος ηλεκτρικής ενέργειας, την αποτελεσματική διαχείριση των διακυμάνσεων της παραγωγής από ΑΠΕ και την παροχή οικονομικών κινήτρων για επενδύσεις σε ευέλικτους πόρους. Με την αυξανόμενη διείσδυση των ΑΠΕ, η σημασία και η πολυπλοκότητα της Αγοράς Εξισορρόπησης αυξάνεται, απαιτώντας συνεχή εξέλιξη των μηχανισμών της και στενότερη διασυνοριακή συνεργασία μεταξύ των Διαχειριστών Συστημάτων Μεταφοράς.


% =============================================================================
\section{Σκοπός και Στόχοι της Εργασίας}
\label{sec:objectives}

Η παρούσα διπλωματική εργασία έχει ως κεντρικό στόχο την ανάπτυξη ενός ολοκληρωμένου, ανοιχτού κώδικα εργαλείου για την προσομοίωση της Προημερήσιας Αγοράς ηλεκτρικής ενέργειας, με έμφαση στον αλγόριθμο EUPHEMIA που χρησιμοποιείται για την εκκαθάριση της συζευγμένης ευρωπαϊκής αγοράς.

\subsection{Κίνητρο και Αναγκαιότητα}
\label{subsec:motivation}

Η κατανόηση των μηχανισμών λειτουργίας των αγορών ηλεκτρικής ενέργειας αποτελεί κρίσιμο παράγοντα για όλους τους εμπλεκόμενους φορείς, από τους παραγωγούς και τους προμηθευτές έως τους ρυθμιστές και τους ερευνητές. Ωστόσο, η πρόσβαση σε εργαλεία προσομοίωσης που αναπαράγουν με ακρίβεια τη λειτουργία της αγοράς είναι περιορισμένη, καθώς τα υπάρχοντα συστήματα είναι συνήθως κλειστά και ιδιόκτητα.

Η ενεργειακή μετάβαση και η αυξανόμενη διείσδυση των ΑΠΕ δημιουργούν νέες προκλήσεις για τη λειτουργία των αγορών ηλεκτρικής ενέργειας. Η μεταβλητότητα της παραγωγής από ανανεώσιμες πηγές, οι περιορισμοί του δικτύου μεταφοράς και η ανάγκη για διασυνοριακή συνεργασία καθιστούν τη μοντελοποίηση της αγοράς ένα σύνθετο πρόβλημα βελτιστοποίησης. Η δυνατότητα προσομοίωσης διαφορετικών σεναρίων και η κατανόηση των παραγόντων που επηρεάζουν τις τιμές της ηλεκτρικής ενέργειας είναι απαραίτητη για τη λήψη αποφάσεων τόσο σε επιχειρησιακό όσο και σε πολιτικό επίπεδο.

\subsection{Επιστημονική Συνεισφορά}
\label{subsec:contribution}

Η συνεισφορά της παρούσας εργασίας εντοπίζεται σε τρεις κύριους άξονες.

Πρώτον, στην \textbf{ανάπτυξη ολοκληρωμένης υποδομής δεδομένων}. Σχεδιάστηκε και υλοποιήθηκε ένα αυτοματοποιημένο σύστημα συλλογής δεδομένων από την πλατφόρμα διαφάνειας του ENTSO-E, το οποίο αποθηκεύει και διαχειρίζεται δεδομένα για μονάδες παραγωγής, φορτία, προβλέψεις ΑΠΕ, διασυνοριακές μεταφορές και μη διαθεσιμότητες μονάδων. Η υποδομή αυτή επιτρέπει την εκτέλεση προσομοιώσεων με πραγματικά δεδομένα για οποιαδήποτε ζώνη προσφοράς (bidding zone) της Ευρώπης.

Δεύτερον, στην \textbf{υλοποίηση αλγορίθμων βελτιστοποίησης αγοράς}. Αναπτύχθηκε μια σειρά από αλγορίθμους που καλύπτουν το πλήρες φάσμα της λειτουργίας της Προημερήσιας Αγοράς, συμπεριλαμβανομένου του προβλήματος Unit Commitment για τον προσδιορισμό του βέλτιστου προγράμματος λειτουργίας των μονάδων παραγωγής, της μηχανής στρατηγικών υποβολής προσφορών, του βιβλίου εντολών οριακού κόστους (merit-order) που κατασκευάζει ανταγωνιστικό αντιπαράδειγμα της αγοράς από τα θεμελιώδη μεγέθη της (τιμές καυσίμων και CO$_2$, αξία νερού, καθαρές εισαγωγές), του αλγορίθμου MPCC (Mathematical Programming with Complementarity Constraints) για την εκκαθάριση της αγοράς και της μοντελοποίησης δικτύου με περιορισμούς ATC (Available Transfer Capacity) για την πολυζωνική εκκαθάριση της αγοράς.

Τρίτον, στη \textbf{δημιουργία ανοιχτού κώδικα εργαλείου με ποσοτικά επικυρωμένη ακρίβεια}. Ολόκληρο το σύστημα αναπτύχθηκε σε γλώσσα Julia με χρήση του πλαισίου JuMP για τη μοντελοποίηση προβλημάτων βελτιστοποίησης, και ο κώδικας διατίθεται ελεύθερα με άδεια ανοιχτού κώδικα. Το σύστημα αναπαράγει τις τιμές της ελληνικής ζώνης με συντελεστή συσχέτισης 0{,}84 σε ανάλυση δεκαπενταλέπτου επί συνεχούς τετραμήνου (Κεφάλαιο~\ref{ch:results})· εξ όσων είμαστε σε θέση να γνωρίζουμε, πρόκειται για το πρώτο ανοιχτό εργαλείο που μοντελοποιεί την αγορά αυτή με τεκμηριωμένη ακρίβεια αυτού του επιπέδου.

\subsection{Ερευνητικά Ερωτήματα}
\label{subsec:research-questions}

Η εργασία επιχειρεί να απαντήσει στα ακόλουθα ερευνητικά ερωτήματα. Πρώτον, είναι δυνατή η αναπαραγωγή των τιμών της Προημερήσιας Αγοράς με ικανοποιητική ακρίβεια χρησιμοποιώντας δημοσίως διαθέσιμα δεδομένα και αλγορίθμους ανοιχτού κώδικα; Δεύτερον, ποια είναι η επίδραση των διαφορετικών στρατηγικών υποβολής προσφορών στις τελικές τιμές της αγοράς; Τρίτον, πώς επηρεάζουν οι περιορισμοί διασυνοριακής μεταφοράς τη διαμόρφωση των τιμών σε γειτονικές ζώνες προσφοράς;


% =============================================================================
\section{Δομή της Εργασίας}
\label{sec:thesis-structure}

Η παρούσα διπλωματική εργασία οργανώνεται σε έξι κεφάλαια, τα οποία καλύπτουν το θεωρητικό υπόβαθρο, την αρχιτεκτονική του συστήματος, την υλοποίηση, τα αποτελέσματα και τα συμπεράσματα.

Στο \textbf{Κεφάλαιο~\ref{ch:introduction}} (Εισαγωγή), το παρόν κεφάλαιο, παρουσιάζεται το πλαίσιο της εργασίας με αναφορά στην κλιματική αλλαγή, την ενεργειακή μετάβαση και τη δομή των ευρωπαϊκών αγορών ηλεκτρικής ενέργειας. Επίσης, καθορίζονται οι στόχοι και η συνεισφορά της εργασίας.

Στο \textbf{Κεφάλαιο~\ref{ch:theoretical-background}} (Θεωρητικό Υπόβαθρο) αναλύονται οι θεωρητικές βάσεις της εργασίας, συμπεριλαμβανομένων των μεθόδων μαθηματικού προγραμματισμού, του προβλήματος Unit Commitment, του αλγορίθμου EUPHEMIA και των στρατηγικών υποβολής προσφορών.

Στο \textbf{Κεφάλαιο~\ref{ch:system-architecture}} (Αρχιτεκτονική Συστήματος και Δεδομένα) περιγράφεται η αρχιτεκτονική του συστήματος που αναπτύχθηκε, με έμφαση στο pipeline συλλογής δεδομένων από το ENTSO-E, το σχήμα της βάσης δεδομένων και το οικοσύστημα Julia/JuMP.

Στο \textbf{Κεφάλαιο~\ref{ch:implementation}} (Υλοποίηση) παρουσιάζεται αναλυτικά η υλοποίηση των αλγορίθμων, συμπεριλαμβανομένου του solver για το πρόβλημα Unit Commitment, της μηχανής στρατηγικών υποβολής προσφορών, του βιβλίου εντολών οριακού κόστους, του αλγορίθμου MPCC για την εκκαθάριση της αγοράς και της μοντελοποίησης των περιορισμών δικτύου.

Στο \textbf{Κεφάλαιο~\ref{ch:results}} (Αποτελέσματα και Αξιολόγηση) παρουσιάζονται τα αποτελέσματα της εφαρμογής του συστήματος σε πραγματικά δεδομένα, η σύγκριση με τις πραγματικές τιμές της αγοράς και η ανάλυση της απόδοσης του συστήματος.

Στο \textbf{Κεφάλαιο~\ref{ch:conclusions}} (Συμπεράσματα) συνοψίζονται τα κύρια ευρήματα της εργασίας, αναφέρονται οι περιορισμοί και προτείνονται κατευθύνσεις για μελλοντική έρευνα.

Τέλος, στο \textbf{Παράρτημα} παρατίθενται συμπληρωματικά στοιχεία και σύνδεσμοι προς τον κώδικα της εφαρμογής.

% =============================================================================

\chapter{Εισαγωγή}
\label{ch:introduction}

\section{Κλιματική Αλλαγή και μέτρα για την αντιμετώπισή της}
\label{sec:climate-change}

Η Κλιματική Αλλαγή αναδεικνύεται ως η κορυφαία πρόκληση της εποχής μας, συνιστώντας μια υπαρξιακή απειλή για το φυσικό περιβάλλον, την συνοχή των ανθρώπινων κοινωνιών και την παγκόσμια οικονομική σταθερότητα. Η επιστημονική κοινότητα συμφωνεί ότι η ραγδαία αύξηση της μέσης θερμοκρασίας του πλανήτη και η εντατικοποίηση των ακραίων καιρικών φαινομένων οφείλονται, κατά κύριο λόγο, στις ανθρωπογενείς εκπομπές αερίων του θερμοκηπίου~\cite{lynas2021}.

Μέρος αυτών των εκπομπών προέρχεται από την χρήση ορυκτών καυσίμων όπως ο άνθρακας, το πετρέλαιο, ο λιγνίτης και το φυσικό αέριο, για την παραγωγή ηλεκτρικής ενέργειας. Η πρωτοκαθεδρία των ορυκτών καυσίμων παραμένει αδιαμφισβήτητη, αντιπροσωπεύοντας συνολικά σχεδόν το 60\% της παγκόσμιας παραγωγής ηλεκτρικής ενέργειας για το έτος 2024. Η εν λόγω αναλογία υπογραμμίζει την συνεχιζόμενη εξάρτηση των παγκόσμιων ενεργειακών συστημάτων από συμβατικές πηγές που επιτείνουν την κλιματική αλλαγή και τις επιπτώσεις αυτής. Αξίζει όμως να σημειωθεί ότι το ενεργειακό μίγμα βελτιώνεται σταδιακά, καθότι οι Ανανεώσιμες Πηγές Ενέργειας (ΑΠΕ) ανήλθαν συνολικά στο ένα τρίτο (33\%) της ηλεκτροπαραγωγής, όπου μαζί με την πυρηνική ενέργεια (9\%), άγγιξαν για πρώτη φορά το 40\% της παγκόσμιας ηλεκτροπαραγωγής ενέργειας~\cite{iea2024}.

\subsection{Παγκόσμιο επίπεδο}
\label{subsec:global-level}

Η Σύμβαση-Πλαίσιο των Ηνωμένων Εθνών για την Κλιματική Αλλαγή (UNFCCC), που συμφωνήθηκε το 1992 στη Διάσκεψη του Ρίο, αποτελεί τη βασική διεθνή συνθήκη για την αντιμετώπιση της κλιματικής αλλαγής. Στόχος της είναι η πρόληψη της επικίνδυνης ανθρωπογενούς παρέμβασης στο παγκόσμιο κλιματικό σύστημα. Από αυτήν, προέκυψαν συμφωνίες όπως το Πρωτόκολλο του Κιότο (1997), το οποίο εισήγαγε νομικά δεσμευτικούς στόχους μείωσης εκπομπών για τις ανεπτυγμένες χώρες και καθιέρωσε μηχανισμούς όπως η εμπορία δικαιωμάτων εκπομπών (Emissions Trading), ο μηχανισμός ``καθαρής'' ανάπτυξης (Clean Development Mechanism) και ο μηχανισμός κοινής εφαρμογής (Joint Implementation)~\cite{ypen-kyoto}. Το Πρωτόκολλο του Κιότο ήταν πρωτοπόρο για την εποχή, αλλά η αποτελεσματικότητά του περιορίστηκε από τη μη συμμετοχή ή αποχώρηση μεγάλων εκπομπέων όπως οι ΗΠΑ, η Ιαπωνία και ο Καναδάς~\cite{bassetti2022}.

Το 2015, η Συμφωνία του Παρισιού υιοθετήθηκε ως η νέα παγκόσμια συμφωνία για το κλίμα. Ο πρωταρχικός της στόχος είναι να συγκρατήσει ``την αύξηση της μέσης παγκόσμιας θερμοκρασίας πολύ κάτω από τους 2°C σε σχέση με τα προβιομηχανικά επίπεδα'' και να καταβάλει προσπάθειες ``να περιορίσει την αύξηση της θερμοκρασίας στον 1,5°C σε σχέση με τα προβιομηχανικά επίπεδα''~\cite{unfccc-paris}.

\subsection{Ευρωπαϊκό επίπεδο}
\label{subsec:european-level}

Πέραν των ανωτέρω συμφωνιών, η αντιμετώπιση της κλιματικής αλλαγής σε ευρωπαϊκό επίπεδο και πιο συγκεκριμένα όσον αφορά την Ευρωπαϊκή Ένωση, εστιάζει γύρω από τον στόχο της κλιματικής ουδετερότητας έως το 2050 και τις ενδιάμεσες δεσμεύσεις μείωσης των καθαρών εκπομπών τουλάχιστον 55\% έως το 2030 και τουλάχιστον 90\% έως το 2040, σε σχέση με τα επίπεδα του 1990. Η πρωτοβουλία-πλαίσιο που συγκεντρώνει αυτά τα πολιτικά μέτρα είναι το Ευρωπαϊκό Σχέδιο για το Περιβάλλον (European Green Deal), το οποίο προβλέπει αναθεώρηση της υπάρχουσας νομοθεσίας και εισαγωγή νέων πολιτικών σε τομείς όπως η ενεργειακή μετάβαση, η κυκλική οικονομία, η ανακαίνιση κτιρίων, η βιοποικιλότητα και η αγροτική πολιτική~\cite{ec-green-deal}.

Το πακέτο πολιτικών του Green Deal συνδέει ρητά την ενεργειακή αποδέσμευση από τα ορυκτά καύσιμα με μέτρα οικονομικής στήριξης, επενδύσεις σε τεχνολογικές καινοτομίες και κοινωνικά μέτρα για την αντιμετώπιση των επιπτώσεων στη συνοχή των περιοχών και των εργαζομένων. Η νομοθετική δέσμευση ενισχύεται από τον Ευρωπαϊκό Νόμο για το Κλίμα, ο οποίος καθιστά τον στόχο της ουδετερότητας νομικά δεσμευτικό και απαιτεί σαφείς ενδιάμεσους στόχους και μηχανισμούς παρακολούθησης για τα κράτη-μέλη~\cite{eu-climate-law}.

Ως αποτέλεσμα, η ΕΕ με το Green Deal συνδυάζει ρυθμιστικά εργαλεία (όπως αναθεωρήσεις αγοράς δικαιωμάτων εκπομπών, πρότυπα ενεργειακής απόδοσης και κανονισμούς για τη μεταφορά και τη βιομηχανία) με χρηματοδοτικά μέσα (ταμεία δίκαιης μετάβασης, προγράμματα επενδύσεων) και πολιτικές υποστήριξης των ΑΠΕ και της έρευνας. Η ευρωπαϊκή προσέγγιση προωθεί τόσο τη μείωση εκπομπών όσο και την ενίσχυση της κλιματικής ανθεκτικότητας των υποδομών και των κοινοτήτων, με έμφαση στη συνοχή των πολιτικών σε όλα τα επίπεδα διακυβέρνησης~\cite{sikora2020}.

\subsection{Εθνικό επίπεδο}
\label{subsec:national-level}

Σε εθνικό επίπεδο, η Ελλάδα, ως μέλος της ΕΕ, έχει υιοθετήσει ένα πλαίσιο πολιτικών και στρατηγικών που ευθυγραμμίζονται πλήρως με τους στόχους της Ευρωπαϊκής Ένωσης για την κλιματική ουδετερότητα έως το 2050 και τη μείωση των εκπομπών έως το 2030, ενσωματώνοντας μέτρα για την ενεργειακή μετάβαση, την απολιγνιτοποίηση, την επιτάχυνση των ΑΠΕ και την αναβάθμιση της ενεργειακής απόδοσης των κτιρίων.

Το παραπάνω πλαίσιο είναι το αναθεωρημένο Εθνικό Σχέδιο για την Ενέργεια και το Κλίμα (ΕΣΕΚ) το οποίο κυρώθηκε το 2024. Σε αυτό το έγγραφο αναγνωρίζονται οι επιπτώσεις που η κλιματική αλλαγή θα επιφέρει ειδικά στην Ελλάδα, όπως η ``μειωμένη διαθεσιμότητα υδάτων'' η οποία ``θέτει σε κίνδυνο τη λειτουργία θερμοηλεκτρικών μονάδων. Τα δίκτυα μεταφοράς και διανομής ηλεκτρικής ενέργειας, οι σταθμοί υψηλής τάσης, και λοιπά ενεργειακά δίκτυα και εγκαταστάσεις είναι τρωτά σε ακραία καιρικά φαινόμενα, ενώ οι παράκτιες ενεργειακές υποδομές, απειλούνται από την άνοδο της στάθμης της θάλασσας. Η αύξηση της μέσης θερμοκρασίας αναμένεται να μειώσει τις ενεργειακές ανάγκες για θέρμανση κατά τη χειμερινή περίοδο και να αυξήσει τις ανάγκες για ψύξη κατά τη θερινή περίοδο. Οι μεταβολές στη ζήτηση θα οδηγήσουν σε μεγαλύτερη διακύμανση φορτίων δοκιμάζοντας το σύστημα ηλεκτρικής ενέργειας. Επίσης, η κλιματική αλλαγή αναμένεται να επηρεάσει την απόδοση των συστημάτων ΑΠΕ''~\cite{esek2024}.

\subsection{Τεχνητή νοημοσύνη και ενεργειακή κατανάλωση}
\label{subsec:ai-energy}

Η ραγδαία επέκταση της Τεχνητής Νοημοσύνης (AI) και του υπολογιστικού νέφους (cloud computing) θέτει μία αυξανόμενη πρόκληση για την κλιματική αλλαγή, καθώς οδηγεί σε πρωτοφανή αύξηση της κατανάλωσης ενέργειας των κέντρων δεδομένων (data centers). Τα κέντρα δεδομένων, τα οποία ήδη καταναλώνουν περίπου το 1\% της παγκόσμιας ηλεκτρικής ενέργειας, αναμένεται να διπλασιάσουν αυτόν τον αριθμό έως το 2030, κυρίως λόγω των εξαιρετικά απαιτητικών εργασιών της τεχνητής νοημοσύνης, όπως η εκπαίδευση μεγάλων γλωσσικών μοντέλων (Large Language Models - LLMs). 

Η εκπαίδευση και η χρήση αυτών των μοντέλων απαιτεί την εντατική χρήση ενεργοβόρου υλικού όπως GPUs και TPUs, καθώς και μεγάλη ανάγκη για συστήματα ψύξης, τα οποία ευθύνονται για σχεδόν το 40\% της συνολικής ενεργειακής κατανάλωσης. Αυτό σημαίνει ότι η περιβαλλοντική επίπτωση αυξάνεται σημαντικά, εκτός εάν υπάρξει μια σημαντική στροφή σε ανανεώσιμες πηγές ενέργειας και προηγμένες τεχνικές ψύξης. Κατά συνέπεια, η έγκαιρη αντιμετώπιση της ενεργειακής ζήτησης της AI είναι κρίσιμη για τη διασφάλιση της βιωσιμότητας της ψηφιακής υποδομής και τον μετριασμό των αρνητικών επιπτώσεων στο περιβάλλον~\cite{chauhan2024}.


% =============================================================================
\section{Συστήματα Ηλεκτρικής Ενέργειας}
\label{sec:power-systems}

Τα συστήματα ηλεκτρικής ενέργειας αποτελούν το σύνολο των εγκαταστάσεων και μέσων για την παροχή ηλεκτρικής ενέργειας σε εξυπηρετούμενες περιοχές κατανάλωσης. Για την εύρυθμη λειτουργία ενός συστήματος ηλεκτρικής ενέργειας είναι απαραίτητη η παροχή ηλεκτρικής ενέργειας οπουδήποτε υπάρχει ζήτηση με το ελάχιστο δυνατό κόστος και τις ελάχιστες οικολογικές επιπτώσεις, εξασφαλίζοντας σταθερή συχνότητα, σταθερή τάση και υψηλή αξιοπιστία τροφοδότησης~\cite{vournas-kontaxis}.

Η τροφοδότηση των καταναλωτών με ηλεκτρική ενέργεια προϋποθέτει τρεις ξεχωριστές λειτουργίες του συστήματος ηλεκτρικής ενέργειας: την παραγωγή, τη μεταφορά, και τη διανομή. Η παραγόμενη ηλεκτρική ενέργεια πρέπει ιδανικά να απορροφάται από την κατανάλωση σε ίδιο χρόνο, καθώς η αποθήκευσή της παραμένει δύσκολη και τις περισσότερες φορές οικονομικά ασύμφορη, παρά την πρόοδο στην τεχνολογία αποθήκευσης ηλεκτρικής ενέργειας τα τελευταία χρόνια~\cite{hiesl2020}.

% TODO: Εισαγωγή σχήματος συστήματος ηλεκτρικής ενέργειας
\begin{figure}[htbp]
\centering
% \includegraphics[width=0.9\textwidth]{figures/power_system_overview.pdf}
\fbox{\parbox{0.8\textwidth}{\centering\vspace{2cm}[Σχήμα: Γενική δομή συστήματος ηλεκτρικής ενέργειας]\vspace{2cm}}}
\caption{Γενική δομή συστήματος ηλεκτρικής ενέργειας με τις τρεις βασικές λειτουργίες: παραγωγή, μεταφορά και διανομή.}
\label{fig:power-system-overview}
\end{figure}

\subsection{Παραγωγή Ηλεκτρικής Ενέργειας}
\label{subsec:generation}

Η παραγωγή ηλεκτρικής ενέργειας αφορά τη μετατροπή διαφόρων μορφών πρωτογενούς ενέργειας σε ηλεκτρική. Τα εργοστάσια παραγωγής διακρίνονται με βάση την τεχνολογία και την πηγή ενέργειας που χρησιμοποιούν.

Οι \textbf{θερμοηλεκτρικοί σταθμοί παραγωγής} αποτελούν τη συμβατική μορφή παραγωγής ηλεκτρικής ενέργειας. Σε αυτούς, η θερμική ενέργεια που προέρχεται από την καύση ορυκτών καυσίμων (λιγνίτης, φυσικό αέριο, πετρέλαιο) ή από πυρηνικές αντιδράσεις μετατρέπεται σε μηχανική μέσω ατμοστροβίλων ή αεριοστροβίλων, και στη συνέχεια σε ηλεκτρική μέσω γεννητριών. Οι μονάδες αυτές χαρακτηρίζονται από υψηλή διαθεσιμότητα και ελεγχόμενη παραγωγή, αλλά συνεπάγονται σημαντικές εκπομπές αερίων του θερμοκηπίου.

Οι \textbf{υδροηλεκτρικοί σταθμοί} αξιοποιούν τη δυναμική ενέργεια του νερού για την κίνηση υδροστροβίλων. Διακρίνονται σε συμβατικούς (με φράγμα και ταμιευτήρα) και αντλησιοταμιευτικούς, οι οποίοι λειτουργούν και ως συστήματα αποθήκευσης ενέργειας.

Οι \textbf{Ανανεώσιμες Πηγές Ενέργειας (ΑΠΕ)} περιλαμβάνουν τεχνολογίες όπως τα αιολικά πάρκα, τα φωτοβολταϊκά συστήματα, τη γεωθερμία και τη βιομάζα. Οι ΑΠΕ χαρακτηρίζονται από μηδενικές ή ελάχιστες εκπομπές κατά τη λειτουργία τους, αλλά η παραγωγή τους εξαρτάται από τις καιρικές συνθήκες, καθιστώντας τη μεταβλητή και εν μέρει απρόβλεπτη.

\subsection{Μεταφορά Ηλεκτρικής Ενέργειας}
\label{subsec:transmission}

Η μεταφορά ηλεκτρικής ενέργειας αφορά τη μεταφορά μεγάλων ποσοτήτων ηλεκτρικής ενέργειας από τους σταθμούς παραγωγής προς τα κέντρα κατανάλωσης μέσω του δικτύου υψηλής τάσης. Η υψηλή τάση (150kV, 400kV) χρησιμοποιείται για τη μείωση των απωλειών μεταφοράς, καθώς για δεδομένη ισχύ, η αύξηση της τάσης οδηγεί σε μείωση του ρεύματος και συνεπώς των απωλειών λόγω θέρμανσης των αγωγών.

Αρμόδιος για την αξιόπιστη μεταφορά ηλεκτρικής ενέργειας στην ελληνική επικράτεια είναι ο Ανεξάρτητος Διαχειριστής Μεταφοράς Ηλεκτρικής Ενέργειας (ΑΔΜΗΕ). Ο ΑΔΜΗΕ διαχειρίζεται το Διασυνδεδεμένο Σύστημα Μεταφοράς, το οποίο περιλαμβάνει βασικά στοιχεία όπως εναέριες γραμμές μεταφοράς 400kV, 150kV και 66kV, υπόγειες και υποβρύχιες καλωδιακές γραμμές 150kV και 400kV (AC και DC), υποσταθμούς 150/20kV, καθώς και Κέντρα Υπερυψηλής Τάσης (ΚΥΤ) 400/150kV.

Το Ελληνικό Σύστημα Μεταφοράς λειτουργεί παράλληλα με το διασυνδεδεμένο Ευρωπαϊκό Σύστημα, υπό τον ευρύτερο συντονισμό του Ευρωπαϊκού Δικτύου Διαχειριστών Συστημάτων Μεταφοράς Ηλεκτρικής Ενέργειας (European Network of Transmission System Operators for Electricity - ENTSO-E). Η σύνδεση αυτή πραγματοποιείται μέσω διασυνδετικών γραμμών μεταφοράς 400kV, που ενώνουν το ελληνικό σύστημα με εκείνα της Αλβανίας, της Βουλγαρίας, της Βόρειας Μακεδονίας και της Τουρκίας. Επιπλέον, υπάρχει ασύγχρονη διασύνδεση με την Ιταλία μέσω υποβρυχίου συνδέσμου συνεχούς ρεύματος τάσης 400kV~\cite{admie-system}.

\subsection{Διανομή Ηλεκτρικής Ενέργειας}
\label{subsec:distribution}

Από τους υποσταθμούς υποβιβασμού τάσης ξεκινά η διανομή ηλεκτρικής ενέργειας, το σύνολο των διαδικασιών που απαιτούνται για τον αξιόπιστο διαμοιρασμό της στους μετρητές των καταναλωτών, διατηρώντας την τάση στους ζυγούς και τις ροές στις γραμμές του δικτύου εντός επιτρεπτών ορίων.

Αρμόδιος για τη διαχείριση του δικτύου διανομής ηλεκτρικής ενέργειας στην ελληνική επικράτεια είναι ο Διαχειριστής Ελληνικού Δικτύου Διανομής Ηλεκτρικής Ενέργειας (ΔΕΔΔΗΕ). Ο ΔΕΔΔΗΕ είναι επίσης υπεύθυνος για τη διαχείριση των αγορών των Μη Διασυνδεδεμένων Νησιών (ΜΔΝ)~\cite{nomos4001}.

Εντός του δικτύου διανομής ηλεκτρικής ενέργειας έχει ξεκινήσει και η ανάπτυξη των τεχνολογιών έξυπνων δικτύων (smart grids) με χρήση έξυπνων μετρητών (smart meters), συστημάτων εποπτικού ελέγχου (SCADA) και ολοκληρωμένων πληροφοριακών συστημάτων διανομής (DMS), που λειτουργούν καταλυτικά στην εκμετάλλευση των οικιακών εγκαταστάσεων ΑΠΕ, συστημάτων αποθήκευσης καθώς και τις τεχνολογίες Vehicle to Grid (V2G)~\cite{fang2012}.

Μέσω της αμφίδρομης ροής πληροφοριών και ενέργειας, τα έξυπνα δίκτυα επιτρέπουν στους διαχειριστές δικτύου να παρακολουθούν σε πραγματικό χρόνο την κατανάλωση, να εντοπίζουν βλάβες άμεσα και να διαχειρίζονται καλύτερα την κατανομή φορτίου. Επιπλέον, παρέχουν δυνατότητες συμμετοχής των τελικών καταναλωτών στην αγορά μέσω demand response, όπου οι χρήστες προσαρμόζουν τη ζήτησή τους ανάλογα με τα σήματα της αγοράς ή τις τιμές ενέργειας. Η μετάβαση προς τα έξυπνα δίκτυα είναι αναγκαία για την υλοποίηση των στόχων της ενεργειακής μετάβασης, την ενίσχυση της ενεργειακής απόδοσης και την επίτευξη ενός βιώσιμου, ψηφιοποιημένου και αποκεντρωμένου ενεργειακού συστήματος.


% =============================================================================
\section{Αγορές Ηλεκτρικής Ενέργειας}
\label{sec:energy-markets}

Για την εξασφάλιση της ισορροπίας προσφοράς-ζήτησης και της απρόσκοπτης παροχής ενέργειας στους καταναλωτές στο ελάχιστο δυνατό κόστος, τα συστήματα ηλεκτρικής ενέργειας ενθυλακώνονται σε αγορές ηλεκτρικής ενέργειας. Οι αγορές ηλεκτρικής ενέργειας χαρακτηρίζονται από ιδιαιτερότητες σε σχέση με άλλες αγορές εμπορευμάτων ή καταναλωτικών αγαθών, λόγω του υψηλού κόστους αποθήκευσης της ηλεκτρικής ενέργειας~\cite{mo2025}.

% TODO: Εισαγωγή σχήματος δομής αγοράς
\begin{figure}[htbp]
\centering
% \includegraphics[width=0.9\textwidth]{figures/market_structure.pdf}
\fbox{\parbox{0.8\textwidth}{\centering\vspace{2cm}[Σχήμα: Δομή αγορών ηλεκτρικής ενέργειας]\vspace{2cm}}}
\caption{Χρονική αλληλουχία των αγορών ηλεκτρικής ενέργειας, από την προθεσμιακή αγορά έως την αγορά εξισορρόπησης.}
\label{fig:market-structure}
\end{figure}

\subsection{Μονοπωλιακή Αγορά}
\label{subsec:monopoly-market}

Ιστορικά, η ανάπτυξη του τομέα της ηλεκτρικής ενέργειας απαιτούσε τεράστια κεφαλαιουχική δαπάνη, μακροχρόνιο σχεδιασμό και υψηλό επίπεδο τεχνικής εξειδίκευσης, στοιχεία που καθιστούσαν πρακτικά αδύνατη τη συμμετοχή ιδιωτών. Για τον λόγο αυτό, το κράτος ή κρατικά ελεγχόμενες επιχειρήσεις αναλάμβαναν την ευθύνη για την κατασκευή, λειτουργία και διαχείριση του συστήματος παραγωγής, μεταφοράς και διανομής ηλεκτρικής ενέργειας. Το μοντέλο που διαμορφώθηκε ήταν μονοπωλιακό, με χαρακτηριστικό γνώρισμα την κάθετη ολοκλήρωση, δηλαδή τον έλεγχο όλων των σταδίων της αλυσίδας αξίας --- από την παραγωγή έως την προμήθεια στον τελικό καταναλωτή --- από έναν και μόνο φορέα. Στην ελληνική επικράτεια το ρόλο αυτό είχε αναλάβει η Δημόσια Επιχείρηση Ηλεκτρισμού (ΔΕΗ).

Σε αυτό το πλαίσιο απουσίας ανταγωνισμού οι καταναλωτές δεν είχαν τη δυνατότητα επιλογής παρόχου. Οι τιμές της ηλεκτρικής ενέργειας καθορίζονταν διοικητικά από τις αρμόδιες ρυθμιστικές αρχές, με βάση το κόστος και ένα σταθερό περιθώριο κέρδους. Επιπλέον, οι επενδύσεις και η επέκταση του δικτύου σχεδιάζονταν κεντρικά, με στόχο την κάλυψη της ζήτησης και την ασφάλεια εφοδιασμού. Παρά τα πλεονεκτήματα σταθερότητας και ελέγχου, το μονοπωλιακό αυτό μοντέλο συχνά οδηγούσε σε αναποτελεσματικότητα, καθυστερήσεις στην υιοθέτηση καινοτομιών και υψηλότερο κόστος για τον τελικό καταναλωτή~\cite{vlados2021}.

Συγκεκριμένα, η ρύθμιση ενός μονοπωλιακού φορέα, όπως μια κρατική ή ρυθμιζόμενη επιχείρηση ηλεκτρικής ενέργειας, αποτελεί χαρακτηριστικό πρόβλημα εντολέα-εντολοδόχου (principal-agent). Ο ρυθμιστής (εντολέας) επιδιώκει τη μεγιστοποίηση του κοινωνικού πλεονάσματος, ενώ η επιχείρηση (εντολοδόχος) στοχεύει στη μεγιστοποίηση του κέρδους. Ωστόσο, ο ρυθμιστής δεν διαθέτει πλήρη πληροφόρηση για το κόστος και τα πραγματικά δεδομένα λειτουργίας της επιχείρησης, η οποία έχει κίνητρο να τα υπερεκτιμά, προβάλλοντας αυξημένες ανάγκες για ασφάλεια και επενδύσεις. Η ασυμμετρία πληροφόρησης αυτή δημιουργεί προβλήματα όπως ηθικός κίνδυνος και αρνητική επιλογή, οδηγώντας σε αποκλίσεις από την αποδοτικότητα του τέλειου ανταγωνισμού~\cite{grossman-hart1992}.

Για την αντιμετώπιση αυτών των στρεβλώσεων, ο ρυθμιστής μπορεί είτε να μειώσει την ασυμμετρία πληροφόρησης, μέσω ελέγχων ή συγκριτικής αξιολόγησης (benchmarking), είτε να εφαρμόσει στρατηγικές κινήτρων, ώστε να ενισχύσει την αποδοτικότητα της επιχείρησης. Παρά τις προσπάθειες αυτές, η ισορροπία που επιτυγχάνεται μεταξύ ρυθμιστή και επιχείρησης δεν φτάνει ποτέ το επίπεδο αποδοτικότητας που προκύπτει σε καθεστώς τέλειου ανταγωνισμού, με αποτέλεσμα την ύπαρξη κόστους πρακτόρευσης (agency cost) για την κοινωνία~\cite{panda-leepsa2017}.

\subsection{Άνοιγμα της αγοράς}
\label{subsec:market-liberalization}

Για την αντιμετώπιση των προβλημάτων αποδοτικότητας που προκύπτουν λόγω του κόστους πρακτόρευσης στο μονοπωλιακό μοντέλο αγοράς, από τη δεκαετία του 1990 η Ευρωπαϊκή Ένωση ξεκίνησε μια διαδικασία απελευθέρωσης των αγορών ηλεκτρικής ενέργειας, με στόχο τη δημιουργία μιας ενιαίας, ανταγωνιστικής αγοράς. Τα κύρια στάδια αυτής της διαδικασίας περιλαμβάνουν το Πρώτο Ενεργειακό Πακέτο (1996), όπου η Οδηγία 96/92/ΕΚ έθεσε τα θεμέλια για το άνοιγμα των αγορών, εισάγοντας την έννοια του διαχωρισμού των δραστηριοτήτων και την πρόσβαση τρίτων στο δίκτυο. Το Δεύτερο Ενεργειακό Πακέτο (2003) με τις Οδηγίες 2003/54/ΕΚ και 2003/55/ΕΚ επιτάχυνε το άνοιγμα της αγοράς, καθιστώντας υποχρεωτικό το νομικό διαχωρισμό των δραστηριοτήτων μεταφοράς και διανομής από την παραγωγή και προμήθεια. Το Τρίτο Ενεργειακό Πακέτο (2009) εισήγαγε αυστηρότερους κανόνες για το διαχωρισμό της ιδιοκτησίας, ενίσχυσε τις αρμοδιότητες των εθνικών ρυθμιστικών αρχών και ίδρυσε τον Οργανισμό Συνεργασίας των Ρυθμιστικών Αρχών Ενέργειας (ACER). Τέλος, το Τέταρτο Ενεργειακό Πακέτο (2019), γνωστό και ως ``Καθαρή Ενέργεια για όλους τους Ευρωπαίους'', εστιάζει στην ενσωμάτωση των ανανεώσιμων πηγών ενέργειας, την ενεργειακή αποδοτικότητα και νέους κανόνες για την αγορά ηλεκτρικής ενέργειας~\cite{eikeland-duffield2011, eu-energy-packages}.

Με βάση τα παραπάνω, το άνοιγμα της αγοράς οδήγησε στον ιδιοκτησιακό διαχωρισμό των δραστηριοτήτων, τη δημιουργία ανταγωνιστικών αγορών χονδρικής και λιανικής, την ελευθερία επιλογής προμηθευτή για τους καταναλωτές και την ανάπτυξη νέων επιχειρηματικών μοντέλων και υπηρεσιών.

\subsection{Φορείς της αγοράς}
\label{subsec:market-participants}

Στο σύγχρονο μοντέλο των απελευθερωμένων αγορών ηλεκτρικής ενέργειας, διάφοροι φορείς διαδραματίζουν καθοριστικό ρόλο. Οι \textbf{Παραγωγοί} είναι εταιρείες που διαθέτουν και λειτουργούν μονάδες παραγωγής ηλεκτρικής ενέργειας, περιλαμβάνοντας τόσο συμβατικούς παραγωγούς (θερμικές μονάδες) όσο και παραγωγούς ΑΠΕ. Ο \textbf{Διαχειριστής Συστήματος Μεταφοράς} (TSO - Transmission System Operator) είναι υπεύθυνος για τη λειτουργία, συντήρηση και ανάπτυξη του συστήματος μεταφοράς υψηλής τάσης, καθώς και για την εξισορρόπηση του συστήματος σε πραγματικό χρόνο. Για την Ελλάδα, αυτός ο φορέας είναι ο ΑΔΜΗΕ. Ο \textbf{Διαχειριστής Δικτύου Διανομής} (DSO - Distribution System Operator) είναι υπεύθυνος για τη λειτουργία, συντήρηση και ανάπτυξη του δικτύου διανομής μέσης και χαμηλής τάσης. Για την Ελλάδα, αυτός ο φορέας είναι ο ΔΕΔΔΗΕ.

Οι \textbf{Προμηθευτές} αγοράζουν ηλεκτρική ενέργεια από την αγορά χονδρικής και την πωλούν στους τελικούς καταναλωτές. Οι \textbf{Έμποροι} (Traders) αγοράζουν και πωλούν ηλεκτρική ενέργεια στις χονδρικές αγορές με σκοπό το κέρδος, χωρίς να είναι απαραίτητα παραγωγοί ή προμηθευτές. Τα \textbf{Χρηματιστήρια Ενέργειας} (Power Exchanges) είναι οργανωμένες αγορές όπου διαπραγματεύονται προϊόντα ηλεκτρικής ενέργειας. Παραδείγματα περιλαμβάνουν το EPEX SPOT, το Nord Pool, το GME και για την Ελλάδα το HENEX (EnExGroup). Ο \textbf{Διαχειριστής Αγοράς} είναι υπεύθυνος για την οργάνωση και λειτουργία των διαφόρων τμημάτων της αγοράς ηλεκτρικής ενέργειας. Οι \textbf{Ρυθμιστικές Αρχές} εποπτεύουν τη λειτουργία της αγοράς, θεσπίζουν κανόνες και διασφαλίζουν τον υγιή ανταγωνισμό. Σε ευρωπαϊκό επίπεδο, ο ACER συντονίζει τις εθνικές ρυθμιστικές αρχές. Οι \textbf{Καταναλωτές} περιλαμβάνουν βιομηχανικούς, εμπορικούς και οικιακούς χρήστες. Στις σύγχρονες αγορές, οι καταναλωτές μπορούν επίσης να συμμετέχουν ενεργά στην αγορά μέσω της απόκρισης ζήτησης και της αυτοπαραγωγής.

\subsection{Προθεσμιακή Αγορά - OTC}
\label{subsec:forward-market}

Η Προθεσμιακή Αγορά αποτελεί το πρώτο στάδιο της χονδρικής αγοράς ηλεκτρικής ενέργειας, όπου οι συμμετέχοντες μπορούν να διαπραγματευτούν και να συνάψουν συμβόλαια για μελλοντική παράδοση ηλεκτρικής ενέργειας. Διακρίνεται σε \textbf{Οργανωμένη Προθεσμιακή Αγορά}, η οποία λειτουργεί μέσω χρηματιστηρίων ενέργειας (π.χ. EEX - European Energy Exchange) όπου διαπραγματεύονται τυποποιημένα προθεσμιακά συμβόλαια (futures, forwards, options), και σε \textbf{Εξωχρηματιστηριακή Αγορά} (OTC - Over The Counter), όπου οι συναλλαγές πραγματοποιούνται απευθείας μεταξύ των αντισυμβαλλομένων, με συμβόλαια προσαρμοσμένα στις ανάγκες τους.

Τα βασικά χαρακτηριστικά της Προθεσμιακής Αγοράς περιλαμβάνουν τον χρονικό ορίζοντα (τα συμβόλαια μπορεί να αφορούν παράδοση ηλεκτρικής ενέργειας από μερικές ημέρες έως αρκετά χρόνια στο μέλλον), τη διαχείριση κινδύνου (επιτρέπει στους συμμετέχοντες να αντισταθμίσουν τον κίνδυνο μεταβολής των τιμών), την ευελιξία (ιδιαίτερα στην αγορά OTC, τα συμβόλαια μπορούν να προσαρμοστούν ως προς τον όγκο, τη διάρκεια και τους όρους παράδοσης), τη ρευστότητα (στις οργανωμένες αγορές, τα τυποποιημένα προϊόντα έχουν συνήθως υψηλότερη ρευστότητα) και τον πιστωτικό κίνδυνο (στις συναλλαγές OTC υπάρχει κίνδυνος αθέτησης από τον αντισυμβαλλόμενο).

Η Προθεσμιακή Αγορά διαδραματίζει σημαντικό ρόλο στη διαμόρφωση μακροπρόθεσμων τιμολογιακών σημάτων, συμβάλλει στην οικονομική σταθερότητα των συμμετεχόντων και παρέχει ενδείξεις για μελλοντικές επενδύσεις στον τομέα της ηλεκτρικής ενέργειας.

\subsection{Προημερήσια Αγορά Ηλεκτρικής Ενέργειας}
\label{subsec:day-ahead-market}

Η Προημερήσια Αγορά (Day-Ahead Market - DAM) αποτελεί κεντρικό πυλώνα της χονδρικής αγοράς ηλεκτρικής ενέργειας στην Ευρώπη. Σε αυτή την αγορά, οι συμμετέχοντες υποβάλλουν προσφορές για αγορά και πώληση ηλεκτρικής ενέργειας για κάθε ώρα της επόμενης ημέρας.

Η αγορά λειτουργεί συνήθως με \textbf{μηχανισμό δημοπρασίας μοναδικής τιμής} (uniform pricing auction), όπου όλες οι αποδεκτές προσφορές εκκαθαρίζονται στην οριακή τιμή συστήματος (System Marginal Price - SMP). Όσον αφορά το χρονοδιάγραμμα, οι προσφορές υποβάλλονται μέχρι μια συγκεκριμένη ώρα (π.χ. 12:00 CET) και τα αποτελέσματα ανακοινώνονται λίγο αργότερα, καθορίζοντας το πρόγραμμα παραγωγής και κατανάλωσης για την επόμενη ημέρα. Τα προϊόντα συνήθως αφορούν ωριαίες ή ημίωρες ποσότητες ηλεκτρικής ενέργειας, ενώ σε ορισμένες αγορές υπάρχουν και πιο σύνθετα προϊόντα όπως μπλοκ προσφορές και συνδεδεμένες προσφορές.

Στην Ευρώπη, οι περισσότερες εθνικές προημερήσιες αγορές είναι συζευγμένες μέσω του μηχανισμού \textbf{Price Coupling of Regions} (PCR), που βελτιστοποιεί τη χρήση των διασυνδέσεων και προωθεί τη σύγκλιση των τιμών. Για την επίλυση της συζευγμένης προημερήσιας αγοράς χρησιμοποιείται ο \textbf{αλγόριθμος EUPHEMIA} (EU Pan-European Hybrid Electricity Market Integration Algorithm), ο οποίος λαμβάνει υπόψη τις προσφορές αγοράς και πώλησης, τους περιορισμούς μεταφοράς και τα διαθέσιμα περιθώρια των διασυνδέσεων~\cite{euphemia-public-description}.

Η Προημερήσια Αγορά παρέχει ένα αξιόπιστο σημείο αναφοράς για τις τιμές ηλεκτρικής ενέργειας και μία βάση για τη διαμόρφωση του προγράμματος λειτουργίας του συστήματος. Επίσης αποτελεί έναν μηχανισμό για την αποτελεσματική κατανομή των πόρων και των διασυνοριακών ικανοτήτων μεταφοράς, καθώς και ένα σημαντικό εργαλείο για τη βελτιστοποίηση του κόστους παραγωγής σε πανευρωπαϊκό επίπεδο.

Οι τιμές που προκύπτουν από την Προημερήσια Αγορά αντανακλούν το οριακό κόστος παραγωγής ηλεκτρικής ενέργειας και επηρεάζονται από παράγοντες όπως η ζήτηση, η διαθεσιμότητα των μονάδων παραγωγής, η παραγωγή από ΑΠΕ και οι περιορισμοί του δικτύου.

\subsection{Ενδοημερήσια Αγορά Ηλεκτρικής Ενέργειας}
\label{subsec:intraday-market}

Η Ενδοημερήσια Αγορά (Intraday Market - IDM) επιτρέπει στους συμμετέχοντες να προσαρμόσουν τις θέσεις τους μετά το κλείσιμο της Προημερήσιας Αγοράς και πριν την πραγματική ώρα παράδοσης της ηλεκτρικής ενέργειας. Παρέχει τη δυνατότητα διαπραγμάτευσης ηλεκτρικής ενέργειας σε χρονικό ορίζοντα ``σχεδόν πραγματικού χρόνου''.

Σε αντίθεση με την Προημερήσια Αγορά που λειτουργεί με δημοπρασία, η Ενδοημερήσια Αγορά συνήθως λειτουργεί με μηχανισμό \textbf{συνεχούς διαπραγμάτευσης} (continuous trading), όπου οι συναλλαγές πραγματοποιούνται συνεχώς μέχρι το κλείσιμο της πύλης (gate closure) που είναι συνήθως 5-60 λεπτά πριν την παράδοση. Η αγορά επιτρέπει τη διαχείριση αβεβαιότητας και την αντιμετώπιση απρόβλεπτων γεγονότων όπως βλάβες μονάδων παραγωγής, μεταβολές στην πρόβλεψη παραγωγής ΑΠΕ ή απροσδόκητες διακυμάνσεις στη ζήτηση. Το Ευρωπαϊκό μοντέλο \textbf{XBID} (Cross-Border Intraday) παρέχει μια ενιαία πλατφόρμα για διασυνοριακές ενδοημερήσιες συναλλαγές σε όλη την Ευρώπη, ενισχύοντας τη ρευστότητα και την αποτελεσματικότητα της αγοράς.

Η Ενδοημερήσια Αγορά προσφέρει βελτιωμένη αποτελεσματικότητα του συστήματος, καθώς οι συμμετέχοντες μπορούν να προσαρμόσουν τις θέσεις τους με βάση πιο πρόσφατες πληροφορίες. Αποτελεί ένα σημαντικό εργαλείο για την ενσωμάτωση μεταβλητών ΑΠΕ, επιτρέποντας την αντιμετώπιση σφαλμάτων πρόβλεψης και βοηθάει στη μείωση της ανάγκης για ενεργοποίηση πόρων εξισορρόπησης, καθώς αποκλίσεις μπορούν να διορθωθούν μέσω της αγοράς αυτής. Η σημασία της Ενδοημερήσιας Αγοράς αυξάνεται συνεχώς, ιδιαίτερα με την αυξανόμενη διείσδυση των ΑΠΕ και τις προκλήσεις που αυτή δημιουργεί στην πρόβλεψη της παραγωγής.

\subsection{Αγορά Εξισορρόπησης}
\label{subsec:balancing-market}

Η Αγορά Εξισορρόπησης (Balancing Market) αποτελεί το τελευταίο χρονικά στάδιο των αγορών ηλεκτρικής ενέργειας και λειτουργεί σε πραγματικό ή κοντά σε πραγματικό χρόνο. Ο βασικός της στόχος είναι να διασφαλίσει την ισορροπία μεταξύ παραγωγής και κατανάλωσης ηλεκτρικής ενέργειας ανά πάσα στιγμή, διατηρώντας τη συχνότητα του συστήματος στα επιτρεπτά όρια.

Η αγορά διαχειρίζεται από τον Διαχειριστή του Συστήματος Μεταφοράς (TSO), ο οποίος είναι υπεύθυνος για την εξισορρόπηση του συστήματος σε πραγματικό χρόνο. Τα προϊόντα εξισορρόπησης περιλαμβάνουν τις Εφεδρείες Διατήρησης Συχνότητας (FCR - Frequency Containment Reserves) με αυτόματη απόκριση εντός δευτερολέπτων, τις Εφεδρείες Αποκατάστασης Συχνότητας (FRR - Frequency Restoration Reserves) που διακρίνονται σε αυτόματες (aFRR) και χειροκίνητες (mFRR), καθώς και τις Εφεδρείες Αντικατάστασης (RR - Replacement Reserves) για την αποκατάσταση των προηγούμενων εφεδρειών~\cite{entsoe-balancing}.

Για την αποτελεσματικότερη λειτουργία της αγοράς, έχουν αναπτυχθεί πανευρωπαϊκές πλατφόρμες εξισορρόπησης όπως η πλατφόρμα PICASSO για την ανταλλαγή aFRR, η πλατφόρμα MARI για την ανταλλαγή mFRR και η TERRE για την ανταλλαγή RR. Η τιμολόγηση ενέργειας εξισορρόπησης μπορεί να βασίζεται είτε στην τιμή προσφοράς (pay-as-bid) είτε στην οριακή τιμή (pay-as-cleared), με την τελευταία να προτιμάται στις περισσότερες ευρωπαϊκές αγορές. Οι συμμετέχοντες που αποκλίνουν από το πρόγραμμά τους χρεώνονται ή πιστώνονται με την τιμή απόκλισης (imbalance price).

Η Αγορά Εξισορρόπησης είναι καθοριστικής σημασίας για τη διασφάλιση της αξιοπιστίας και ασφάλειας του συστήματος ηλεκτρικής ενέργειας, την αποτελεσματική διαχείριση των διακυμάνσεων της παραγωγής από ΑΠΕ και την παροχή οικονομικών κινήτρων για επενδύσεις σε ευέλικτους πόρους. Με την αυξανόμενη διείσδυση των ΑΠΕ, η σημασία και η πολυπλοκότητα της Αγοράς Εξισορρόπησης αυξάνεται, απαιτώντας συνεχή εξέλιξη των μηχανισμών της και στενότερη διασυνοριακή συνεργασία μεταξύ των Διαχειριστών Συστημάτων Μεταφοράς.


% =============================================================================
\section{Σκοπός και Στόχοι της Εργασίας}
\label{sec:objectives}

Η παρούσα διπλωματική εργασία έχει ως κεντρικό στόχο την ανάπτυξη ενός ολοκληρωμένου, ανοιχτού κώδικα εργαλείου για την προσομοίωση της Προημερήσιας Αγοράς ηλεκτρικής ενέργειας, με έμφαση στον αλγόριθμο EUPHEMIA που χρησιμοποιείται για την εκκαθάριση της συζευγμένης ευρωπαϊκής αγοράς.

\subsection{Κίνητρο και Αναγκαιότητα}
\label{subsec:motivation}

Η κατανόηση των μηχανισμών λειτουργίας των αγορών ηλεκτρικής ενέργειας αποτελεί κρίσιμο παράγοντα για όλους τους εμπλεκόμενους φορείς, από τους παραγωγούς και τους προμηθευτές έως τους ρυθμιστές και τους ερευνητές. Ωστόσο, η πρόσβαση σε εργαλεία προσομοίωσης που αναπαράγουν με ακρίβεια τη λειτουργία της αγοράς είναι περιορισμένη, καθώς τα υπάρχοντα συστήματα είναι συνήθως κλειστά και ιδιόκτητα.

Η ενεργειακή μετάβαση και η αυξανόμενη διείσδυση των ΑΠΕ δημιουργούν νέες προκλήσεις για τη λειτουργία των αγορών ηλεκτρικής ενέργειας. Η μεταβλητότητα της παραγωγής από ανανεώσιμες πηγές, οι περιορισμοί του δικτύου μεταφοράς και η ανάγκη για διασυνοριακή συνεργασία καθιστούν τη μοντελοποίηση της αγοράς ένα σύνθετο πρόβλημα βελτιστοποίησης. Η δυνατότητα προσομοίωσης διαφορετικών σεναρίων και η κατανόηση των παραγόντων που επηρεάζουν τις τιμές της ηλεκτρικής ενέργειας είναι απαραίτητη για τη λήψη αποφάσεων τόσο σε επιχειρησιακό όσο και σε πολιτικό επίπεδο.

\subsection{Επιστημονική Συνεισφορά}
\label{subsec:contribution}

Η συνεισφορά της παρούσας εργασίας εντοπίζεται σε τρεις κύριους άξονες.

Πρώτον, στην \textbf{ανάπτυξη ολοκληρωμένης υποδομής δεδομένων}. Σχεδιάστηκε και υλοποιήθηκε ένα αυτοματοποιημένο σύστημα συλλογής δεδομένων από την πλατφόρμα διαφάνειας του ENTSO-E, το οποίο αποθηκεύει και διαχειρίζεται δεδομένα για μονάδες παραγωγής, φορτία, προβλέψεις ΑΠΕ, διασυνοριακές μεταφορές και μη διαθεσιμότητες μονάδων. Η υποδομή αυτή επιτρέπει την εκτέλεση προσομοιώσεων με πραγματικά δεδομένα για οποιαδήποτε ζώνη προσφοράς (bidding zone) της Ευρώπης.

Δεύτερον, στην \textbf{υλοποίηση αλγορίθμων βελτιστοποίησης αγοράς}. Αναπτύχθηκε μια σειρά από αλγορίθμους που καλύπτουν το πλήρες φάσμα της λειτουργίας της Προημερήσιας Αγοράς, συμπεριλαμβανομένου του προβλήματος Unit Commitment για τον προσδιορισμό του βέλτιστου προγράμματος λειτουργίας των μονάδων παραγωγής, της μηχανής στρατηγικών υποβολής προσφορών, του βιβλίου εντολών οριακού κόστους (merit-order) που κατασκευάζει ανταγωνιστικό αντιπαράδειγμα της αγοράς από τα θεμελιώδη μεγέθη της (τιμές καυσίμων και CO$_2$, αξία νερού, καθαρές εισαγωγές), του αλγορίθμου MPCC (Mathematical Programming with Complementarity Constraints) για την εκκαθάριση της αγοράς και της μοντελοποίησης δικτύου με περιορισμούς ATC (Available Transfer Capacity) για την πολυζωνική εκκαθάριση της αγοράς.

Τρίτον, στη \textbf{δημιουργία ανοιχτού κώδικα εργαλείου με ποσοτικά επικυρωμένη ακρίβεια}. Ολόκληρο το σύστημα αναπτύχθηκε σε γλώσσα Julia με χρήση του πλαισίου JuMP για τη μοντελοποίηση προβλημάτων βελτιστοποίησης, και ο κώδικας διατίθεται ελεύθερα με άδεια ανοιχτού κώδικα. Το σύστημα αναπαράγει τις τιμές της ελληνικής ζώνης με συντελεστή συσχέτισης 0{,}84 σε ανάλυση δεκαπενταλέπτου επί συνεχούς τετραμήνου (Κεφάλαιο~\ref{ch:results})· εξ όσων είμαστε σε θέση να γνωρίζουμε, πρόκειται για το πρώτο ανοιχτό εργαλείο που μοντελοποιεί την αγορά αυτή με τεκμηριωμένη ακρίβεια αυτού του επιπέδου.

\subsection{Ερευνητικά Ερωτήματα}
\label{subsec:research-questions}

Η εργασία επιχειρεί να απαντήσει στα ακόλουθα ερευνητικά ερωτήματα. Πρώτον, είναι δυνατή η αναπαραγωγή των τιμών της Προημερήσιας Αγοράς με ικανοποιητική ακρίβεια χρησιμοποιώντας δημοσίως διαθέσιμα δεδομένα και αλγορίθμους ανοιχτού κώδικα; Δεύτερον, ποια είναι η επίδραση των διαφορετικών στρατηγικών υποβολής προσφορών στις τελικές τιμές της αγοράς; Τρίτον, πώς επηρεάζουν οι περιορισμοί διασυνοριακής μεταφοράς τη διαμόρφωση των τιμών σε γειτονικές ζώνες προσφοράς;


% =============================================================================
\section{Δομή της Εργασίας}
\label{sec:thesis-structure}

Η παρούσα διπλωματική εργασία οργανώνεται σε έξι κεφάλαια, τα οποία καλύπτουν το θεωρητικό υπόβαθρο, την αρχιτεκτονική του συστήματος, την υλοποίηση, τα αποτελέσματα και τα συμπεράσματα.

Στο \textbf{Κεφάλαιο~\ref{ch:introduction}} (Εισαγωγή), το παρόν κεφάλαιο, παρουσιάζεται το πλαίσιο της εργασίας με αναφορά στην κλιματική αλλαγή, την ενεργειακή μετάβαση και τη δομή των ευρωπαϊκών αγορών ηλεκτρικής ενέργειας. Επίσης, καθορίζονται οι στόχοι και η συνεισφορά της εργασίας.

Στο \textbf{Κεφάλαιο~\ref{ch:theoretical-background}} (Θεωρητικό Υπόβαθρο) αναλύονται οι θεωρητικές βάσεις της εργασίας, συμπεριλαμβανομένων των μεθόδων μαθηματικού προγραμματισμού, του προβλήματος Unit Commitment, του αλγορίθμου EUPHEMIA και των στρατηγικών υποβολής προσφορών.

Στο \textbf{Κεφάλαιο~\ref{ch:system-architecture}} (Αρχιτεκτονική Συστήματος και Δεδομένα) περιγράφεται η αρχιτεκτονική του συστήματος που αναπτύχθηκε, με έμφαση στο pipeline συλλογής δεδομένων από το ENTSO-E, το σχήμα της βάσης δεδομένων και το οικοσύστημα Julia/JuMP.

Στο \textbf{Κεφάλαιο~\ref{ch:implementation}} (Υλοποίηση) παρουσιάζεται αναλυτικά η υλοποίηση των αλγορίθμων, συμπεριλαμβανομένου του solver για το πρόβλημα Unit Commitment, της μηχανής στρατηγικών υποβολής προσφορών, του βιβλίου εντολών οριακού κόστους, του αλγορίθμου MPCC για την εκκαθάριση της αγοράς και της μοντελοποίησης των περιορισμών δικτύου.

Στο \textbf{Κεφάλαιο~\ref{ch:results}} (Αποτελέσματα και Αξιολόγηση) παρουσιάζονται τα αποτελέσματα της εφαρμογής του συστήματος σε πραγματικά δεδομένα, η σύγκριση με τις πραγματικές τιμές της αγοράς και η ανάλυση της απόδοσης του συστήματος.

Στο \textbf{Κεφάλαιο~\ref{ch:conclusions}} (Συμπεράσματα) συνοψίζονται τα κύρια ευρήματα της εργασίας, αναφέρονται οι περιορισμοί και προτείνονται κατευθύνσεις για μελλοντική έρευνα.

Τέλος, στο \textbf{Παράρτημα} παρατίθενται συμπληρωματικά στοιχεία και σύνδεσμοι προς τον κώδικα της εφαρμογής.

% =============================================================================

\chapter{Εισαγωγή}
\label{ch:introduction}

\section{Κλιματική Αλλαγή και μέτρα για την αντιμετώπισή της}
\label{sec:climate-change}

Η Κλιματική Αλλαγή αναδεικνύεται ως η κορυφαία πρόκληση της εποχής μας, συνιστώντας μια υπαρξιακή απειλή για το φυσικό περιβάλλον, την συνοχή των ανθρώπινων κοινωνιών και την παγκόσμια οικονομική σταθερότητα. Η επιστημονική κοινότητα συμφωνεί ότι η ραγδαία αύξηση της μέσης θερμοκρασίας του πλανήτη και η εντατικοποίηση των ακραίων καιρικών φαινομένων οφείλονται, κατά κύριο λόγο, στις ανθρωπογενείς εκπομπές αερίων του θερμοκηπίου~\cite{lynas2021}.

Μέρος αυτών των εκπομπών προέρχεται από την χρήση ορυκτών καυσίμων όπως ο άνθρακας, το πετρέλαιο, ο λιγνίτης και το φυσικό αέριο, για την παραγωγή ηλεκτρικής ενέργειας. Η πρωτοκαθεδρία των ορυκτών καυσίμων παραμένει αδιαμφισβήτητη, αντιπροσωπεύοντας συνολικά σχεδόν το 60\% της παγκόσμιας παραγωγής ηλεκτρικής ενέργειας για το έτος 2024. Η εν λόγω αναλογία υπογραμμίζει την συνεχιζόμενη εξάρτηση των παγκόσμιων ενεργειακών συστημάτων από συμβατικές πηγές που επιτείνουν την κλιματική αλλαγή και τις επιπτώσεις αυτής. Αξίζει όμως να σημειωθεί ότι το ενεργειακό μίγμα βελτιώνεται σταδιακά, καθότι οι Ανανεώσιμες Πηγές Ενέργειας (ΑΠΕ) ανήλθαν συνολικά στο ένα τρίτο (33\%) της ηλεκτροπαραγωγής, όπου μαζί με την πυρηνική ενέργεια (9\%), άγγιξαν για πρώτη φορά το 40\% της παγκόσμιας ηλεκτροπαραγωγής ενέργειας~\cite{iea2024}.

\subsection{Παγκόσμιο επίπεδο}
\label{subsec:global-level}

Η Σύμβαση-Πλαίσιο των Ηνωμένων Εθνών για την Κλιματική Αλλαγή (UNFCCC), που συμφωνήθηκε το 1992 στη Διάσκεψη του Ρίο, αποτελεί τη βασική διεθνή συνθήκη για την αντιμετώπιση της κλιματικής αλλαγής. Στόχος της είναι η πρόληψη της επικίνδυνης ανθρωπογενούς παρέμβασης στο παγκόσμιο κλιματικό σύστημα. Από αυτήν, προέκυψαν συμφωνίες όπως το Πρωτόκολλο του Κιότο (1997), το οποίο εισήγαγε νομικά δεσμευτικούς στόχους μείωσης εκπομπών για τις ανεπτυγμένες χώρες και καθιέρωσε μηχανισμούς όπως η εμπορία δικαιωμάτων εκπομπών (Emissions Trading), ο μηχανισμός ``καθαρής'' ανάπτυξης (Clean Development Mechanism) και ο μηχανισμός κοινής εφαρμογής (Joint Implementation)~\cite{ypen-kyoto}. Το Πρωτόκολλο του Κιότο ήταν πρωτοπόρο για την εποχή, αλλά η αποτελεσματικότητά του περιορίστηκε από τη μη συμμετοχή ή αποχώρηση μεγάλων εκπομπέων όπως οι ΗΠΑ, η Ιαπωνία και ο Καναδάς~\cite{bassetti2022}.

Το 2015, η Συμφωνία του Παρισιού υιοθετήθηκε ως η νέα παγκόσμια συμφωνία για το κλίμα. Ο πρωταρχικός της στόχος είναι να συγκρατήσει ``την αύξηση της μέσης παγκόσμιας θερμοκρασίας πολύ κάτω από τους 2°C σε σχέση με τα προβιομηχανικά επίπεδα'' και να καταβάλει προσπάθειες ``να περιορίσει την αύξηση της θερμοκρασίας στον 1,5°C σε σχέση με τα προβιομηχανικά επίπεδα''~\cite{unfccc-paris}.

\subsection{Ευρωπαϊκό επίπεδο}
\label{subsec:european-level}

Πέραν των ανωτέρω συμφωνιών, η αντιμετώπιση της κλιματικής αλλαγής σε ευρωπαϊκό επίπεδο και πιο συγκεκριμένα όσον αφορά την Ευρωπαϊκή Ένωση, εστιάζει γύρω από τον στόχο της κλιματικής ουδετερότητας έως το 2050 και τις ενδιάμεσες δεσμεύσεις μείωσης των καθαρών εκπομπών τουλάχιστον 55\% έως το 2030 και τουλάχιστον 90\% έως το 2040, σε σχέση με τα επίπεδα του 1990. Η πρωτοβουλία-πλαίσιο που συγκεντρώνει αυτά τα πολιτικά μέτρα είναι το Ευρωπαϊκό Σχέδιο για το Περιβάλλον (European Green Deal), το οποίο προβλέπει αναθεώρηση της υπάρχουσας νομοθεσίας και εισαγωγή νέων πολιτικών σε τομείς όπως η ενεργειακή μετάβαση, η κυκλική οικονομία, η ανακαίνιση κτιρίων, η βιοποικιλότητα και η αγροτική πολιτική~\cite{ec-green-deal}.

Το πακέτο πολιτικών του Green Deal συνδέει ρητά την ενεργειακή αποδέσμευση από τα ορυκτά καύσιμα με μέτρα οικονομικής στήριξης, επενδύσεις σε τεχνολογικές καινοτομίες και κοινωνικά μέτρα για την αντιμετώπιση των επιπτώσεων στη συνοχή των περιοχών και των εργαζομένων. Η νομοθετική δέσμευση ενισχύεται από τον Ευρωπαϊκό Νόμο για το Κλίμα, ο οποίος καθιστά τον στόχο της ουδετερότητας νομικά δεσμευτικό και απαιτεί σαφείς ενδιάμεσους στόχους και μηχανισμούς παρακολούθησης για τα κράτη-μέλη~\cite{eu-climate-law}.

Ως αποτέλεσμα, η ΕΕ με το Green Deal συνδυάζει ρυθμιστικά εργαλεία (όπως αναθεωρήσεις αγοράς δικαιωμάτων εκπομπών, πρότυπα ενεργειακής απόδοσης και κανονισμούς για τη μεταφορά και τη βιομηχανία) με χρηματοδοτικά μέσα (ταμεία δίκαιης μετάβασης, προγράμματα επενδύσεων) και πολιτικές υποστήριξης των ΑΠΕ και της έρευνας. Η ευρωπαϊκή προσέγγιση προωθεί τόσο τη μείωση εκπομπών όσο και την ενίσχυση της κλιματικής ανθεκτικότητας των υποδομών και των κοινοτήτων, με έμφαση στη συνοχή των πολιτικών σε όλα τα επίπεδα διακυβέρνησης~\cite{sikora2020}.

\subsection{Εθνικό επίπεδο}
\label{subsec:national-level}

Σε εθνικό επίπεδο, η Ελλάδα, ως μέλος της ΕΕ, έχει υιοθετήσει ένα πλαίσιο πολιτικών και στρατηγικών που ευθυγραμμίζονται πλήρως με τους στόχους της Ευρωπαϊκής Ένωσης για την κλιματική ουδετερότητα έως το 2050 και τη μείωση των εκπομπών έως το 2030, ενσωματώνοντας μέτρα για την ενεργειακή μετάβαση, την απολιγνιτοποίηση, την επιτάχυνση των ΑΠΕ και την αναβάθμιση της ενεργειακής απόδοσης των κτιρίων.

Το παραπάνω πλαίσιο είναι το αναθεωρημένο Εθνικό Σχέδιο για την Ενέργεια και το Κλίμα (ΕΣΕΚ) το οποίο κυρώθηκε το 2024. Σε αυτό το έγγραφο αναγνωρίζονται οι επιπτώσεις που η κλιματική αλλαγή θα επιφέρει ειδικά στην Ελλάδα, όπως η ``μειωμένη διαθεσιμότητα υδάτων'' η οποία ``θέτει σε κίνδυνο τη λειτουργία θερμοηλεκτρικών μονάδων. Τα δίκτυα μεταφοράς και διανομής ηλεκτρικής ενέργειας, οι σταθμοί υψηλής τάσης, και λοιπά ενεργειακά δίκτυα και εγκαταστάσεις είναι τρωτά σε ακραία καιρικά φαινόμενα, ενώ οι παράκτιες ενεργειακές υποδομές, απειλούνται από την άνοδο της στάθμης της θάλασσας. Η αύξηση της μέσης θερμοκρασίας αναμένεται να μειώσει τις ενεργειακές ανάγκες για θέρμανση κατά τη χειμερινή περίοδο και να αυξήσει τις ανάγκες για ψύξη κατά τη θερινή περίοδο. Οι μεταβολές στη ζήτηση θα οδηγήσουν σε μεγαλύτερη διακύμανση φορτίων δοκιμάζοντας το σύστημα ηλεκτρικής ενέργειας. Επίσης, η κλιματική αλλαγή αναμένεται να επηρεάσει την απόδοση των συστημάτων ΑΠΕ''~\cite{esek2024}.

\subsection{Τεχνητή νοημοσύνη και ενεργειακή κατανάλωση}
\label{subsec:ai-energy}

Η ραγδαία επέκταση της Τεχνητής Νοημοσύνης (AI) και του υπολογιστικού νέφους (cloud computing) θέτει μία αυξανόμενη πρόκληση για την κλιματική αλλαγή, καθώς οδηγεί σε πρωτοφανή αύξηση της κατανάλωσης ενέργειας των κέντρων δεδομένων (data centers). Τα κέντρα δεδομένων, τα οποία ήδη καταναλώνουν περίπου το 1\% της παγκόσμιας ηλεκτρικής ενέργειας, αναμένεται να διπλασιάσουν αυτόν τον αριθμό έως το 2030, κυρίως λόγω των εξαιρετικά απαιτητικών εργασιών της τεχνητής νοημοσύνης, όπως η εκπαίδευση μεγάλων γλωσσικών μοντέλων (Large Language Models - LLMs). 

Η εκπαίδευση και η χρήση αυτών των μοντέλων απαιτεί την εντατική χρήση ενεργοβόρου υλικού όπως GPUs και TPUs, καθώς και μεγάλη ανάγκη για συστήματα ψύξης, τα οποία ευθύνονται για σχεδόν το 40\% της συνολικής ενεργειακής κατανάλωσης. Αυτό σημαίνει ότι η περιβαλλοντική επίπτωση αυξάνεται σημαντικά, εκτός εάν υπάρξει μια σημαντική στροφή σε ανανεώσιμες πηγές ενέργειας και προηγμένες τεχνικές ψύξης. Κατά συνέπεια, η έγκαιρη αντιμετώπιση της ενεργειακής ζήτησης της AI είναι κρίσιμη για τη διασφάλιση της βιωσιμότητας της ψηφιακής υποδομής και τον μετριασμό των αρνητικών επιπτώσεων στο περιβάλλον~\cite{chauhan2024}.


% =============================================================================
\section{Συστήματα Ηλεκτρικής Ενέργειας}
\label{sec:power-systems}

Τα συστήματα ηλεκτρικής ενέργειας αποτελούν το σύνολο των εγκαταστάσεων και μέσων για την παροχή ηλεκτρικής ενέργειας σε εξυπηρετούμενες περιοχές κατανάλωσης. Για την εύρυθμη λειτουργία ενός συστήματος ηλεκτρικής ενέργειας είναι απαραίτητη η παροχή ηλεκτρικής ενέργειας οπουδήποτε υπάρχει ζήτηση με το ελάχιστο δυνατό κόστος και τις ελάχιστες οικολογικές επιπτώσεις, εξασφαλίζοντας σταθερή συχνότητα, σταθερή τάση και υψηλή αξιοπιστία τροφοδότησης~\cite{vournas-kontaxis}.

Η τροφοδότηση των καταναλωτών με ηλεκτρική ενέργεια προϋποθέτει τρεις ξεχωριστές λειτουργίες του συστήματος ηλεκτρικής ενέργειας: την παραγωγή, τη μεταφορά, και τη διανομή. Η παραγόμενη ηλεκτρική ενέργεια πρέπει ιδανικά να απορροφάται από την κατανάλωση σε ίδιο χρόνο, καθώς η αποθήκευσή της παραμένει δύσκολη και τις περισσότερες φορές οικονομικά ασύμφορη, παρά την πρόοδο στην τεχνολογία αποθήκευσης ηλεκτρικής ενέργειας τα τελευταία χρόνια~\cite{hiesl2020}.

% TODO: Εισαγωγή σχήματος συστήματος ηλεκτρικής ενέργειας
\begin{figure}[htbp]
\centering
% \includegraphics[width=0.9\textwidth]{figures/power_system_overview.pdf}
\fbox{\parbox{0.8\textwidth}{\centering\vspace{2cm}[Σχήμα: Γενική δομή συστήματος ηλεκτρικής ενέργειας]\vspace{2cm}}}
\caption{Γενική δομή συστήματος ηλεκτρικής ενέργειας με τις τρεις βασικές λειτουργίες: παραγωγή, μεταφορά και διανομή.}
\label{fig:power-system-overview}
\end{figure}

\subsection{Παραγωγή Ηλεκτρικής Ενέργειας}
\label{subsec:generation}

Η παραγωγή ηλεκτρικής ενέργειας αφορά τη μετατροπή διαφόρων μορφών πρωτογενούς ενέργειας σε ηλεκτρική. Τα εργοστάσια παραγωγής διακρίνονται με βάση την τεχνολογία και την πηγή ενέργειας που χρησιμοποιούν.

Οι \textbf{θερμοηλεκτρικοί σταθμοί παραγωγής} αποτελούν τη συμβατική μορφή παραγωγής ηλεκτρικής ενέργειας. Σε αυτούς, η θερμική ενέργεια που προέρχεται από την καύση ορυκτών καυσίμων (λιγνίτης, φυσικό αέριο, πετρέλαιο) ή από πυρηνικές αντιδράσεις μετατρέπεται σε μηχανική μέσω ατμοστροβίλων ή αεριοστροβίλων, και στη συνέχεια σε ηλεκτρική μέσω γεννητριών. Οι μονάδες αυτές χαρακτηρίζονται από υψηλή διαθεσιμότητα και ελεγχόμενη παραγωγή, αλλά συνεπάγονται σημαντικές εκπομπές αερίων του θερμοκηπίου.

Οι \textbf{υδροηλεκτρικοί σταθμοί} αξιοποιούν τη δυναμική ενέργεια του νερού για την κίνηση υδροστροβίλων. Διακρίνονται σε συμβατικούς (με φράγμα και ταμιευτήρα) και αντλησιοταμιευτικούς, οι οποίοι λειτουργούν και ως συστήματα αποθήκευσης ενέργειας.

Οι \textbf{Ανανεώσιμες Πηγές Ενέργειας (ΑΠΕ)} περιλαμβάνουν τεχνολογίες όπως τα αιολικά πάρκα, τα φωτοβολταϊκά συστήματα, τη γεωθερμία και τη βιομάζα. Οι ΑΠΕ χαρακτηρίζονται από μηδενικές ή ελάχιστες εκπομπές κατά τη λειτουργία τους, αλλά η παραγωγή τους εξαρτάται από τις καιρικές συνθήκες, καθιστώντας τη μεταβλητή και εν μέρει απρόβλεπτη.

\subsection{Μεταφορά Ηλεκτρικής Ενέργειας}
\label{subsec:transmission}

Η μεταφορά ηλεκτρικής ενέργειας αφορά τη μεταφορά μεγάλων ποσοτήτων ηλεκτρικής ενέργειας από τους σταθμούς παραγωγής προς τα κέντρα κατανάλωσης μέσω του δικτύου υψηλής τάσης. Η υψηλή τάση (150kV, 400kV) χρησιμοποιείται για τη μείωση των απωλειών μεταφοράς, καθώς για δεδομένη ισχύ, η αύξηση της τάσης οδηγεί σε μείωση του ρεύματος και συνεπώς των απωλειών λόγω θέρμανσης των αγωγών.

Αρμόδιος για την αξιόπιστη μεταφορά ηλεκτρικής ενέργειας στην ελληνική επικράτεια είναι ο Ανεξάρτητος Διαχειριστής Μεταφοράς Ηλεκτρικής Ενέργειας (ΑΔΜΗΕ). Ο ΑΔΜΗΕ διαχειρίζεται το Διασυνδεδεμένο Σύστημα Μεταφοράς, το οποίο περιλαμβάνει βασικά στοιχεία όπως εναέριες γραμμές μεταφοράς 400kV, 150kV και 66kV, υπόγειες και υποβρύχιες καλωδιακές γραμμές 150kV και 400kV (AC και DC), υποσταθμούς 150/20kV, καθώς και Κέντρα Υπερυψηλής Τάσης (ΚΥΤ) 400/150kV.

Το Ελληνικό Σύστημα Μεταφοράς λειτουργεί παράλληλα με το διασυνδεδεμένο Ευρωπαϊκό Σύστημα, υπό τον ευρύτερο συντονισμό του Ευρωπαϊκού Δικτύου Διαχειριστών Συστημάτων Μεταφοράς Ηλεκτρικής Ενέργειας (European Network of Transmission System Operators for Electricity - ENTSO-E). Η σύνδεση αυτή πραγματοποιείται μέσω διασυνδετικών γραμμών μεταφοράς 400kV, που ενώνουν το ελληνικό σύστημα με εκείνα της Αλβανίας, της Βουλγαρίας, της Βόρειας Μακεδονίας και της Τουρκίας. Επιπλέον, υπάρχει ασύγχρονη διασύνδεση με την Ιταλία μέσω υποβρυχίου συνδέσμου συνεχούς ρεύματος τάσης 400kV~\cite{admie-system}.

\subsection{Διανομή Ηλεκτρικής Ενέργειας}
\label{subsec:distribution}

Από τους υποσταθμούς υποβιβασμού τάσης ξεκινά η διανομή ηλεκτρικής ενέργειας, το σύνολο των διαδικασιών που απαιτούνται για τον αξιόπιστο διαμοιρασμό της στους μετρητές των καταναλωτών, διατηρώντας την τάση στους ζυγούς και τις ροές στις γραμμές του δικτύου εντός επιτρεπτών ορίων.

Αρμόδιος για τη διαχείριση του δικτύου διανομής ηλεκτρικής ενέργειας στην ελληνική επικράτεια είναι ο Διαχειριστής Ελληνικού Δικτύου Διανομής Ηλεκτρικής Ενέργειας (ΔΕΔΔΗΕ). Ο ΔΕΔΔΗΕ είναι επίσης υπεύθυνος για τη διαχείριση των αγορών των Μη Διασυνδεδεμένων Νησιών (ΜΔΝ)~\cite{nomos4001}.

Εντός του δικτύου διανομής ηλεκτρικής ενέργειας έχει ξεκινήσει και η ανάπτυξη των τεχνολογιών έξυπνων δικτύων (smart grids) με χρήση έξυπνων μετρητών (smart meters), συστημάτων εποπτικού ελέγχου (SCADA) και ολοκληρωμένων πληροφοριακών συστημάτων διανομής (DMS), που λειτουργούν καταλυτικά στην εκμετάλλευση των οικιακών εγκαταστάσεων ΑΠΕ, συστημάτων αποθήκευσης καθώς και τις τεχνολογίες Vehicle to Grid (V2G)~\cite{fang2012}.

Μέσω της αμφίδρομης ροής πληροφοριών και ενέργειας, τα έξυπνα δίκτυα επιτρέπουν στους διαχειριστές δικτύου να παρακολουθούν σε πραγματικό χρόνο την κατανάλωση, να εντοπίζουν βλάβες άμεσα και να διαχειρίζονται καλύτερα την κατανομή φορτίου. Επιπλέον, παρέχουν δυνατότητες συμμετοχής των τελικών καταναλωτών στην αγορά μέσω demand response, όπου οι χρήστες προσαρμόζουν τη ζήτησή τους ανάλογα με τα σήματα της αγοράς ή τις τιμές ενέργειας. Η μετάβαση προς τα έξυπνα δίκτυα είναι αναγκαία για την υλοποίηση των στόχων της ενεργειακής μετάβασης, την ενίσχυση της ενεργειακής απόδοσης και την επίτευξη ενός βιώσιμου, ψηφιοποιημένου και αποκεντρωμένου ενεργειακού συστήματος.


% =============================================================================
\section{Αγορές Ηλεκτρικής Ενέργειας}
\label{sec:energy-markets}

Για την εξασφάλιση της ισορροπίας προσφοράς-ζήτησης και της απρόσκοπτης παροχής ενέργειας στους καταναλωτές στο ελάχιστο δυνατό κόστος, τα συστήματα ηλεκτρικής ενέργειας ενθυλακώνονται σε αγορές ηλεκτρικής ενέργειας. Οι αγορές ηλεκτρικής ενέργειας χαρακτηρίζονται από ιδιαιτερότητες σε σχέση με άλλες αγορές εμπορευμάτων ή καταναλωτικών αγαθών, λόγω του υψηλού κόστους αποθήκευσης της ηλεκτρικής ενέργειας~\cite{mo2025}.

% TODO: Εισαγωγή σχήματος δομής αγοράς
\begin{figure}[htbp]
\centering
% \includegraphics[width=0.9\textwidth]{figures/market_structure.pdf}
\fbox{\parbox{0.8\textwidth}{\centering\vspace{2cm}[Σχήμα: Δομή αγορών ηλεκτρικής ενέργειας]\vspace{2cm}}}
\caption{Χρονική αλληλουχία των αγορών ηλεκτρικής ενέργειας, από την προθεσμιακή αγορά έως την αγορά εξισορρόπησης.}
\label{fig:market-structure}
\end{figure}

\subsection{Μονοπωλιακή Αγορά}
\label{subsec:monopoly-market}

Ιστορικά, η ανάπτυξη του τομέα της ηλεκτρικής ενέργειας απαιτούσε τεράστια κεφαλαιουχική δαπάνη, μακροχρόνιο σχεδιασμό και υψηλό επίπεδο τεχνικής εξειδίκευσης, στοιχεία που καθιστούσαν πρακτικά αδύνατη τη συμμετοχή ιδιωτών. Για τον λόγο αυτό, το κράτος ή κρατικά ελεγχόμενες επιχειρήσεις αναλάμβαναν την ευθύνη για την κατασκευή, λειτουργία και διαχείριση του συστήματος παραγωγής, μεταφοράς και διανομής ηλεκτρικής ενέργειας. Το μοντέλο που διαμορφώθηκε ήταν μονοπωλιακό, με χαρακτηριστικό γνώρισμα την κάθετη ολοκλήρωση, δηλαδή τον έλεγχο όλων των σταδίων της αλυσίδας αξίας --- από την παραγωγή έως την προμήθεια στον τελικό καταναλωτή --- από έναν και μόνο φορέα. Στην ελληνική επικράτεια το ρόλο αυτό είχε αναλάβει η Δημόσια Επιχείρηση Ηλεκτρισμού (ΔΕΗ).

Σε αυτό το πλαίσιο απουσίας ανταγωνισμού οι καταναλωτές δεν είχαν τη δυνατότητα επιλογής παρόχου. Οι τιμές της ηλεκτρικής ενέργειας καθορίζονταν διοικητικά από τις αρμόδιες ρυθμιστικές αρχές, με βάση το κόστος και ένα σταθερό περιθώριο κέρδους. Επιπλέον, οι επενδύσεις και η επέκταση του δικτύου σχεδιάζονταν κεντρικά, με στόχο την κάλυψη της ζήτησης και την ασφάλεια εφοδιασμού. Παρά τα πλεονεκτήματα σταθερότητας και ελέγχου, το μονοπωλιακό αυτό μοντέλο συχνά οδηγούσε σε αναποτελεσματικότητα, καθυστερήσεις στην υιοθέτηση καινοτομιών και υψηλότερο κόστος για τον τελικό καταναλωτή~\cite{vlados2021}.

Συγκεκριμένα, η ρύθμιση ενός μονοπωλιακού φορέα, όπως μια κρατική ή ρυθμιζόμενη επιχείρηση ηλεκτρικής ενέργειας, αποτελεί χαρακτηριστικό πρόβλημα εντολέα-εντολοδόχου (principal-agent). Ο ρυθμιστής (εντολέας) επιδιώκει τη μεγιστοποίηση του κοινωνικού πλεονάσματος, ενώ η επιχείρηση (εντολοδόχος) στοχεύει στη μεγιστοποίηση του κέρδους. Ωστόσο, ο ρυθμιστής δεν διαθέτει πλήρη πληροφόρηση για το κόστος και τα πραγματικά δεδομένα λειτουργίας της επιχείρησης, η οποία έχει κίνητρο να τα υπερεκτιμά, προβάλλοντας αυξημένες ανάγκες για ασφάλεια και επενδύσεις. Η ασυμμετρία πληροφόρησης αυτή δημιουργεί προβλήματα όπως ηθικός κίνδυνος και αρνητική επιλογή, οδηγώντας σε αποκλίσεις από την αποδοτικότητα του τέλειου ανταγωνισμού~\cite{grossman-hart1992}.

Για την αντιμετώπιση αυτών των στρεβλώσεων, ο ρυθμιστής μπορεί είτε να μειώσει την ασυμμετρία πληροφόρησης, μέσω ελέγχων ή συγκριτικής αξιολόγησης (benchmarking), είτε να εφαρμόσει στρατηγικές κινήτρων, ώστε να ενισχύσει την αποδοτικότητα της επιχείρησης. Παρά τις προσπάθειες αυτές, η ισορροπία που επιτυγχάνεται μεταξύ ρυθμιστή και επιχείρησης δεν φτάνει ποτέ το επίπεδο αποδοτικότητας που προκύπτει σε καθεστώς τέλειου ανταγωνισμού, με αποτέλεσμα την ύπαρξη κόστους πρακτόρευσης (agency cost) για την κοινωνία~\cite{panda-leepsa2017}.

\subsection{Άνοιγμα της αγοράς}
\label{subsec:market-liberalization}

Για την αντιμετώπιση των προβλημάτων αποδοτικότητας που προκύπτουν λόγω του κόστους πρακτόρευσης στο μονοπωλιακό μοντέλο αγοράς, από τη δεκαετία του 1990 η Ευρωπαϊκή Ένωση ξεκίνησε μια διαδικασία απελευθέρωσης των αγορών ηλεκτρικής ενέργειας, με στόχο τη δημιουργία μιας ενιαίας, ανταγωνιστικής αγοράς. Τα κύρια στάδια αυτής της διαδικασίας περιλαμβάνουν το Πρώτο Ενεργειακό Πακέτο (1996), όπου η Οδηγία 96/92/ΕΚ έθεσε τα θεμέλια για το άνοιγμα των αγορών, εισάγοντας την έννοια του διαχωρισμού των δραστηριοτήτων και την πρόσβαση τρίτων στο δίκτυο. Το Δεύτερο Ενεργειακό Πακέτο (2003) με τις Οδηγίες 2003/54/ΕΚ και 2003/55/ΕΚ επιτάχυνε το άνοιγμα της αγοράς, καθιστώντας υποχρεωτικό το νομικό διαχωρισμό των δραστηριοτήτων μεταφοράς και διανομής από την παραγωγή και προμήθεια. Το Τρίτο Ενεργειακό Πακέτο (2009) εισήγαγε αυστηρότερους κανόνες για το διαχωρισμό της ιδιοκτησίας, ενίσχυσε τις αρμοδιότητες των εθνικών ρυθμιστικών αρχών και ίδρυσε τον Οργανισμό Συνεργασίας των Ρυθμιστικών Αρχών Ενέργειας (ACER). Τέλος, το Τέταρτο Ενεργειακό Πακέτο (2019), γνωστό και ως ``Καθαρή Ενέργεια για όλους τους Ευρωπαίους'', εστιάζει στην ενσωμάτωση των ανανεώσιμων πηγών ενέργειας, την ενεργειακή αποδοτικότητα και νέους κανόνες για την αγορά ηλεκτρικής ενέργειας~\cite{eikeland-duffield2011, eu-energy-packages}.

Με βάση τα παραπάνω, το άνοιγμα της αγοράς οδήγησε στον ιδιοκτησιακό διαχωρισμό των δραστηριοτήτων, τη δημιουργία ανταγωνιστικών αγορών χονδρικής και λιανικής, την ελευθερία επιλογής προμηθευτή για τους καταναλωτές και την ανάπτυξη νέων επιχειρηματικών μοντέλων και υπηρεσιών.

\subsection{Φορείς της αγοράς}
\label{subsec:market-participants}

Στο σύγχρονο μοντέλο των απελευθερωμένων αγορών ηλεκτρικής ενέργειας, διάφοροι φορείς διαδραματίζουν καθοριστικό ρόλο. Οι \textbf{Παραγωγοί} είναι εταιρείες που διαθέτουν και λειτουργούν μονάδες παραγωγής ηλεκτρικής ενέργειας, περιλαμβάνοντας τόσο συμβατικούς παραγωγούς (θερμικές μονάδες) όσο και παραγωγούς ΑΠΕ. Ο \textbf{Διαχειριστής Συστήματος Μεταφοράς} (TSO - Transmission System Operator) είναι υπεύθυνος για τη λειτουργία, συντήρηση και ανάπτυξη του συστήματος μεταφοράς υψηλής τάσης, καθώς και για την εξισορρόπηση του συστήματος σε πραγματικό χρόνο. Για την Ελλάδα, αυτός ο φορέας είναι ο ΑΔΜΗΕ. Ο \textbf{Διαχειριστής Δικτύου Διανομής} (DSO - Distribution System Operator) είναι υπεύθυνος για τη λειτουργία, συντήρηση και ανάπτυξη του δικτύου διανομής μέσης και χαμηλής τάσης. Για την Ελλάδα, αυτός ο φορέας είναι ο ΔΕΔΔΗΕ.

Οι \textbf{Προμηθευτές} αγοράζουν ηλεκτρική ενέργεια από την αγορά χονδρικής και την πωλούν στους τελικούς καταναλωτές. Οι \textbf{Έμποροι} (Traders) αγοράζουν και πωλούν ηλεκτρική ενέργεια στις χονδρικές αγορές με σκοπό το κέρδος, χωρίς να είναι απαραίτητα παραγωγοί ή προμηθευτές. Τα \textbf{Χρηματιστήρια Ενέργειας} (Power Exchanges) είναι οργανωμένες αγορές όπου διαπραγματεύονται προϊόντα ηλεκτρικής ενέργειας. Παραδείγματα περιλαμβάνουν το EPEX SPOT, το Nord Pool, το GME και για την Ελλάδα το HENEX (EnExGroup). Ο \textbf{Διαχειριστής Αγοράς} είναι υπεύθυνος για την οργάνωση και λειτουργία των διαφόρων τμημάτων της αγοράς ηλεκτρικής ενέργειας. Οι \textbf{Ρυθμιστικές Αρχές} εποπτεύουν τη λειτουργία της αγοράς, θεσπίζουν κανόνες και διασφαλίζουν τον υγιή ανταγωνισμό. Σε ευρωπαϊκό επίπεδο, ο ACER συντονίζει τις εθνικές ρυθμιστικές αρχές. Οι \textbf{Καταναλωτές} περιλαμβάνουν βιομηχανικούς, εμπορικούς και οικιακούς χρήστες. Στις σύγχρονες αγορές, οι καταναλωτές μπορούν επίσης να συμμετέχουν ενεργά στην αγορά μέσω της απόκρισης ζήτησης και της αυτοπαραγωγής.

\subsection{Προθεσμιακή Αγορά - OTC}
\label{subsec:forward-market}

Η Προθεσμιακή Αγορά αποτελεί το πρώτο στάδιο της χονδρικής αγοράς ηλεκτρικής ενέργειας, όπου οι συμμετέχοντες μπορούν να διαπραγματευτούν και να συνάψουν συμβόλαια για μελλοντική παράδοση ηλεκτρικής ενέργειας. Διακρίνεται σε \textbf{Οργανωμένη Προθεσμιακή Αγορά}, η οποία λειτουργεί μέσω χρηματιστηρίων ενέργειας (π.χ. EEX - European Energy Exchange) όπου διαπραγματεύονται τυποποιημένα προθεσμιακά συμβόλαια (futures, forwards, options), και σε \textbf{Εξωχρηματιστηριακή Αγορά} (OTC - Over The Counter), όπου οι συναλλαγές πραγματοποιούνται απευθείας μεταξύ των αντισυμβαλλομένων, με συμβόλαια προσαρμοσμένα στις ανάγκες τους.

Τα βασικά χαρακτηριστικά της Προθεσμιακής Αγοράς περιλαμβάνουν τον χρονικό ορίζοντα (τα συμβόλαια μπορεί να αφορούν παράδοση ηλεκτρικής ενέργειας από μερικές ημέρες έως αρκετά χρόνια στο μέλλον), τη διαχείριση κινδύνου (επιτρέπει στους συμμετέχοντες να αντισταθμίσουν τον κίνδυνο μεταβολής των τιμών), την ευελιξία (ιδιαίτερα στην αγορά OTC, τα συμβόλαια μπορούν να προσαρμοστούν ως προς τον όγκο, τη διάρκεια και τους όρους παράδοσης), τη ρευστότητα (στις οργανωμένες αγορές, τα τυποποιημένα προϊόντα έχουν συνήθως υψηλότερη ρευστότητα) και τον πιστωτικό κίνδυνο (στις συναλλαγές OTC υπάρχει κίνδυνος αθέτησης από τον αντισυμβαλλόμενο).

Η Προθεσμιακή Αγορά διαδραματίζει σημαντικό ρόλο στη διαμόρφωση μακροπρόθεσμων τιμολογιακών σημάτων, συμβάλλει στην οικονομική σταθερότητα των συμμετεχόντων και παρέχει ενδείξεις για μελλοντικές επενδύσεις στον τομέα της ηλεκτρικής ενέργειας.

\subsection{Προημερήσια Αγορά Ηλεκτρικής Ενέργειας}
\label{subsec:day-ahead-market}

Η Προημερήσια Αγορά (Day-Ahead Market - DAM) αποτελεί κεντρικό πυλώνα της χονδρικής αγοράς ηλεκτρικής ενέργειας στην Ευρώπη. Σε αυτή την αγορά, οι συμμετέχοντες υποβάλλουν προσφορές για αγορά και πώληση ηλεκτρικής ενέργειας για κάθε ώρα της επόμενης ημέρας.

Η αγορά λειτουργεί συνήθως με \textbf{μηχανισμό δημοπρασίας μοναδικής τιμής} (uniform pricing auction), όπου όλες οι αποδεκτές προσφορές εκκαθαρίζονται στην οριακή τιμή συστήματος (System Marginal Price - SMP). Όσον αφορά το χρονοδιάγραμμα, οι προσφορές υποβάλλονται μέχρι μια συγκεκριμένη ώρα (π.χ. 12:00 CET) και τα αποτελέσματα ανακοινώνονται λίγο αργότερα, καθορίζοντας το πρόγραμμα παραγωγής και κατανάλωσης για την επόμενη ημέρα. Τα προϊόντα συνήθως αφορούν ωριαίες ή ημίωρες ποσότητες ηλεκτρικής ενέργειας, ενώ σε ορισμένες αγορές υπάρχουν και πιο σύνθετα προϊόντα όπως μπλοκ προσφορές και συνδεδεμένες προσφορές.

Στην Ευρώπη, οι περισσότερες εθνικές προημερήσιες αγορές είναι συζευγμένες μέσω του μηχανισμού \textbf{Price Coupling of Regions} (PCR), που βελτιστοποιεί τη χρήση των διασυνδέσεων και προωθεί τη σύγκλιση των τιμών. Για την επίλυση της συζευγμένης προημερήσιας αγοράς χρησιμοποιείται ο \textbf{αλγόριθμος EUPHEMIA} (EU Pan-European Hybrid Electricity Market Integration Algorithm), ο οποίος λαμβάνει υπόψη τις προσφορές αγοράς και πώλησης, τους περιορισμούς μεταφοράς και τα διαθέσιμα περιθώρια των διασυνδέσεων~\cite{euphemia-public-description}.

Η Προημερήσια Αγορά παρέχει ένα αξιόπιστο σημείο αναφοράς για τις τιμές ηλεκτρικής ενέργειας και μία βάση για τη διαμόρφωση του προγράμματος λειτουργίας του συστήματος. Επίσης αποτελεί έναν μηχανισμό για την αποτελεσματική κατανομή των πόρων και των διασυνοριακών ικανοτήτων μεταφοράς, καθώς και ένα σημαντικό εργαλείο για τη βελτιστοποίηση του κόστους παραγωγής σε πανευρωπαϊκό επίπεδο.

Οι τιμές που προκύπτουν από την Προημερήσια Αγορά αντανακλούν το οριακό κόστος παραγωγής ηλεκτρικής ενέργειας και επηρεάζονται από παράγοντες όπως η ζήτηση, η διαθεσιμότητα των μονάδων παραγωγής, η παραγωγή από ΑΠΕ και οι περιορισμοί του δικτύου.

\subsection{Ενδοημερήσια Αγορά Ηλεκτρικής Ενέργειας}
\label{subsec:intraday-market}

Η Ενδοημερήσια Αγορά (Intraday Market - IDM) επιτρέπει στους συμμετέχοντες να προσαρμόσουν τις θέσεις τους μετά το κλείσιμο της Προημερήσιας Αγοράς και πριν την πραγματική ώρα παράδοσης της ηλεκτρικής ενέργειας. Παρέχει τη δυνατότητα διαπραγμάτευσης ηλεκτρικής ενέργειας σε χρονικό ορίζοντα ``σχεδόν πραγματικού χρόνου''.

Σε αντίθεση με την Προημερήσια Αγορά που λειτουργεί με δημοπρασία, η Ενδοημερήσια Αγορά συνήθως λειτουργεί με μηχανισμό \textbf{συνεχούς διαπραγμάτευσης} (continuous trading), όπου οι συναλλαγές πραγματοποιούνται συνεχώς μέχρι το κλείσιμο της πύλης (gate closure) που είναι συνήθως 5-60 λεπτά πριν την παράδοση. Η αγορά επιτρέπει τη διαχείριση αβεβαιότητας και την αντιμετώπιση απρόβλεπτων γεγονότων όπως βλάβες μονάδων παραγωγής, μεταβολές στην πρόβλεψη παραγωγής ΑΠΕ ή απροσδόκητες διακυμάνσεις στη ζήτηση. Το Ευρωπαϊκό μοντέλο \textbf{XBID} (Cross-Border Intraday) παρέχει μια ενιαία πλατφόρμα για διασυνοριακές ενδοημερήσιες συναλλαγές σε όλη την Ευρώπη, ενισχύοντας τη ρευστότητα και την αποτελεσματικότητα της αγοράς.

Η Ενδοημερήσια Αγορά προσφέρει βελτιωμένη αποτελεσματικότητα του συστήματος, καθώς οι συμμετέχοντες μπορούν να προσαρμόσουν τις θέσεις τους με βάση πιο πρόσφατες πληροφορίες. Αποτελεί ένα σημαντικό εργαλείο για την ενσωμάτωση μεταβλητών ΑΠΕ, επιτρέποντας την αντιμετώπιση σφαλμάτων πρόβλεψης και βοηθάει στη μείωση της ανάγκης για ενεργοποίηση πόρων εξισορρόπησης, καθώς αποκλίσεις μπορούν να διορθωθούν μέσω της αγοράς αυτής. Η σημασία της Ενδοημερήσιας Αγοράς αυξάνεται συνεχώς, ιδιαίτερα με την αυξανόμενη διείσδυση των ΑΠΕ και τις προκλήσεις που αυτή δημιουργεί στην πρόβλεψη της παραγωγής.

\subsection{Αγορά Εξισορρόπησης}
\label{subsec:balancing-market}

Η Αγορά Εξισορρόπησης (Balancing Market) αποτελεί το τελευταίο χρονικά στάδιο των αγορών ηλεκτρικής ενέργειας και λειτουργεί σε πραγματικό ή κοντά σε πραγματικό χρόνο. Ο βασικός της στόχος είναι να διασφαλίσει την ισορροπία μεταξύ παραγωγής και κατανάλωσης ηλεκτρικής ενέργειας ανά πάσα στιγμή, διατηρώντας τη συχνότητα του συστήματος στα επιτρεπτά όρια.

Η αγορά διαχειρίζεται από τον Διαχειριστή του Συστήματος Μεταφοράς (TSO), ο οποίος είναι υπεύθυνος για την εξισορρόπηση του συστήματος σε πραγματικό χρόνο. Τα προϊόντα εξισορρόπησης περιλαμβάνουν τις Εφεδρείες Διατήρησης Συχνότητας (FCR - Frequency Containment Reserves) με αυτόματη απόκριση εντός δευτερολέπτων, τις Εφεδρείες Αποκατάστασης Συχνότητας (FRR - Frequency Restoration Reserves) που διακρίνονται σε αυτόματες (aFRR) και χειροκίνητες (mFRR), καθώς και τις Εφεδρείες Αντικατάστασης (RR - Replacement Reserves) για την αποκατάσταση των προηγούμενων εφεδρειών~\cite{entsoe-balancing}.

Για την αποτελεσματικότερη λειτουργία της αγοράς, έχουν αναπτυχθεί πανευρωπαϊκές πλατφόρμες εξισορρόπησης όπως η πλατφόρμα PICASSO για την ανταλλαγή aFRR, η πλατφόρμα MARI για την ανταλλαγή mFRR και η TERRE για την ανταλλαγή RR. Η τιμολόγηση ενέργειας εξισορρόπησης μπορεί να βασίζεται είτε στην τιμή προσφοράς (pay-as-bid) είτε στην οριακή τιμή (pay-as-cleared), με την τελευταία να προτιμάται στις περισσότερες ευρωπαϊκές αγορές. Οι συμμετέχοντες που αποκλίνουν από το πρόγραμμά τους χρεώνονται ή πιστώνονται με την τιμή απόκλισης (imbalance price).

Η Αγορά Εξισορρόπησης είναι καθοριστικής σημασίας για τη διασφάλιση της αξιοπιστίας και ασφάλειας του συστήματος ηλεκτρικής ενέργειας, την αποτελεσματική διαχείριση των διακυμάνσεων της παραγωγής από ΑΠΕ και την παροχή οικονομικών κινήτρων για επενδύσεις σε ευέλικτους πόρους. Με την αυξανόμενη διείσδυση των ΑΠΕ, η σημασία και η πολυπλοκότητα της Αγοράς Εξισορρόπησης αυξάνεται, απαιτώντας συνεχή εξέλιξη των μηχανισμών της και στενότερη διασυνοριακή συνεργασία μεταξύ των Διαχειριστών Συστημάτων Μεταφοράς.


% =============================================================================
\section{Σκοπός και Στόχοι της Εργασίας}
\label{sec:objectives}

Η παρούσα διπλωματική εργασία έχει ως κεντρικό στόχο την ανάπτυξη ενός ολοκληρωμένου, ανοιχτού κώδικα εργαλείου για την προσομοίωση της Προημερήσιας Αγοράς ηλεκτρικής ενέργειας, με έμφαση στον αλγόριθμο EUPHEMIA που χρησιμοποιείται για την εκκαθάριση της συζευγμένης ευρωπαϊκής αγοράς.

\subsection{Κίνητρο και Αναγκαιότητα}
\label{subsec:motivation}

Η κατανόηση των μηχανισμών λειτουργίας των αγορών ηλεκτρικής ενέργειας αποτελεί κρίσιμο παράγοντα για όλους τους εμπλεκόμενους φορείς, από τους παραγωγούς και τους προμηθευτές έως τους ρυθμιστές και τους ερευνητές. Ωστόσο, η πρόσβαση σε εργαλεία προσομοίωσης που αναπαράγουν με ακρίβεια τη λειτουργία της αγοράς είναι περιορισμένη, καθώς τα υπάρχοντα συστήματα είναι συνήθως κλειστά και ιδιόκτητα.

Η ενεργειακή μετάβαση και η αυξανόμενη διείσδυση των ΑΠΕ δημιουργούν νέες προκλήσεις για τη λειτουργία των αγορών ηλεκτρικής ενέργειας. Η μεταβλητότητα της παραγωγής από ανανεώσιμες πηγές, οι περιορισμοί του δικτύου μεταφοράς και η ανάγκη για διασυνοριακή συνεργασία καθιστούν τη μοντελοποίηση της αγοράς ένα σύνθετο πρόβλημα βελτιστοποίησης. Η δυνατότητα προσομοίωσης διαφορετικών σεναρίων και η κατανόηση των παραγόντων που επηρεάζουν τις τιμές της ηλεκτρικής ενέργειας είναι απαραίτητη για τη λήψη αποφάσεων τόσο σε επιχειρησιακό όσο και σε πολιτικό επίπεδο.

\subsection{Επιστημονική Συνεισφορά}
\label{subsec:contribution}

Η συνεισφορά της παρούσας εργασίας εντοπίζεται σε τρεις κύριους άξονες.

Πρώτον, στην \textbf{ανάπτυξη ολοκληρωμένης υποδομής δεδομένων}. Σχεδιάστηκε και υλοποιήθηκε ένα αυτοματοποιημένο σύστημα συλλογής δεδομένων από την πλατφόρμα διαφάνειας του ENTSO-E, το οποίο αποθηκεύει και διαχειρίζεται δεδομένα για μονάδες παραγωγής, φορτία, προβλέψεις ΑΠΕ, διασυνοριακές μεταφορές και μη διαθεσιμότητες μονάδων. Η υποδομή αυτή επιτρέπει την εκτέλεση προσομοιώσεων με πραγματικά δεδομένα για οποιαδήποτε ζώνη προσφοράς (bidding zone) της Ευρώπης.

Δεύτερον, στην \textbf{υλοποίηση αλγορίθμων βελτιστοποίησης αγοράς}. Αναπτύχθηκε μια σειρά από αλγορίθμους που καλύπτουν το πλήρες φάσμα της λειτουργίας της Προημερήσιας Αγοράς, συμπεριλαμβανομένου του προβλήματος Unit Commitment για τον προσδιορισμό του βέλτιστου προγράμματος λειτουργίας των μονάδων παραγωγής, της μηχανής στρατηγικών υποβολής προσφορών, του βιβλίου εντολών οριακού κόστους (merit-order) που κατασκευάζει ανταγωνιστικό αντιπαράδειγμα της αγοράς από τα θεμελιώδη μεγέθη της (τιμές καυσίμων και CO$_2$, αξία νερού, καθαρές εισαγωγές), του αλγορίθμου MPCC (Mathematical Programming with Complementarity Constraints) για την εκκαθάριση της αγοράς και της μοντελοποίησης δικτύου με περιορισμούς ATC (Available Transfer Capacity) για την πολυζωνική εκκαθάριση της αγοράς.

Τρίτον, στη \textbf{δημιουργία ανοιχτού κώδικα εργαλείου με ποσοτικά επικυρωμένη ακρίβεια}. Ολόκληρο το σύστημα αναπτύχθηκε σε γλώσσα Julia με χρήση του πλαισίου JuMP για τη μοντελοποίηση προβλημάτων βελτιστοποίησης, και ο κώδικας διατίθεται ελεύθερα με άδεια ανοιχτού κώδικα. Το σύστημα αναπαράγει τις τιμές της ελληνικής ζώνης με συντελεστή συσχέτισης 0{,}84 σε ανάλυση δεκαπενταλέπτου επί συνεχούς τετραμήνου (Κεφάλαιο~\ref{ch:results})· εξ όσων είμαστε σε θέση να γνωρίζουμε, πρόκειται για το πρώτο ανοιχτό εργαλείο που μοντελοποιεί την αγορά αυτή με τεκμηριωμένη ακρίβεια αυτού του επιπέδου.

\subsection{Ερευνητικά Ερωτήματα}
\label{subsec:research-questions}

Η εργασία επιχειρεί να απαντήσει στα ακόλουθα ερευνητικά ερωτήματα. Πρώτον, είναι δυνατή η αναπαραγωγή των τιμών της Προημερήσιας Αγοράς με ικανοποιητική ακρίβεια χρησιμοποιώντας δημοσίως διαθέσιμα δεδομένα και αλγορίθμους ανοιχτού κώδικα; Δεύτερον, ποια είναι η επίδραση των διαφορετικών στρατηγικών υποβολής προσφορών στις τελικές τιμές της αγοράς; Τρίτον, πώς επηρεάζουν οι περιορισμοί διασυνοριακής μεταφοράς τη διαμόρφωση των τιμών σε γειτονικές ζώνες προσφοράς;


% =============================================================================
\section{Δομή της Εργασίας}
\label{sec:thesis-structure}

Η παρούσα διπλωματική εργασία οργανώνεται σε έξι κεφάλαια, τα οποία καλύπτουν το θεωρητικό υπόβαθρο, την αρχιτεκτονική του συστήματος, την υλοποίηση, τα αποτελέσματα και τα συμπεράσματα.

Στο \textbf{Κεφάλαιο~\ref{ch:introduction}} (Εισαγωγή), το παρόν κεφάλαιο, παρουσιάζεται το πλαίσιο της εργασίας με αναφορά στην κλιματική αλλαγή, την ενεργειακή μετάβαση και τη δομή των ευρωπαϊκών αγορών ηλεκτρικής ενέργειας. Επίσης, καθορίζονται οι στόχοι και η συνεισφορά της εργασίας.

Στο \textbf{Κεφάλαιο~\ref{ch:theoretical-background}} (Θεωρητικό Υπόβαθρο) αναλύονται οι θεωρητικές βάσεις της εργασίας, συμπεριλαμβανομένων των μεθόδων μαθηματικού προγραμματισμού, του προβλήματος Unit Commitment, του αλγορίθμου EUPHEMIA και των στρατηγικών υποβολής προσφορών.

Στο \textbf{Κεφάλαιο~\ref{ch:system-architecture}} (Αρχιτεκτονική Συστήματος και Δεδομένα) περιγράφεται η αρχιτεκτονική του συστήματος που αναπτύχθηκε, με έμφαση στο pipeline συλλογής δεδομένων από το ENTSO-E, το σχήμα της βάσης δεδομένων και το οικοσύστημα Julia/JuMP.

Στο \textbf{Κεφάλαιο~\ref{ch:implementation}} (Υλοποίηση) παρουσιάζεται αναλυτικά η υλοποίηση των αλγορίθμων, συμπεριλαμβανομένου του solver για το πρόβλημα Unit Commitment, της μηχανής στρατηγικών υποβολής προσφορών, του βιβλίου εντολών οριακού κόστους, του αλγορίθμου MPCC για την εκκαθάριση της αγοράς και της μοντελοποίησης των περιορισμών δικτύου.

Στο \textbf{Κεφάλαιο~\ref{ch:results}} (Αποτελέσματα και Αξιολόγηση) παρουσιάζονται τα αποτελέσματα της εφαρμογής του συστήματος σε πραγματικά δεδομένα, η σύγκριση με τις πραγματικές τιμές της αγοράς και η ανάλυση της απόδοσης του συστήματος.

Στο \textbf{Κεφάλαιο~\ref{ch:conclusions}} (Συμπεράσματα) συνοψίζονται τα κύρια ευρήματα της εργασίας, αναφέρονται οι περιορισμοί και προτείνονται κατευθύνσεις για μελλοντική έρευνα.

Τέλος, στο \textbf{Παράρτημα} παρατίθενται συμπληρωματικά στοιχεία και σύνδεσμοι προς τον κώδικα της εφαρμογής.

% =============================================================================

\chapter{Εισαγωγή}
\label{ch:introduction}

\section{Κλιματική Αλλαγή και μέτρα για την αντιμετώπισή της}
\label{sec:climate-change}

Η Κλιματική Αλλαγή αναδεικνύεται ως η κορυφαία πρόκληση της εποχής μας, συνιστώντας μια υπαρξιακή απειλή για το φυσικό περιβάλλον, την συνοχή των ανθρώπινων κοινωνιών και την παγκόσμια οικονομική σταθερότητα. Η επιστημονική κοινότητα συμφωνεί ότι η ραγδαία αύξηση της μέσης θερμοκρασίας του πλανήτη και η εντατικοποίηση των ακραίων καιρικών φαινομένων οφείλονται, κατά κύριο λόγο, στις ανθρωπογενείς εκπομπές αερίων του θερμοκηπίου~\cite{lynas2021}.

Μέρος αυτών των εκπομπών προέρχεται από την χρήση ορυκτών καυσίμων όπως ο άνθρακας, το πετρέλαιο, ο λιγνίτης και το φυσικό αέριο, για την παραγωγή ηλεκτρικής ενέργειας. Η πρωτοκαθεδρία των ορυκτών καυσίμων παραμένει αδιαμφισβήτητη, αντιπροσωπεύοντας συνολικά σχεδόν το 60\% της παγκόσμιας παραγωγής ηλεκτρικής ενέργειας για το έτος 2024. Η εν λόγω αναλογία υπογραμμίζει την συνεχιζόμενη εξάρτηση των παγκόσμιων ενεργειακών συστημάτων από συμβατικές πηγές που επιτείνουν την κλιματική αλλαγή και τις επιπτώσεις αυτής. Αξίζει όμως να σημειωθεί ότι το ενεργειακό μίγμα βελτιώνεται σταδιακά, καθότι οι Ανανεώσιμες Πηγές Ενέργειας (ΑΠΕ) ανήλθαν συνολικά στο ένα τρίτο (33\%) της ηλεκτροπαραγωγής, όπου μαζί με την πυρηνική ενέργεια (9\%), άγγιξαν για πρώτη φορά το 40\% της παγκόσμιας ηλεκτροπαραγωγής ενέργειας~\cite{iea2024}.

\subsection{Παγκόσμιο επίπεδο}
\label{subsec:global-level}

Η Σύμβαση-Πλαίσιο των Ηνωμένων Εθνών για την Κλιματική Αλλαγή (UNFCCC), που συμφωνήθηκε το 1992 στη Διάσκεψη του Ρίο, αποτελεί τη βασική διεθνή συνθήκη για την αντιμετώπιση της κλιματικής αλλαγής. Στόχος της είναι η πρόληψη της επικίνδυνης ανθρωπογενούς παρέμβασης στο παγκόσμιο κλιματικό σύστημα. Από αυτήν, προέκυψαν συμφωνίες όπως το Πρωτόκολλο του Κιότο (1997), το οποίο εισήγαγε νομικά δεσμευτικούς στόχους μείωσης εκπομπών για τις ανεπτυγμένες χώρες και καθιέρωσε μηχανισμούς όπως η εμπορία δικαιωμάτων εκπομπών (Emissions Trading), ο μηχανισμός ``καθαρής'' ανάπτυξης (Clean Development Mechanism) και ο μηχανισμός κοινής εφαρμογής (Joint Implementation)~\cite{ypen-kyoto}. Το Πρωτόκολλο του Κιότο ήταν πρωτοπόρο για την εποχή, αλλά η αποτελεσματικότητά του περιορίστηκε από τη μη συμμετοχή ή αποχώρηση μεγάλων εκπομπέων όπως οι ΗΠΑ, η Ιαπωνία και ο Καναδάς~\cite{bassetti2022}.

Το 2015, η Συμφωνία του Παρισιού υιοθετήθηκε ως η νέα παγκόσμια συμφωνία για το κλίμα. Ο πρωταρχικός της στόχος είναι να συγκρατήσει ``την αύξηση της μέσης παγκόσμιας θερμοκρασίας πολύ κάτω από τους 2°C σε σχέση με τα προβιομηχανικά επίπεδα'' και να καταβάλει προσπάθειες ``να περιορίσει την αύξηση της θερμοκρασίας στον 1,5°C σε σχέση με τα προβιομηχανικά επίπεδα''~\cite{unfccc-paris}.

\subsection{Ευρωπαϊκό επίπεδο}
\label{subsec:european-level}

Πέραν των ανωτέρω συμφωνιών, η αντιμετώπιση της κλιματικής αλλαγής σε ευρωπαϊκό επίπεδο και πιο συγκεκριμένα όσον αφορά την Ευρωπαϊκή Ένωση, εστιάζει γύρω από τον στόχο της κλιματικής ουδετερότητας έως το 2050 και τις ενδιάμεσες δεσμεύσεις μείωσης των καθαρών εκπομπών τουλάχιστον 55\% έως το 2030 και τουλάχιστον 90\% έως το 2040, σε σχέση με τα επίπεδα του 1990. Η πρωτοβουλία-πλαίσιο που συγκεντρώνει αυτά τα πολιτικά μέτρα είναι το Ευρωπαϊκό Σχέδιο για το Περιβάλλον (European Green Deal), το οποίο προβλέπει αναθεώρηση της υπάρχουσας νομοθεσίας και εισαγωγή νέων πολιτικών σε τομείς όπως η ενεργειακή μετάβαση, η κυκλική οικονομία, η ανακαίνιση κτιρίων, η βιοποικιλότητα και η αγροτική πολιτική~\cite{ec-green-deal}.

Το πακέτο πολιτικών του Green Deal συνδέει ρητά την ενεργειακή αποδέσμευση από τα ορυκτά καύσιμα με μέτρα οικονομικής στήριξης, επενδύσεις σε τεχνολογικές καινοτομίες και κοινωνικά μέτρα για την αντιμετώπιση των επιπτώσεων στη συνοχή των περιοχών και των εργαζομένων. Η νομοθετική δέσμευση ενισχύεται από τον Ευρωπαϊκό Νόμο για το Κλίμα, ο οποίος καθιστά τον στόχο της ουδετερότητας νομικά δεσμευτικό και απαιτεί σαφείς ενδιάμεσους στόχους και μηχανισμούς παρακολούθησης για τα κράτη-μέλη~\cite{eu-climate-law}.

Ως αποτέλεσμα, η ΕΕ με το Green Deal συνδυάζει ρυθμιστικά εργαλεία (όπως αναθεωρήσεις αγοράς δικαιωμάτων εκπομπών, πρότυπα ενεργειακής απόδοσης και κανονισμούς για τη μεταφορά και τη βιομηχανία) με χρηματοδοτικά μέσα (ταμεία δίκαιης μετάβασης, προγράμματα επενδύσεων) και πολιτικές υποστήριξης των ΑΠΕ και της έρευνας. Η ευρωπαϊκή προσέγγιση προωθεί τόσο τη μείωση εκπομπών όσο και την ενίσχυση της κλιματικής ανθεκτικότητας των υποδομών και των κοινοτήτων, με έμφαση στη συνοχή των πολιτικών σε όλα τα επίπεδα διακυβέρνησης~\cite{sikora2020}.

\subsection{Εθνικό επίπεδο}
\label{subsec:national-level}

Σε εθνικό επίπεδο, η Ελλάδα, ως μέλος της ΕΕ, έχει υιοθετήσει ένα πλαίσιο πολιτικών και στρατηγικών που ευθυγραμμίζονται πλήρως με τους στόχους της Ευρωπαϊκής Ένωσης για την κλιματική ουδετερότητα έως το 2050 και τη μείωση των εκπομπών έως το 2030, ενσωματώνοντας μέτρα για την ενεργειακή μετάβαση, την απολιγνιτοποίηση, την επιτάχυνση των ΑΠΕ και την αναβάθμιση της ενεργειακής απόδοσης των κτιρίων.

Το παραπάνω πλαίσιο είναι το αναθεωρημένο Εθνικό Σχέδιο για την Ενέργεια και το Κλίμα (ΕΣΕΚ) το οποίο κυρώθηκε το 2024. Σε αυτό το έγγραφο αναγνωρίζονται οι επιπτώσεις που η κλιματική αλλαγή θα επιφέρει ειδικά στην Ελλάδα, όπως η ``μειωμένη διαθεσιμότητα υδάτων'' η οποία ``θέτει σε κίνδυνο τη λειτουργία θερμοηλεκτρικών μονάδων. Τα δίκτυα μεταφοράς και διανομής ηλεκτρικής ενέργειας, οι σταθμοί υψηλής τάσης, και λοιπά ενεργειακά δίκτυα και εγκαταστάσεις είναι τρωτά σε ακραία καιρικά φαινόμενα, ενώ οι παράκτιες ενεργειακές υποδομές, απειλούνται από την άνοδο της στάθμης της θάλασσας. Η αύξηση της μέσης θερμοκρασίας αναμένεται να μειώσει τις ενεργειακές ανάγκες για θέρμανση κατά τη χειμερινή περίοδο και να αυξήσει τις ανάγκες για ψύξη κατά τη θερινή περίοδο. Οι μεταβολές στη ζήτηση θα οδηγήσουν σε μεγαλύτερη διακύμανση φορτίων δοκιμάζοντας το σύστημα ηλεκτρικής ενέργειας. Επίσης, η κλιματική αλλαγή αναμένεται να επηρεάσει την απόδοση των συστημάτων ΑΠΕ''~\cite{esek2024}.

\subsection{Τεχνητή νοημοσύνη και ενεργειακή κατανάλωση}
\label{subsec:ai-energy}

Η ραγδαία επέκταση της Τεχνητής Νοημοσύνης (AI) και του υπολογιστικού νέφους (cloud computing) θέτει μία αυξανόμενη πρόκληση για την κλιματική αλλαγή, καθώς οδηγεί σε πρωτοφανή αύξηση της κατανάλωσης ενέργειας των κέντρων δεδομένων (data centers). Τα κέντρα δεδομένων, τα οποία ήδη καταναλώνουν περίπου το 1\% της παγκόσμιας ηλεκτρικής ενέργειας, αναμένεται να διπλασιάσουν αυτόν τον αριθμό έως το 2030, κυρίως λόγω των εξαιρετικά απαιτητικών εργασιών της τεχνητής νοημοσύνης, όπως η εκπαίδευση μεγάλων γλωσσικών μοντέλων (Large Language Models - LLMs). 

Η εκπαίδευση και η χρήση αυτών των μοντέλων απαιτεί την εντατική χρήση ενεργοβόρου υλικού όπως GPUs και TPUs, καθώς και μεγάλη ανάγκη για συστήματα ψύξης, τα οποία ευθύνονται για σχεδόν το 40\% της συνολικής ενεργειακής κατανάλωσης. Αυτό σημαίνει ότι η περιβαλλοντική επίπτωση αυξάνεται σημαντικά, εκτός εάν υπάρξει μια σημαντική στροφή σε ανανεώσιμες πηγές ενέργειας και προηγμένες τεχνικές ψύξης. Κατά συνέπεια, η έγκαιρη αντιμετώπιση της ενεργειακής ζήτησης της AI είναι κρίσιμη για τη διασφάλιση της βιωσιμότητας της ψηφιακής υποδομής και τον μετριασμό των αρνητικών επιπτώσεων στο περιβάλλον~\cite{chauhan2024}.


% =============================================================================
\section{Συστήματα Ηλεκτρικής Ενέργειας}
\label{sec:power-systems}

Τα συστήματα ηλεκτρικής ενέργειας αποτελούν το σύνολο των εγκαταστάσεων και μέσων για την παροχή ηλεκτρικής ενέργειας σε εξυπηρετούμενες περιοχές κατανάλωσης. Για την εύρυθμη λειτουργία ενός συστήματος ηλεκτρικής ενέργειας είναι απαραίτητη η παροχή ηλεκτρικής ενέργειας οπουδήποτε υπάρχει ζήτηση με το ελάχιστο δυνατό κόστος και τις ελάχιστες οικολογικές επιπτώσεις, εξασφαλίζοντας σταθερή συχνότητα, σταθερή τάση και υψηλή αξιοπιστία τροφοδότησης~\cite{vournas-kontaxis}.

Η τροφοδότηση των καταναλωτών με ηλεκτρική ενέργεια προϋποθέτει τρεις ξεχωριστές λειτουργίες του συστήματος ηλεκτρικής ενέργειας: την παραγωγή, τη μεταφορά, και τη διανομή. Η παραγόμενη ηλεκτρική ενέργεια πρέπει ιδανικά να απορροφάται από την κατανάλωση σε ίδιο χρόνο, καθώς η αποθήκευσή της παραμένει δύσκολη και τις περισσότερες φορές οικονομικά ασύμφορη, παρά την πρόοδο στην τεχνολογία αποθήκευσης ηλεκτρικής ενέργειας τα τελευταία χρόνια~\cite{hiesl2020}.

% TODO: Εισαγωγή σχήματος συστήματος ηλεκτρικής ενέργειας
\begin{figure}[htbp]
\centering
% \includegraphics[width=0.9\textwidth]{figures/power_system_overview.pdf}
\fbox{\parbox{0.8\textwidth}{\centering\vspace{2cm}[Σχήμα: Γενική δομή συστήματος ηλεκτρικής ενέργειας]\vspace{2cm}}}
\caption{Γενική δομή συστήματος ηλεκτρικής ενέργειας με τις τρεις βασικές λειτουργίες: παραγωγή, μεταφορά και διανομή.}
\label{fig:power-system-overview}
\end{figure}

\subsection{Παραγωγή Ηλεκτρικής Ενέργειας}
\label{subsec:generation}

Η παραγωγή ηλεκτρικής ενέργειας αφορά τη μετατροπή διαφόρων μορφών πρωτογενούς ενέργειας σε ηλεκτρική. Τα εργοστάσια παραγωγής διακρίνονται με βάση την τεχνολογία και την πηγή ενέργειας που χρησιμοποιούν.

Οι \textbf{θερμοηλεκτρικοί σταθμοί παραγωγής} αποτελούν τη συμβατική μορφή παραγωγής ηλεκτρικής ενέργειας. Σε αυτούς, η θερμική ενέργεια που προέρχεται από την καύση ορυκτών καυσίμων (λιγνίτης, φυσικό αέριο, πετρέλαιο) ή από πυρηνικές αντιδράσεις μετατρέπεται σε μηχανική μέσω ατμοστροβίλων ή αεριοστροβίλων, και στη συνέχεια σε ηλεκτρική μέσω γεννητριών. Οι μονάδες αυτές χαρακτηρίζονται από υψηλή διαθεσιμότητα και ελεγχόμενη παραγωγή, αλλά συνεπάγονται σημαντικές εκπομπές αερίων του θερμοκηπίου.

Οι \textbf{υδροηλεκτρικοί σταθμοί} αξιοποιούν τη δυναμική ενέργεια του νερού για την κίνηση υδροστροβίλων. Διακρίνονται σε συμβατικούς (με φράγμα και ταμιευτήρα) και αντλησιοταμιευτικούς, οι οποίοι λειτουργούν και ως συστήματα αποθήκευσης ενέργειας.

Οι \textbf{Ανανεώσιμες Πηγές Ενέργειας (ΑΠΕ)} περιλαμβάνουν τεχνολογίες όπως τα αιολικά πάρκα, τα φωτοβολταϊκά συστήματα, τη γεωθερμία και τη βιομάζα. Οι ΑΠΕ χαρακτηρίζονται από μηδενικές ή ελάχιστες εκπομπές κατά τη λειτουργία τους, αλλά η παραγωγή τους εξαρτάται από τις καιρικές συνθήκες, καθιστώντας τη μεταβλητή και εν μέρει απρόβλεπτη.

\subsection{Μεταφορά Ηλεκτρικής Ενέργειας}
\label{subsec:transmission}

Η μεταφορά ηλεκτρικής ενέργειας αφορά τη μεταφορά μεγάλων ποσοτήτων ηλεκτρικής ενέργειας από τους σταθμούς παραγωγής προς τα κέντρα κατανάλωσης μέσω του δικτύου υψηλής τάσης. Η υψηλή τάση (150kV, 400kV) χρησιμοποιείται για τη μείωση των απωλειών μεταφοράς, καθώς για δεδομένη ισχύ, η αύξηση της τάσης οδηγεί σε μείωση του ρεύματος και συνεπώς των απωλειών λόγω θέρμανσης των αγωγών.

Αρμόδιος για την αξιόπιστη μεταφορά ηλεκτρικής ενέργειας στην ελληνική επικράτεια είναι ο Ανεξάρτητος Διαχειριστής Μεταφοράς Ηλεκτρικής Ενέργειας (ΑΔΜΗΕ). Ο ΑΔΜΗΕ διαχειρίζεται το Διασυνδεδεμένο Σύστημα Μεταφοράς, το οποίο περιλαμβάνει βασικά στοιχεία όπως εναέριες γραμμές μεταφοράς 400kV, 150kV και 66kV, υπόγειες και υποβρύχιες καλωδιακές γραμμές 150kV και 400kV (AC και DC), υποσταθμούς 150/20kV, καθώς και Κέντρα Υπερυψηλής Τάσης (ΚΥΤ) 400/150kV.

Το Ελληνικό Σύστημα Μεταφοράς λειτουργεί παράλληλα με το διασυνδεδεμένο Ευρωπαϊκό Σύστημα, υπό τον ευρύτερο συντονισμό του Ευρωπαϊκού Δικτύου Διαχειριστών Συστημάτων Μεταφοράς Ηλεκτρικής Ενέργειας (European Network of Transmission System Operators for Electricity - ENTSO-E). Η σύνδεση αυτή πραγματοποιείται μέσω διασυνδετικών γραμμών μεταφοράς 400kV, που ενώνουν το ελληνικό σύστημα με εκείνα της Αλβανίας, της Βουλγαρίας, της Βόρειας Μακεδονίας και της Τουρκίας. Επιπλέον, υπάρχει ασύγχρονη διασύνδεση με την Ιταλία μέσω υποβρυχίου συνδέσμου συνεχούς ρεύματος τάσης 400kV~\cite{admie-system}.

\subsection{Διανομή Ηλεκτρικής Ενέργειας}
\label{subsec:distribution}

Από τους υποσταθμούς υποβιβασμού τάσης ξεκινά η διανομή ηλεκτρικής ενέργειας, το σύνολο των διαδικασιών που απαιτούνται για τον αξιόπιστο διαμοιρασμό της στους μετρητές των καταναλωτών, διατηρώντας την τάση στους ζυγούς και τις ροές στις γραμμές του δικτύου εντός επιτρεπτών ορίων.

Αρμόδιος για τη διαχείριση του δικτύου διανομής ηλεκτρικής ενέργειας στην ελληνική επικράτεια είναι ο Διαχειριστής Ελληνικού Δικτύου Διανομής Ηλεκτρικής Ενέργειας (ΔΕΔΔΗΕ). Ο ΔΕΔΔΗΕ είναι επίσης υπεύθυνος για τη διαχείριση των αγορών των Μη Διασυνδεδεμένων Νησιών (ΜΔΝ)~\cite{nomos4001}.

Εντός του δικτύου διανομής ηλεκτρικής ενέργειας έχει ξεκινήσει και η ανάπτυξη των τεχνολογιών έξυπνων δικτύων (smart grids) με χρήση έξυπνων μετρητών (smart meters), συστημάτων εποπτικού ελέγχου (SCADA) και ολοκληρωμένων πληροφοριακών συστημάτων διανομής (DMS), που λειτουργούν καταλυτικά στην εκμετάλλευση των οικιακών εγκαταστάσεων ΑΠΕ, συστημάτων αποθήκευσης καθώς και τις τεχνολογίες Vehicle to Grid (V2G)~\cite{fang2012}.

Μέσω της αμφίδρομης ροής πληροφοριών και ενέργειας, τα έξυπνα δίκτυα επιτρέπουν στους διαχειριστές δικτύου να παρακολουθούν σε πραγματικό χρόνο την κατανάλωση, να εντοπίζουν βλάβες άμεσα και να διαχειρίζονται καλύτερα την κατανομή φορτίου. Επιπλέον, παρέχουν δυνατότητες συμμετοχής των τελικών καταναλωτών στην αγορά μέσω demand response, όπου οι χρήστες προσαρμόζουν τη ζήτησή τους ανάλογα με τα σήματα της αγοράς ή τις τιμές ενέργειας. Η μετάβαση προς τα έξυπνα δίκτυα είναι αναγκαία για την υλοποίηση των στόχων της ενεργειακής μετάβασης, την ενίσχυση της ενεργειακής απόδοσης και την επίτευξη ενός βιώσιμου, ψηφιοποιημένου και αποκεντρωμένου ενεργειακού συστήματος.


% =============================================================================
\section{Αγορές Ηλεκτρικής Ενέργειας}
\label{sec:energy-markets}

Για την εξασφάλιση της ισορροπίας προσφοράς-ζήτησης και της απρόσκοπτης παροχής ενέργειας στους καταναλωτές στο ελάχιστο δυνατό κόστος, τα συστήματα ηλεκτρικής ενέργειας ενθυλακώνονται σε αγορές ηλεκτρικής ενέργειας. Οι αγορές ηλεκτρικής ενέργειας χαρακτηρίζονται από ιδιαιτερότητες σε σχέση με άλλες αγορές εμπορευμάτων ή καταναλωτικών αγαθών, λόγω του υψηλού κόστους αποθήκευσης της ηλεκτρικής ενέργειας~\cite{mo2025}.

% TODO: Εισαγωγή σχήματος δομής αγοράς
\begin{figure}[htbp]
\centering
% \includegraphics[width=0.9\textwidth]{figures/market_structure.pdf}
\fbox{\parbox{0.8\textwidth}{\centering\vspace{2cm}[Σχήμα: Δομή αγορών ηλεκτρικής ενέργειας]\vspace{2cm}}}
\caption{Χρονική αλληλουχία των αγορών ηλεκτρικής ενέργειας, από την προθεσμιακή αγορά έως την αγορά εξισορρόπησης.}
\label{fig:market-structure}
\end{figure}

\subsection{Μονοπωλιακή Αγορά}
\label{subsec:monopoly-market}

Ιστορικά, η ανάπτυξη του τομέα της ηλεκτρικής ενέργειας απαιτούσε τεράστια κεφαλαιουχική δαπάνη, μακροχρόνιο σχεδιασμό και υψηλό επίπεδο τεχνικής εξειδίκευσης, στοιχεία που καθιστούσαν πρακτικά αδύνατη τη συμμετοχή ιδιωτών. Για τον λόγο αυτό, το κράτος ή κρατικά ελεγχόμενες επιχειρήσεις αναλάμβαναν την ευθύνη για την κατασκευή, λειτουργία και διαχείριση του συστήματος παραγωγής, μεταφοράς και διανομής ηλεκτρικής ενέργειας. Το μοντέλο που διαμορφώθηκε ήταν μονοπωλιακό, με χαρακτηριστικό γνώρισμα την κάθετη ολοκλήρωση, δηλαδή τον έλεγχο όλων των σταδίων της αλυσίδας αξίας --- από την παραγωγή έως την προμήθεια στον τελικό καταναλωτή --- από έναν και μόνο φορέα. Στην ελληνική επικράτεια το ρόλο αυτό είχε αναλάβει η Δημόσια Επιχείρηση Ηλεκτρισμού (ΔΕΗ).

Σε αυτό το πλαίσιο απουσίας ανταγωνισμού οι καταναλωτές δεν είχαν τη δυνατότητα επιλογής παρόχου. Οι τιμές της ηλεκτρικής ενέργειας καθορίζονταν διοικητικά από τις αρμόδιες ρυθμιστικές αρχές, με βάση το κόστος και ένα σταθερό περιθώριο κέρδους. Επιπλέον, οι επενδύσεις και η επέκταση του δικτύου σχεδιάζονταν κεντρικά, με στόχο την κάλυψη της ζήτησης και την ασφάλεια εφοδιασμού. Παρά τα πλεονεκτήματα σταθερότητας και ελέγχου, το μονοπωλιακό αυτό μοντέλο συχνά οδηγούσε σε αναποτελεσματικότητα, καθυστερήσεις στην υιοθέτηση καινοτομιών και υψηλότερο κόστος για τον τελικό καταναλωτή~\cite{vlados2021}.

Συγκεκριμένα, η ρύθμιση ενός μονοπωλιακού φορέα, όπως μια κρατική ή ρυθμιζόμενη επιχείρηση ηλεκτρικής ενέργειας, αποτελεί χαρακτηριστικό πρόβλημα εντολέα-εντολοδόχου (principal-agent). Ο ρυθμιστής (εντολέας) επιδιώκει τη μεγιστοποίηση του κοινωνικού πλεονάσματος, ενώ η επιχείρηση (εντολοδόχος) στοχεύει στη μεγιστοποίηση του κέρδους. Ωστόσο, ο ρυθμιστής δεν διαθέτει πλήρη πληροφόρηση για το κόστος και τα πραγματικά δεδομένα λειτουργίας της επιχείρησης, η οποία έχει κίνητρο να τα υπερεκτιμά, προβάλλοντας αυξημένες ανάγκες για ασφάλεια και επενδύσεις. Η ασυμμετρία πληροφόρησης αυτή δημιουργεί προβλήματα όπως ηθικός κίνδυνος και αρνητική επιλογή, οδηγώντας σε αποκλίσεις από την αποδοτικότητα του τέλειου ανταγωνισμού~\cite{grossman-hart1992}.

Για την αντιμετώπιση αυτών των στρεβλώσεων, ο ρυθμιστής μπορεί είτε να μειώσει την ασυμμετρία πληροφόρησης, μέσω ελέγχων ή συγκριτικής αξιολόγησης (benchmarking), είτε να εφαρμόσει στρατηγικές κινήτρων, ώστε να ενισχύσει την αποδοτικότητα της επιχείρησης. Παρά τις προσπάθειες αυτές, η ισορροπία που επιτυγχάνεται μεταξύ ρυθμιστή και επιχείρησης δεν φτάνει ποτέ το επίπεδο αποδοτικότητας που προκύπτει σε καθεστώς τέλειου ανταγωνισμού, με αποτέλεσμα την ύπαρξη κόστους πρακτόρευσης (agency cost) για την κοινωνία~\cite{panda-leepsa2017}.

\subsection{Άνοιγμα της αγοράς}
\label{subsec:market-liberalization}

Για την αντιμετώπιση των προβλημάτων αποδοτικότητας που προκύπτουν λόγω του κόστους πρακτόρευσης στο μονοπωλιακό μοντέλο αγοράς, από τη δεκαετία του 1990 η Ευρωπαϊκή Ένωση ξεκίνησε μια διαδικασία απελευθέρωσης των αγορών ηλεκτρικής ενέργειας, με στόχο τη δημιουργία μιας ενιαίας, ανταγωνιστικής αγοράς. Τα κύρια στάδια αυτής της διαδικασίας περιλαμβάνουν το Πρώτο Ενεργειακό Πακέτο (1996), όπου η Οδηγία 96/92/ΕΚ έθεσε τα θεμέλια για το άνοιγμα των αγορών, εισάγοντας την έννοια του διαχωρισμού των δραστηριοτήτων και την πρόσβαση τρίτων στο δίκτυο. Το Δεύτερο Ενεργειακό Πακέτο (2003) με τις Οδηγίες 2003/54/ΕΚ και 2003/55/ΕΚ επιτάχυνε το άνοιγμα της αγοράς, καθιστώντας υποχρεωτικό το νομικό διαχωρισμό των δραστηριοτήτων μεταφοράς και διανομής από την παραγωγή και προμήθεια. Το Τρίτο Ενεργειακό Πακέτο (2009) εισήγαγε αυστηρότερους κανόνες για το διαχωρισμό της ιδιοκτησίας, ενίσχυσε τις αρμοδιότητες των εθνικών ρυθμιστικών αρχών και ίδρυσε τον Οργανισμό Συνεργασίας των Ρυθμιστικών Αρχών Ενέργειας (ACER). Τέλος, το Τέταρτο Ενεργειακό Πακέτο (2019), γνωστό και ως ``Καθαρή Ενέργεια για όλους τους Ευρωπαίους'', εστιάζει στην ενσωμάτωση των ανανεώσιμων πηγών ενέργειας, την ενεργειακή αποδοτικότητα και νέους κανόνες για την αγορά ηλεκτρικής ενέργειας~\cite{eikeland-duffield2011, eu-energy-packages}.

Με βάση τα παραπάνω, το άνοιγμα της αγοράς οδήγησε στον ιδιοκτησιακό διαχωρισμό των δραστηριοτήτων, τη δημιουργία ανταγωνιστικών αγορών χονδρικής και λιανικής, την ελευθερία επιλογής προμηθευτή για τους καταναλωτές και την ανάπτυξη νέων επιχειρηματικών μοντέλων και υπηρεσιών.

\subsection{Φορείς της αγοράς}
\label{subsec:market-participants}

Στο σύγχρονο μοντέλο των απελευθερωμένων αγορών ηλεκτρικής ενέργειας, διάφοροι φορείς διαδραματίζουν καθοριστικό ρόλο. Οι \textbf{Παραγωγοί} είναι εταιρείες που διαθέτουν και λειτουργούν μονάδες παραγωγής ηλεκτρικής ενέργειας, περιλαμβάνοντας τόσο συμβατικούς παραγωγούς (θερμικές μονάδες) όσο και παραγωγούς ΑΠΕ. Ο \textbf{Διαχειριστής Συστήματος Μεταφοράς} (TSO - Transmission System Operator) είναι υπεύθυνος για τη λειτουργία, συντήρηση και ανάπτυξη του συστήματος μεταφοράς υψηλής τάσης, καθώς και για την εξισορρόπηση του συστήματος σε πραγματικό χρόνο. Για την Ελλάδα, αυτός ο φορέας είναι ο ΑΔΜΗΕ. Ο \textbf{Διαχειριστής Δικτύου Διανομής} (DSO - Distribution System Operator) είναι υπεύθυνος για τη λειτουργία, συντήρηση και ανάπτυξη του δικτύου διανομής μέσης και χαμηλής τάσης. Για την Ελλάδα, αυτός ο φορέας είναι ο ΔΕΔΔΗΕ.

Οι \textbf{Προμηθευτές} αγοράζουν ηλεκτρική ενέργεια από την αγορά χονδρικής και την πωλούν στους τελικούς καταναλωτές. Οι \textbf{Έμποροι} (Traders) αγοράζουν και πωλούν ηλεκτρική ενέργεια στις χονδρικές αγορές με σκοπό το κέρδος, χωρίς να είναι απαραίτητα παραγωγοί ή προμηθευτές. Τα \textbf{Χρηματιστήρια Ενέργειας} (Power Exchanges) είναι οργανωμένες αγορές όπου διαπραγματεύονται προϊόντα ηλεκτρικής ενέργειας. Παραδείγματα περιλαμβάνουν το EPEX SPOT, το Nord Pool, το GME και για την Ελλάδα το HENEX (EnExGroup). Ο \textbf{Διαχειριστής Αγοράς} είναι υπεύθυνος για την οργάνωση και λειτουργία των διαφόρων τμημάτων της αγοράς ηλεκτρικής ενέργειας. Οι \textbf{Ρυθμιστικές Αρχές} εποπτεύουν τη λειτουργία της αγοράς, θεσπίζουν κανόνες και διασφαλίζουν τον υγιή ανταγωνισμό. Σε ευρωπαϊκό επίπεδο, ο ACER συντονίζει τις εθνικές ρυθμιστικές αρχές. Οι \textbf{Καταναλωτές} περιλαμβάνουν βιομηχανικούς, εμπορικούς και οικιακούς χρήστες. Στις σύγχρονες αγορές, οι καταναλωτές μπορούν επίσης να συμμετέχουν ενεργά στην αγορά μέσω της απόκρισης ζήτησης και της αυτοπαραγωγής.

\subsection{Προθεσμιακή Αγορά - OTC}
\label{subsec:forward-market}

Η Προθεσμιακή Αγορά αποτελεί το πρώτο στάδιο της χονδρικής αγοράς ηλεκτρικής ενέργειας, όπου οι συμμετέχοντες μπορούν να διαπραγματευτούν και να συνάψουν συμβόλαια για μελλοντική παράδοση ηλεκτρικής ενέργειας. Διακρίνεται σε \textbf{Οργανωμένη Προθεσμιακή Αγορά}, η οποία λειτουργεί μέσω χρηματιστηρίων ενέργειας (π.χ. EEX - European Energy Exchange) όπου διαπραγματεύονται τυποποιημένα προθεσμιακά συμβόλαια (futures, forwards, options), και σε \textbf{Εξωχρηματιστηριακή Αγορά} (OTC - Over The Counter), όπου οι συναλλαγές πραγματοποιούνται απευθείας μεταξύ των αντισυμβαλλομένων, με συμβόλαια προσαρμοσμένα στις ανάγκες τους.

Τα βασικά χαρακτηριστικά της Προθεσμιακής Αγοράς περιλαμβάνουν τον χρονικό ορίζοντα (τα συμβόλαια μπορεί να αφορούν παράδοση ηλεκτρικής ενέργειας από μερικές ημέρες έως αρκετά χρόνια στο μέλλον), τη διαχείριση κινδύνου (επιτρέπει στους συμμετέχοντες να αντισταθμίσουν τον κίνδυνο μεταβολής των τιμών), την ευελιξία (ιδιαίτερα στην αγορά OTC, τα συμβόλαια μπορούν να προσαρμοστούν ως προς τον όγκο, τη διάρκεια και τους όρους παράδοσης), τη ρευστότητα (στις οργανωμένες αγορές, τα τυποποιημένα προϊόντα έχουν συνήθως υψηλότερη ρευστότητα) και τον πιστωτικό κίνδυνο (στις συναλλαγές OTC υπάρχει κίνδυνος αθέτησης από τον αντισυμβαλλόμενο).

Η Προθεσμιακή Αγορά διαδραματίζει σημαντικό ρόλο στη διαμόρφωση μακροπρόθεσμων τιμολογιακών σημάτων, συμβάλλει στην οικονομική σταθερότητα των συμμετεχόντων και παρέχει ενδείξεις για μελλοντικές επενδύσεις στον τομέα της ηλεκτρικής ενέργειας.

\subsection{Προημερήσια Αγορά Ηλεκτρικής Ενέργειας}
\label{subsec:day-ahead-market}

Η Προημερήσια Αγορά (Day-Ahead Market - DAM) αποτελεί κεντρικό πυλώνα της χονδρικής αγοράς ηλεκτρικής ενέργειας στην Ευρώπη. Σε αυτή την αγορά, οι συμμετέχοντες υποβάλλουν προσφορές για αγορά και πώληση ηλεκτρικής ενέργειας για κάθε ώρα της επόμενης ημέρας.

Η αγορά λειτουργεί συνήθως με \textbf{μηχανισμό δημοπρασίας μοναδικής τιμής} (uniform pricing auction), όπου όλες οι αποδεκτές προσφορές εκκαθαρίζονται στην οριακή τιμή συστήματος (System Marginal Price - SMP). Όσον αφορά το χρονοδιάγραμμα, οι προσφορές υποβάλλονται μέχρι μια συγκεκριμένη ώρα (π.χ. 12:00 CET) και τα αποτελέσματα ανακοινώνονται λίγο αργότερα, καθορίζοντας το πρόγραμμα παραγωγής και κατανάλωσης για την επόμενη ημέρα. Τα προϊόντα συνήθως αφορούν ωριαίες ή ημίωρες ποσότητες ηλεκτρικής ενέργειας, ενώ σε ορισμένες αγορές υπάρχουν και πιο σύνθετα προϊόντα όπως μπλοκ προσφορές και συνδεδεμένες προσφορές.

Στην Ευρώπη, οι περισσότερες εθνικές προημερήσιες αγορές είναι συζευγμένες μέσω του μηχανισμού \textbf{Price Coupling of Regions} (PCR), που βελτιστοποιεί τη χρήση των διασυνδέσεων και προωθεί τη σύγκλιση των τιμών. Για την επίλυση της συζευγμένης προημερήσιας αγοράς χρησιμοποιείται ο \textbf{αλγόριθμος EUPHEMIA} (EU Pan-European Hybrid Electricity Market Integration Algorithm), ο οποίος λαμβάνει υπόψη τις προσφορές αγοράς και πώλησης, τους περιορισμούς μεταφοράς και τα διαθέσιμα περιθώρια των διασυνδέσεων~\cite{euphemia-public-description}.

Η Προημερήσια Αγορά παρέχει ένα αξιόπιστο σημείο αναφοράς για τις τιμές ηλεκτρικής ενέργειας και μία βάση για τη διαμόρφωση του προγράμματος λειτουργίας του συστήματος. Επίσης αποτελεί έναν μηχανισμό για την αποτελεσματική κατανομή των πόρων και των διασυνοριακών ικανοτήτων μεταφοράς, καθώς και ένα σημαντικό εργαλείο για τη βελτιστοποίηση του κόστους παραγωγής σε πανευρωπαϊκό επίπεδο.

Οι τιμές που προκύπτουν από την Προημερήσια Αγορά αντανακλούν το οριακό κόστος παραγωγής ηλεκτρικής ενέργειας και επηρεάζονται από παράγοντες όπως η ζήτηση, η διαθεσιμότητα των μονάδων παραγωγής, η παραγωγή από ΑΠΕ και οι περιορισμοί του δικτύου.

\subsection{Ενδοημερήσια Αγορά Ηλεκτρικής Ενέργειας}
\label{subsec:intraday-market}

Η Ενδοημερήσια Αγορά (Intraday Market - IDM) επιτρέπει στους συμμετέχοντες να προσαρμόσουν τις θέσεις τους μετά το κλείσιμο της Προημερήσιας Αγοράς και πριν την πραγματική ώρα παράδοσης της ηλεκτρικής ενέργειας. Παρέχει τη δυνατότητα διαπραγμάτευσης ηλεκτρικής ενέργειας σε χρονικό ορίζοντα ``σχεδόν πραγματικού χρόνου''.

Σε αντίθεση με την Προημερήσια Αγορά που λειτουργεί με δημοπρασία, η Ενδοημερήσια Αγορά συνήθως λειτουργεί με μηχανισμό \textbf{συνεχούς διαπραγμάτευσης} (continuous trading), όπου οι συναλλαγές πραγματοποιούνται συνεχώς μέχρι το κλείσιμο της πύλης (gate closure) που είναι συνήθως 5-60 λεπτά πριν την παράδοση. Η αγορά επιτρέπει τη διαχείριση αβεβαιότητας και την αντιμετώπιση απρόβλεπτων γεγονότων όπως βλάβες μονάδων παραγωγής, μεταβολές στην πρόβλεψη παραγωγής ΑΠΕ ή απροσδόκητες διακυμάνσεις στη ζήτηση. Το Ευρωπαϊκό μοντέλο \textbf{XBID} (Cross-Border Intraday) παρέχει μια ενιαία πλατφόρμα για διασυνοριακές ενδοημερήσιες συναλλαγές σε όλη την Ευρώπη, ενισχύοντας τη ρευστότητα και την αποτελεσματικότητα της αγοράς.

Η Ενδοημερήσια Αγορά προσφέρει βελτιωμένη αποτελεσματικότητα του συστήματος, καθώς οι συμμετέχοντες μπορούν να προσαρμόσουν τις θέσεις τους με βάση πιο πρόσφατες πληροφορίες. Αποτελεί ένα σημαντικό εργαλείο για την ενσωμάτωση μεταβλητών ΑΠΕ, επιτρέποντας την αντιμετώπιση σφαλμάτων πρόβλεψης και βοηθάει στη μείωση της ανάγκης για ενεργοποίηση πόρων εξισορρόπησης, καθώς αποκλίσεις μπορούν να διορθωθούν μέσω της αγοράς αυτής. Η σημασία της Ενδοημερήσιας Αγοράς αυξάνεται συνεχώς, ιδιαίτερα με την αυξανόμενη διείσδυση των ΑΠΕ και τις προκλήσεις που αυτή δημιουργεί στην πρόβλεψη της παραγωγής.

\subsection{Αγορά Εξισορρόπησης}
\label{subsec:balancing-market}

Η Αγορά Εξισορρόπησης (Balancing Market) αποτελεί το τελευταίο χρονικά στάδιο των αγορών ηλεκτρικής ενέργειας και λειτουργεί σε πραγματικό ή κοντά σε πραγματικό χρόνο. Ο βασικός της στόχος είναι να διασφαλίσει την ισορροπία μεταξύ παραγωγής και κατανάλωσης ηλεκτρικής ενέργειας ανά πάσα στιγμή, διατηρώντας τη συχνότητα του συστήματος στα επιτρεπτά όρια.

Η αγορά διαχειρίζεται από τον Διαχειριστή του Συστήματος Μεταφοράς (TSO), ο οποίος είναι υπεύθυνος για την εξισορρόπηση του συστήματος σε πραγματικό χρόνο. Τα προϊόντα εξισορρόπησης περιλαμβάνουν τις Εφεδρείες Διατήρησης Συχνότητας (FCR - Frequency Containment Reserves) με αυτόματη απόκριση εντός δευτερολέπτων, τις Εφεδρείες Αποκατάστασης Συχνότητας (FRR - Frequency Restoration Reserves) που διακρίνονται σε αυτόματες (aFRR) και χειροκίνητες (mFRR), καθώς και τις Εφεδρείες Αντικατάστασης (RR - Replacement Reserves) για την αποκατάσταση των προηγούμενων εφεδρειών~\cite{entsoe-balancing}.

Για την αποτελεσματικότερη λειτουργία της αγοράς, έχουν αναπτυχθεί πανευρωπαϊκές πλατφόρμες εξισορρόπησης όπως η πλατφόρμα PICASSO για την ανταλλαγή aFRR, η πλατφόρμα MARI για την ανταλλαγή mFRR και η TERRE για την ανταλλαγή RR. Η τιμολόγηση ενέργειας εξισορρόπησης μπορεί να βασίζεται είτε στην τιμή προσφοράς (pay-as-bid) είτε στην οριακή τιμή (pay-as-cleared), με την τελευταία να προτιμάται στις περισσότερες ευρωπαϊκές αγορές. Οι συμμετέχοντες που αποκλίνουν από το πρόγραμμά τους χρεώνονται ή πιστώνονται με την τιμή απόκλισης (imbalance price).

Η Αγορά Εξισορρόπησης είναι καθοριστικής σημασίας για τη διασφάλιση της αξιοπιστίας και ασφάλειας του συστήματος ηλεκτρικής ενέργειας, την αποτελεσματική διαχείριση των διακυμάνσεων της παραγωγής από ΑΠΕ και την παροχή οικονομικών κινήτρων για επενδύσεις σε ευέλικτους πόρους. Με την αυξανόμενη διείσδυση των ΑΠΕ, η σημασία και η πολυπλοκότητα της Αγοράς Εξισορρόπησης αυξάνεται, απαιτώντας συνεχή εξέλιξη των μηχανισμών της και στενότερη διασυνοριακή συνεργασία μεταξύ των Διαχειριστών Συστημάτων Μεταφοράς.


% =============================================================================
\section{Σκοπός και Στόχοι της Εργασίας}
\label{sec:objectives}

Η παρούσα διπλωματική εργασία έχει ως κεντρικό στόχο την ανάπτυξη ενός ολοκληρωμένου, ανοιχτού κώδικα εργαλείου για την προσομοίωση της Προημερήσιας Αγοράς ηλεκτρικής ενέργειας, με έμφαση στον αλγόριθμο EUPHEMIA που χρησιμοποιείται για την εκκαθάριση της συζευγμένης ευρωπαϊκής αγοράς.

\subsection{Κίνητρο και Αναγκαιότητα}
\label{subsec:motivation}

Η κατανόηση των μηχανισμών λειτουργίας των αγορών ηλεκτρικής ενέργειας αποτελεί κρίσιμο παράγοντα για όλους τους εμπλεκόμενους φορείς, από τους παραγωγούς και τους προμηθευτές έως τους ρυθμιστές και τους ερευνητές. Ωστόσο, η πρόσβαση σε εργαλεία προσομοίωσης που αναπαράγουν με ακρίβεια τη λειτουργία της αγοράς είναι περιορισμένη, καθώς τα υπάρχοντα συστήματα είναι συνήθως κλειστά και ιδιόκτητα.

Η ενεργειακή μετάβαση και η αυξανόμενη διείσδυση των ΑΠΕ δημιουργούν νέες προκλήσεις για τη λειτουργία των αγορών ηλεκτρικής ενέργειας. Η μεταβλητότητα της παραγωγής από ανανεώσιμες πηγές, οι περιορισμοί του δικτύου μεταφοράς και η ανάγκη για διασυνοριακή συνεργασία καθιστούν τη μοντελοποίηση της αγοράς ένα σύνθετο πρόβλημα βελτιστοποίησης. Η δυνατότητα προσομοίωσης διαφορετικών σεναρίων και η κατανόηση των παραγόντων που επηρεάζουν τις τιμές της ηλεκτρικής ενέργειας είναι απαραίτητη για τη λήψη αποφάσεων τόσο σε επιχειρησιακό όσο και σε πολιτικό επίπεδο.

\subsection{Επιστημονική Συνεισφορά}
\label{subsec:contribution}

Η συνεισφορά της παρούσας εργασίας εντοπίζεται σε τρεις κύριους άξονες.

Πρώτον, στην \textbf{ανάπτυξη ολοκληρωμένης υποδομής δεδομένων}. Σχεδιάστηκε και υλοποιήθηκε ένα αυτοματοποιημένο σύστημα συλλογής δεδομένων από την πλατφόρμα διαφάνειας του ENTSO-E, το οποίο αποθηκεύει και διαχειρίζεται δεδομένα για μονάδες παραγωγής, φορτία, προβλέψεις ΑΠΕ, διασυνοριακές μεταφορές και μη διαθεσιμότητες μονάδων. Η υποδομή αυτή επιτρέπει την εκτέλεση προσομοιώσεων με πραγματικά δεδομένα για οποιαδήποτε ζώνη προσφοράς (bidding zone) της Ευρώπης.

Δεύτερον, στην \textbf{υλοποίηση αλγορίθμων βελτιστοποίησης αγοράς}. Αναπτύχθηκε μια σειρά από αλγορίθμους που καλύπτουν το πλήρες φάσμα της λειτουργίας της Προημερήσιας Αγοράς, συμπεριλαμβανομένου του προβλήματος Unit Commitment για τον προσδιορισμό του βέλτιστου προγράμματος λειτουργίας των μονάδων παραγωγής, της μηχανής στρατηγικών υποβολής προσφορών, του βιβλίου εντολών οριακού κόστους (merit-order) που κατασκευάζει ανταγωνιστικό αντιπαράδειγμα της αγοράς από τα θεμελιώδη μεγέθη της (τιμές καυσίμων και CO$_2$, αξία νερού, καθαρές εισαγωγές), του αλγορίθμου MPCC (Mathematical Programming with Complementarity Constraints) για την εκκαθάριση της αγοράς και της μοντελοποίησης δικτύου με περιορισμούς ATC (Available Transfer Capacity) για την πολυζωνική εκκαθάριση της αγοράς.

Τρίτον, στη \textbf{δημιουργία ανοιχτού κώδικα εργαλείου με ποσοτικά επικυρωμένη ακρίβεια}. Ολόκληρο το σύστημα αναπτύχθηκε σε γλώσσα Julia με χρήση του πλαισίου JuMP για τη μοντελοποίηση προβλημάτων βελτιστοποίησης, και ο κώδικας διατίθεται ελεύθερα με άδεια ανοιχτού κώδικα. Το σύστημα αναπαράγει τις τιμές της ελληνικής ζώνης με συντελεστή συσχέτισης 0{,}84 σε ανάλυση δεκαπενταλέπτου επί συνεχούς τετραμήνου (Κεφάλαιο~\ref{ch:results})· εξ όσων είμαστε σε θέση να γνωρίζουμε, πρόκειται για το πρώτο ανοιχτό εργαλείο που μοντελοποιεί την αγορά αυτή με τεκμηριωμένη ακρίβεια αυτού του επιπέδου.

\subsection{Ερευνητικά Ερωτήματα}
\label{subsec:research-questions}

Η εργασία επιχειρεί να απαντήσει στα ακόλουθα ερευνητικά ερωτήματα. Πρώτον, είναι δυνατή η αναπαραγωγή των τιμών της Προημερήσιας Αγοράς με ικανοποιητική ακρίβεια χρησιμοποιώντας δημοσίως διαθέσιμα δεδομένα και αλγορίθμους ανοιχτού κώδικα; Δεύτερον, ποια είναι η επίδραση των διαφορετικών στρατηγικών υποβολής προσφορών στις τελικές τιμές της αγοράς; Τρίτον, πώς επηρεάζουν οι περιορισμοί διασυνοριακής μεταφοράς τη διαμόρφωση των τιμών σε γειτονικές ζώνες προσφοράς;


% =============================================================================
\section{Δομή της Εργασίας}
\label{sec:thesis-structure}

Η παρούσα διπλωματική εργασία οργανώνεται σε έξι κεφάλαια, τα οποία καλύπτουν το θεωρητικό υπόβαθρο, την αρχιτεκτονική του συστήματος, την υλοποίηση, τα αποτελέσματα και τα συμπεράσματα.

Στο \textbf{Κεφάλαιο~\ref{ch:introduction}} (Εισαγωγή), το παρόν κεφάλαιο, παρουσιάζεται το πλαίσιο της εργασίας με αναφορά στην κλιματική αλλαγή, την ενεργειακή μετάβαση και τη δομή των ευρωπαϊκών αγορών ηλεκτρικής ενέργειας. Επίσης, καθορίζονται οι στόχοι και η συνεισφορά της εργασίας.

Στο \textbf{Κεφάλαιο~\ref{ch:theoretical-background}} (Θεωρητικό Υπόβαθρο) αναλύονται οι θεωρητικές βάσεις της εργασίας, συμπεριλαμβανομένων των μεθόδων μαθηματικού προγραμματισμού, του προβλήματος Unit Commitment, του αλγορίθμου EUPHEMIA και των στρατηγικών υποβολής προσφορών.

Στο \textbf{Κεφάλαιο~\ref{ch:system-architecture}} (Αρχιτεκτονική Συστήματος και Δεδομένα) περιγράφεται η αρχιτεκτονική του συστήματος που αναπτύχθηκε, με έμφαση στο pipeline συλλογής δεδομένων από το ENTSO-E, το σχήμα της βάσης δεδομένων και το οικοσύστημα Julia/JuMP.

Στο \textbf{Κεφάλαιο~\ref{ch:implementation}} (Υλοποίηση) παρουσιάζεται αναλυτικά η υλοποίηση των αλγορίθμων, συμπεριλαμβανομένου του solver για το πρόβλημα Unit Commitment, της μηχανής στρατηγικών υποβολής προσφορών, του βιβλίου εντολών οριακού κόστους, του αλγορίθμου MPCC για την εκκαθάριση της αγοράς και της μοντελοποίησης των περιορισμών δικτύου.

Στο \textbf{Κεφάλαιο~\ref{ch:results}} (Αποτελέσματα και Αξιολόγηση) παρουσιάζονται τα αποτελέσματα της εφαρμογής του συστήματος σε πραγματικά δεδομένα, η σύγκριση με τις πραγματικές τιμές της αγοράς και η ανάλυση της απόδοσης του συστήματος.

Στο \textbf{Κεφάλαιο~\ref{ch:conclusions}} (Συμπεράσματα) συνοψίζονται τα κύρια ευρήματα της εργασίας, αναφέρονται οι περιορισμοί και προτείνονται κατευθύνσεις για μελλοντική έρευνα.

Τέλος, στο \textbf{Παράρτημα} παρατίθενται συμπληρωματικά στοιχεία και σύνδεσμοι προς τον κώδικα της εφαρμογής.
